\documentclass[12pt,letterpaper]{article}

% --- Paquetes Principales ---
\usepackage[utf8]{inputenc}
\usepackage[T1]{fontenc}
\usepackage[spanish,es-tabla]{babel}
\usepackage[margin=2.5cm]{geometry}
\usepackage[skip=0.5em plus 0.1em minus 0.1em]{parskip}
\usepackage{graphicx}
\usepackage{amsmath,amssymb}
\usepackage{booktabs}
\usepackage{float}
\usepackage{tabularx}
\usepackage{array}
\usepackage{longtable}
\usepackage{fancyhdr}
\usepackage{titlesec}
\usepackage{lmodern}
\usepackage[expansion=false]{microtype}
\usepackage[table]{xcolor}
\usepackage{placeins}
\usepackage[hidelinks,pdfencoding=auto,psdextra]{hyperref}
% --- Paquete de Dibujo TikZ y sus Librerías ---
\usepackage{tikz}
\usetikzlibrary{shapes.geometric, arrows.meta, positioning, calc, backgrounds, babel}
\usepackage{caption}
\usepackage{etoolbox}
\usepackage{enumitem}
\usepackage{xurl}
\definecolor{docInk}{HTML}{183653}
\definecolor{docAccent}{HTML}{496A89}
\definecolor{docLine}{HTML}{8A9BAA}
\definecolor{docHead}{HTML}{E7EEF4}
\definecolor{docFill}{HTML}{F2F5F8}
\definecolor{docHigh}{HTML}{DCE6EF}
\colorlet{primary}{docInk}
\colorlet{secondary}{docAccent}
\captionsetup{font=small,labelfont={bf,color=docInk},skip=7pt}
\captionsetup[table]{position=top}
\arrayrulecolor{docLine}
\setlength{\tabcolsep}{4pt}
\setlength{\LTpre}{0.6em}
\setlength{\LTpost}{0.6em}
\setlength{\textfloatsep}{12pt plus 3pt minus 2pt}
\setlength{\intextsep}{10pt plus 3pt minus 2pt}
\AtBeginEnvironment{longtable}{\small\renewcommand{\arraystretch}{1.12}}
\setlist{topsep=0.35em,itemsep=0.25em,parsep=0pt,leftmargin=*}
\setlength{\emergencystretch}{2em}
\widowpenalty=10000
\clubpenalty=10000
\raggedbottom
\setcounter{tocdepth}{3}
\setcounter{topnumber}{3}
\setcounter{totalnumber}{5}
\makeatletter
\setlength{\@fptop}{0pt}
\makeatother

% --- Tolerancia y Control de Flotantes ---
\renewcommand{\topfraction}{0.9}
\renewcommand{\bottomfraction}{0.85}
\renewcommand{\textfraction}{0.08}
\renewcommand{\floatpagefraction}{0.65}

% --- Colores del Diagrama BPMN ---
\colorlet{bpmnBlue}{docFill}
\colorlet{bpmnBlueBorder}{docAccent}
\colorlet{bpmnTeal}{docFill}
\colorlet{bpmnTealBorder}{docAccent}
\colorlet{bpmnGreen}{docFill}
\colorlet{bpmnGreenBorder}{docAccent}
\colorlet{bpmnRed}{docFill}
\colorlet{bpmnRedBorder}{docAccent}
\colorlet{bpmnOrange}{docFill}
\colorlet{bpmnOrangeBorder}{docAccent}

% --- Ayudas para la sección 1.3 (necesidades) ---
\newcommand{\nec}[1]{\mbox{NEC-#1}}                       % evita partir "NEC-01" en dos renglones
\newcolumntype{L}[1]{>{\raggedright\arraybackslash}p{#1}}  % columna de ancho fijo alineada a la izquierda
\setlength{\LTcapwidth}{\textwidth}                        % títulos de longtable a todo el ancho

% --- Colores ---
\colorlet{primary}{docInk}
\colorlet{secondary}{docAccent}

% --- Configuración de Encabezados y Pies de Página ---
\pagestyle{fancy}
\setlength{\headheight}{14.5pt}
\addtolength{\topmargin}{-2.5pt}
\fancyhf{}
\fancyhead[L]{\small\scshape Sistemas de Información}
\fancyhead[R]{\small\scshape Hito 1}
\fancyfoot[C]{\thepage}
\renewcommand{\headrulewidth}{0.4pt}
\renewcommand{\footrulewidth}{0.2pt}

% --- Formato de Secciones y Subtítulos ---
\titleformat{\section}
  {\normalfont\Large\bfseries\color{primary}}{\thesection}{1em}{}
\titleformat{\subsection}
  {\normalfont\large\bfseries\color{secondary}}{\thesubsection}{1em}{}
\titleformat{\subsubsection}
  {\normalfont\normalsize\bfseries}{\thesubsubsection}{1em}{}


\begin{document}

% ==============================================================================
% PORTADA
% ==============================================================================

\begin{titlepage}
    \centering
    
    % Logo Institucional
    \IfFileExists{unal_logo.png}{%
        \includegraphics[width=0.3\textwidth]{unal_logo.png}%
    }{%
        \fbox{\parbox[c][2cm][c]{0.3\textwidth}{\centering\small Logo institucional\\\texttt{unal\_logo.png}}}%
    }\par
    \vspace{0.8cm}
    % Materia e Institución
    {\scshape\Large Universidad Nacional de Colombia\par}
    {\scshape\large Facultad de Ingeniería\par}
    {\scshape\normalsize Departamento de Ingeniería de Sistemas e Industrial\par}
    \vspace{0.5cm}
    {\Large\textbf{Sistemas de Información}\par}
    
 
    
    % Título y Subtítulo
    {\huge\bfseries\color{primary} Hito 1\par}
    \vspace{0.4cm}
    

    
    {\large
    Sebastian Fernando Gutiérrez Rojas\par
    Ronald Arturo Chávez Cabrera\par
    Camilo Andres Rojas Silva \par
    Angel Inyena Pasuy\par
    Jonathan Felipe López Núñez\par
    Juan David Peña Lozada\par
    }
    
    \vfill
    
    % Docente y Fecha
    {\normalsize\textbf{Docente:} Jhon Alexander Garcia Camargo}\par
    \vspace{0.4cm}
    {\normalsize Bogotá D.C., Colombia\par}
    {\normalsize 2026}
\end{titlepage}

% ==============================================================================
% TABLA DE CONTENIDO
% ==============================================================================
\clearpage
\tableofcontents
\clearpage
\section*{Presentación del proyecto}
\addcontentsline{toc}{section}{Presentación del proyecto}
Loard Construcciones S.A.S. presta servicios de montaje y mantenimiento para el sector de hidrocarburos. Este Hito 1 propone una plataforma de gestión operativa y documental que conecte los cortes de obra, el control de consumos y horas, la evidencia de calidad y las habilitaciones de seguridad y salud en el trabajo (SST). Su objetivo es reducir reprocesos y mejorar la oportunidad y trazabilidad de la información, con prioridad de implementación para el sistema integrado de gestión (SIG), SST y los registros que alimentan el dossier de calidad.

La plataforma se denomina \emph{ERP/SIG operativo}: ERP alude a la planificación de recursos empresariales y SIG al sistema integrado de gestión. En este documento, su alcance es operativo y documental; no incluye contabilidad, liquidación de nómina, facturación tributaria ni pagos. Puede exportar información operativa para los sistemas externos. Los clientes y laboratorios intercambian documentos por canales externos y no disponen de cuentas en la plataforma.

\paragraph{Fuentes y estado de la propuesta.}
Las reglas empresariales se retoman del diagnóstico documental y las validaciones declaradas por el equipo en octubre de 2026. El diseño, las prioridades, los puntos de esfuerzo y los valores de prueba constituyen una propuesta académica; no prueban que el sistema esté implementado ni que se hayan alcanzado las metas. El equipo debe conservar la evidencia de las validaciones y de la muestra documental revisada.

\paragraph{Códigos de trazabilidad.}
\textbf{NEC} identifica una necesidad del negocio; \textbf{RF}, un requisito funcional (qué debe hacer el sistema); \textbf{RNF}, un requisito no funcional (condiciones de calidad y operación); y \textbf{HU}, una historia de usuario. Por ejemplo, NEC-01 designa la necesidad de registrar cortes de obra; los números son identificadores, no niveles de prioridad. \textbf{CB} identifica un cuello de botella del proceso actual y \textbf{RES}, una restricción del proyecto. Las relaciones entre necesidades, requisitos e historias se consultan en la Tabla~\ref{tab:pb-trazabilidad}. Los \textbf{SP} (\emph{story points}) estiman esfuerzo relativo y no equivalen a horas.

\paragraph{Siglas del negocio.}
\textbf{APU}: análisis de precios unitarios; \textbf{END}: ensayos no destructivos (\emph{NDT} en inglés); \textbf{WPS}: especificación del procedimiento de soldadura; \textbf{PQR}: registro de calificación del procedimiento de soldadura; y \textbf{RBAC}: control de acceso basado en roles. Estas siglas se utilizan en el diagnóstico y en los requisitos posteriores.

\section{Diagnóstico y alineación estratégica}
El catálogo inicial define las necesidades que se citarán a lo largo del análisis. Porter y la cadena de valor explican sus presiones y puntos de dolor; la formulación del problema delimita el alcance. Después, SAM, la gobernanza y el cuadro de mando relacionan la propuesta con la estrategia, los responsables y los resultados esperados.

\paragraph{Catálogo de necesidades.}
Cada necesidad sigue la estructura ``la empresa necesita [capacidad] para [beneficio]'' y se vincula con su evidencia en el diagnóstico (Tabla~\ref{tab:nec-catalogo}). La prioridad usa MoSCoW: \textbf{M} (debe), \textbf{S} (debería) y \textbf{C} (podría), reconociendo la implementación prioritaria de la Fase 1 (SIG/SST y soporte operativo) y la Fase 2 (compras, cotizaciones y tablero ampliado).

{\small
\begin{longtable}{@{}L{1.8cm}L{7.2cm}L{4.6cm}L{1.2cm}@{}}
\caption{Catálogo de necesidades y fundamento en el diagnóstico.}\label{tab:nec-catalogo}\\
\toprule\rowcolor{docHead}\textbf{ID} & \textbf{Capacidad necesaria} & \textbf{Fundamento} & \textbf{Prior.}\\\midrule\endfirsthead
\toprule\rowcolor{docHead}\textbf{ID} & \textbf{Capacidad necesaria} & \textbf{Fundamento} & \textbf{Prior.}\\\midrule\endhead
\midrule\multicolumn{4}{r@{}}{\footnotesize Continúa en la página siguiente}\\\endfoot
\bottomrule\endlastfoot
\nec{01} & Registrar cortes semanales o mensuales en PC Windows, a cargo de residentes/supervisores. & Reportes informales de obra. & M\\
\nec{02} & Controlar consumos, saldos y merma presupuestada del 5\%, bajo custodia del residente. & Despacho central; sin almacenista. & M\\
\nec{03} & Conocer estado, ubicación, responsable e historial de herramientas y equipos. & Asignaciones y movimientos dispersos. & S\\
\nec{04} & Medir cantidades, horas, rendimientos y desviaciones respecto de la base cotizada. & Mano de obra y avance sin seguimiento oportuno. & M\\
\nec{05} & Preparar el dossier contractual y libro de tubería, con trazabilidad de ensayos y entrega en el plazo acordado. & Transcripción de END y compilación tardía. & M\\
\nec{06} & Mantener tarifas marco y APU para ítems nuevos, con históricos consultables. & Plantillas y tarifas dispersas. & M\\
\nec{07} & Generar y seguir ofertas con plantilla GC-FR04 e imprevistos parametrizados. & Adaptación manual de propuestas. & S\\
\nec{08} & Registrar caja menor y el ciclo de compras según montos y autorizaciones. & Urgencias y aprobaciones fragmentadas. & M\\
\nec{09} & Consultar avance, horas, materiales y alertas, con exportación a finanzas externas. & Información gerencial consolidada a mano. & M\\
\nec{10} & Analizar concentración por cliente y desviaciones por contrato. & Dependencia comercial de Chevron. & S\\
\nec{11} & Conservar lecciones aprendidas y retroalimentar futuras cotizaciones. & Experiencia no sistematizada. & C\\
\nec{12} & Controlar versiones y migrar selectivamente sin incumplir retenciones documentales. & Histórico de Dropbox y copias obsoletas. & M\\
\nec{13} & Evitar ediciones simultáneas conflictivas mediante reserva y liberación. & Copias en conflicto. & S\\
\nec{14} & Clasificar y recuperar soportes por metadatos y contenido indexado. & Búsqueda por carpetas. & S\\
\nec{15} & Aplicar acceso por rol a conceptos ocupacionales y márgenes; excluir historias clínicas y diagnósticos. & Datos sensibles y reserva médica. & M\\
\nec{16} & Formalizar aprobaciones y firmas autorizadas sobre una versión identificada. & Circulación física y firmas escaneadas. & S\\
\nec{17} & Conservar eventos de consulta, modificación y aprobación con integridad verificable. & Necesidad de evidencia para auditoría. & M\\
\nec{18} & Verificar habilitación por tarea, bloquear asignaciones no autorizadas y alertar vencimientos a 30/15 días. & Cursos y conceptos revisados manualmente. & M\\
\nec{19} & Registrar inducciones y capacitaciones con soporte y vigencia. & Planillas físicas dispersas. & S\\
\nec{20} & Capturar inspecciones/SST en Android sin red y sincronizar sin duplicados, fuera de áreas clasificadas. & Contingencias de hasta 6 h en Buenaventura. & S\\
\end{longtable}
}

\subsection{Análisis de las 5 Fuerzas de Porter}
\label{sec:porter-loard}
El marco de Porter permite examinar presiones externas sobre la rentabilidad del sector~\cite{hbs-cinco-fuerzas}. Se aplica a los servicios de montaje y mantenimiento para hidrocarburos que presta Loard. Las intensidades son valoraciones cualitativas del equipo basadas en el diagnóstico empresarial; no equivalen a un estudio de participación de mercado.

\subsubsection{Poder de negociación de los clientes: alto}
La concentración aproximada del 80\% de la facturación histórica en Chevron genera dependencia comercial. Los contratos marco fijan tarifas y exigen evidencia técnica, seguridad y entregas documentales oportunas. Según el caso, el dossier debe radicarse en un máximo de 10 días calendario desde la aceptación del acta de facturación y el pago tiene una condición comercial de 30 días. El sistema debe ayudar a cumplir los compromisos y documentar su seguimiento; no puede garantizar la aceptación ni el pago del cliente.

\textbf{Implicación:} sustentar el avance y los costos con información trazable y preparar oportunamente los soportes de calidad (NEC-04, NEC-05 y NEC-10).

\subsubsection{Poder de negociación de los proveedores: medio, con mayor presión en urgencias}
La empresa depende de consumibles industriales y de laboratorios de ensayos. La entrega centralizada desde Bogotá y la urgencia de reposición en frentes remotos pueden reducir el tiempo disponible para comparar ofertas. El diagnóstico no aporta concentración de proveedores ni costos de cambio; por tanto, la intensidad se mantiene como valoración preliminar.

\textbf{Implicación:} anticipar consumos y reposiciones y registrar el flujo de compras, respetando la autonomía de caja menor de hasta \$250.000 COP y la autorización central para montos superiores (NEC-02, NEC-03 y NEC-08).

\subsubsection{Rivalidad entre competidores existentes: media-alta}
Se infiere competencia por confiabilidad, experiencia y cumplimiento en la adjudicación de obras. En los contratos marco las tarifas están pactadas, por lo que la productividad de mano de obra y el control de consumos tienen especial relevancia. La intensidad propuesta requiere contrastarse con competidores y resultados de ofertas.

\textbf{Implicación:} conservar históricos de APU y rendimientos y comparar cantidades cotizadas con ejecutadas, sin confundir la reserva de imprevistos con utilidad asegurada (NEC-04, NEC-06 y NEC-07).

\subsubsection{Amenaza de nuevos entrantes: media-baja}
La experiencia, las calificaciones técnicas, los procedimientos de soldadura y las condiciones de seguridad y calidad exigidas por los clientes pueden actuar como barreras de entrada. No se dispone de mediciones de ingreso de nuevos contratistas; esta clasificación es una hipótesis del diagnóstico.

\textbf{Implicación:} mantener disponible y vigente la evidencia técnica y de habilitación del personal (NEC-12, NEC-14, NEC-18 y NEC-19).

\subsubsection{Amenaza de servicios sustitutos: baja-media}
La ejecución por equipos propios del cliente o determinadas soluciones de prefabricación podrían sustituir parte del servicio. Su adopción por los clientes de Loard no se encuentra demostrada en el material recibido. Contratar a otro proveedor del mismo servicio corresponde a rivalidad, no a sustitución.

\textbf{Implicación:} demostrar calidad, cumplimiento y productividad, e incorporar las lecciones de proyectos anteriores a la oferta de servicio (NEC-05 y NEC-11).

\paragraph{Síntesis.}
La dependencia del cliente principal es la presión externa mejor sustentada por el caso. El diagnóstico interno desarrolla cómo la información operativa y documental puede apoyar el cumplimiento de los compromisos, reducir reprocesos y proteger el margen.

\subsection{Cadena de valor y puntos de dolor tecnológicos}
\label{sec:cadena-loard}

La cadena de valor relaciona las actividades de Loard con el servicio final: montajes y mantenimiento de infraestructura con estricta conformidad técnica, seguridad ocupacional y evidencia documental íntegra. El modelo de Porter~\cite{hbs-cadena-valor} se ajusta a la realidad de la empresa, caracterizada por frentes distribuidos, personal de aproximadamente 40 trabajadores y delimitación de fronteras sistémicas.

\subsubsection{Actividades primarias}
{\small
\renewcommand{\arraystretch}{1.15}
\begin{longtable}{@{}L{2.6cm}L{3.5cm}L{5.2cm}L{3.1cm}@{}}
\caption{Actividades primarias y puntos de dolor tecnológicos.}\label{tab:cadena-primarias}\\
\toprule\rowcolor{docHead}
\textbf{Actividad} & \textbf{Valor aportado} & \textbf{Punto de dolor y efecto} & \textbf{Trazabilidad}\\
\midrule\endfirsthead
\toprule\rowcolor{docHead}\textbf{Actividad} & \textbf{Valor aportado} & \textbf{Punto de dolor y efecto} & \textbf{Trazabilidad}\\\midrule\endhead
\bottomrule\endfoot
Logística interna & Recepción, custodia y suministro de consumibles (electrodos, discos, gases) y herramientas en frentes. & Sin almacenista dedicado en obra (el residente asume la custodia directa); despacho centralizado desde Bogotá sin devolución de sobrantes; merma técnica estándar del 5\% sin trazabilidad al registrar o sincronizar la información. Pérdidas o desabastecimiento impactan la obra. & NEC-02 y NEC-03.\\\addlinespace
Operaciones & Montajes mecánicos y mantenimiento en facilidades y plantas con mano de obra calificada (costo más crítico). & Operarios no reportan en sistema; reporte a cargo exclusivo de residentes/supervisores. El reporte no es diario en todas las obras (semanal en proyectos, mensual en marco); seguimiento actual informal por teléfono y WhatsApp; demoras en aprobación por Dirección de Proyectos. Dispositivos móviles restringidos a inspecciones y SST, prohibidos en áreas clasificadas con riesgo de explosión. & NEC-01, NEC-04 y NEC-09.\\\addlinespace
Logística externa & Entrega de obras aceptadas y conformación del Dossier de Calidad para radicación. & Libro de Tubería (\emph{Weld Book} GT-FR64) bajo API 1104 y cartilla de soldadura sobre planos; reportes de END entregados por laboratorios en PDF con transcripción manual. Plazo estricto de entrega de 10 días tras acta de facturación; cobro a 30 días. Dependencia del conocimiento empírico de personal clave (Dilan). & NEC-05; apoyo de NEC-12 y NEC-16.\\\addlinespace
Marketing y ventas & Estructuración de cotizaciones en contratos marco y licitaciones independientes. & Tarifas congeladas en contratos marco y APU nuevos solo para ítems no contemplados; adaptación manual por residentes de plantilla marco oficial y revisión por Director de Proyectos (Yonny). Imprevistos (7--10\%) absorben variaciones de insumos ($\sim 1\%$), lluvias y tiempos muertos. & NEC-06 y NEC-07.\\\addlinespace
Servicio posventa & Cierre contractual, atención de observaciones y retroalimentación técnica. & Falta de registro sistemático de lecciones aprendidas y desviaciones reales de rendimientos frente a los APU iniciales. Dificultad para retroalimentar la plantilla de cotización oficial. & NEC-10 y NEC-11.\\
\end{longtable}}

\subsubsection{Actividades de apoyo}
{\small
\renewcommand{\arraystretch}{1.15}
\begin{longtable}{@{}L{2.6cm}L{3.5cm}L{5.2cm}L{3.1cm}@{}}
\caption{Actividades de apoyo y puntos de dolor tecnológicos.}\label{tab:cadena-apoyo}\\
\toprule\rowcolor{docHead}\textbf{Actividad} & \textbf{Valor aportado} & \textbf{Punto de dolor y efecto} & \textbf{Trazabilidad}\\\midrule\endfirsthead
\toprule\rowcolor{docHead}\textbf{Actividad} & \textbf{Valor aportado} & \textbf{Punto de dolor y efecto} & \textbf{Trazabilidad}\\\midrule\endhead
\bottomrule\endfoot
Infraestructura de la empresa & Dirección estratégica, gestión de proyectos, administración del SIG y frontera financiera. & Supervisión a cargo del Director de Proyectos (Yonny); delimitación limpia fuera de contabilidad, facturación y nómina (pero con necesidad de alimentar reportes de horas y consumos a finanzas). Firma legal de actas y dossieres reservada a Gerente General y Representante Legal mediante imágenes de firma manuscrita autorizada. & NEC-09, NEC-10 y NEC-17.\\\addlinespace
Gestión de recursos humanos y SST & Dotación de personal ($\sim 40$ trabajadores) y control estricto de seguridad y salud en el trabajo. & Coordinación a cargo de Edward y Johana; custodia física de carpetas del personal por Olga. Riesgo crítico por vencimientos de cursos de alturas, espacios confinados o exámenes médicos (exige bloqueo estricto sin excepciones ni dispensas). El repositorio administrativo debe excluir historias clínicas y diagnósticos. SST gestiona conceptos y restricciones ocupacionales; residentes consultan solo habilitación para la tarea y vigencia. & NEC-18 y NEC-19.\\\addlinespace
Desarrollo tecnológico & Plataforma de información documental, conectividad y captura de datos. & Repositorio actual en Dropbox (23.000 archivos históricos); la migración debe restringirse a los últimos 3 años más documentos técnicos permanentes (WPS/PQR), con depuración manual de formatos del SIG. Captura administrativa en PC Windows en oficina; móviles Android para inspecciones puntuales fuera de áreas clasificadas; planes corporativos activos con contingencia $\le 6$ h en Buenaventura. & NEC-12 a NEC-17, NEC-20.\\\addlinespace
Abastecimiento / compras & Suministro de insumos y contratación de ensayos y servicios. & Compras de emergencia con autonomía del residente hasta \$250.000 COP con caja menor (\$3M--\$5M COP); compras mayores a \$250.000 COP autorizadas por oficina central en Bogotá. Flujo formal: residente cotiza, Olga genera orden de compra y Gerencia General autoriza desembolso. Despacho centralizado desde Bogotá sin devoluciones. & NEC-08; relación con NEC-02 y NEC-03.\\
\end{longtable}}

\subsubsection{Mapa sintético de la cadena de valor}
\begin{figure}[!htbp]
\centering
\begin{tikzpicture}[
    box_apoyo/.style={
        draw=primary!50,
        fill=docFill,
        rounded corners=3pt,
        align=center,
        inner sep=4pt,
        text width=15.8cm
    },
    box_primaria/.style={
        draw=docAccent,
        fill=docHead,
        rounded corners=3pt,
        align=center,
        inner sep=3.5pt,
        text width=2.85cm,
        minimum height=2.7cm
    }
]

% Actividades de apoyo (4 bandas horizontales con posicionamiento relativo limpio)
\node (apoyo1) [box_apoyo] {
    \textbf{\small Infraestructura de la empresa}\\
    \footnotesize Dirección de proyectos (Yonny), SIG y frontera financiera; firmas legales Gerencia (NEC-09, 10, 17)
};

\node (apoyo2) [box_apoyo, below=5pt of apoyo1] {
    \textbf{\small Gestión de recursos humanos y SST (Edward, Johana, Olga)}\\
    \footnotesize Bloqueo estricto de vencimientos, alertas 30d/15d y reserva de conceptos y habilitación por tarea (NEC-18, 19)
};

\node (apoyo3) [box_apoyo, below=5pt of apoyo2] {
    \textbf{\small Desarrollo tecnológico}\\
    \footnotesize PC Windows, Android en inspecciones (fuera de áreas clasificadas), Dropbox últimos 3 años + WPS/PQR (NEC-12 a 17, 20)
};

\node (apoyo4) [box_apoyo, below=5pt of apoyo3] {
    \textbf{\small Abastecimiento / compras (Olga, Gerencia)}\\
    \footnotesize Caja menor (\$250k) en obra, compras Bogotá ($>\$250$k), despacho centralizado sin devoluciones (NEC-08)
};

% Actividades primarias (5 columnas alineadas bajo apoyo4 con idéntica altura)
\node (prim1) [box_primaria, anchor=north west] at ([yshift=-8pt]apoyo4.south west) {
    \textbf{\footnotesize Logística interna}\\[3pt]
    {\scriptsize Custodia residente\\merma base 5\%\\NEC-02, NEC-03}
};

\node (prim2) [box_primaria, anchor=north west] at ([xshift=6pt]prim1.north east) {
    \textbf{\footnotesize Operaciones}\\[3pt]
    {\scriptsize Cortes semanales/mensuales\\PC Windows / SST\\NEC-01, 04, 09}
};

\node (prim3) [box_primaria, anchor=north west] at ([xshift=6pt]prim2.north east) {
    \textbf{\footnotesize Logística externa}\\[3pt]
    {\scriptsize Dossier API 1104\\$\le 10$d / cobro 30d\\NEC-05, Dilan}
};

\node (prim4) [box_primaria, anchor=north west] at ([xshift=6pt]prim3.north east) {
    \textbf{\footnotesize Marketing y ventas}\\[3pt]
    {\scriptsize Tarifas fijas marco\\imprevistos 7--10\%\\NEC-06, NEC-07}
};

\node (prim5) [box_primaria, anchor=north west] at ([xshift=6pt]prim4.north east) {
    \textbf{\footnotesize Servicio posventa}\\[3pt]
    {\scriptsize Mano de obra crítica\\lecciones aprendidas\\NEC-10, NEC-11}
};

% Flecha horizontal perfectamente alineada debajo de las cajas primarias
\coordinate (arrow_left) at ([yshift=-9pt]prim1.south -| apoyo4.west);
\coordinate (arrow_right) at ([yshift=-9pt]prim1.south -| apoyo4.east);
\draw[-{Stealth[length=3.5mm,width=2.5mm]}, very thick, draw=primary!80!black] (arrow_left) -- (arrow_right);

% Texto explicativo de valor y margen
\node (valor) [align=center, text width=15.8cm, anchor=north, font=\footnotesize] at ([yshift=-15pt]prim3.south) {
    \textbf{Valor para el cliente:} Obras seguras y entrega de dossier en $\le 10$ días (pago a 30 días).\\[2pt]
    \textbf{Margen protegido:} Control de mano de obra crítica, merma $\le 5\%$ y caja menor en obra.
};

\end{tikzpicture}
\caption{Cadena de valor de Loard adaptada a servicios industriales. El orden sigue el modelo conceptual de Porter; no representa una secuencia de ejecución.}
\label{fig:cadena-loard}
\end{figure}

\paragraph{Relación entre los puntos de dolor.}
La desconexión en operaciones por la falta de un canal estructurado en computadores de oficina y frentes retrasa la consolidación de avances y consumos; ello repercute en la Dirección de Proyectos y en el área financiera, que requiere reportes de horas laboradas y materiales. En calidad y logística externa, la compilación artesanal del Libro de Tubería bajo API 1104 y de los ensayos END en PDF pone en riesgo el cumplimiento del plazo perentorio de 10 días tras el acta de facturación para la entrega del dossier, afectando el recaudo a 30 días. Asimismo, la dispersión en compras y la ausencia de alertas automáticas escalonadas en SST (30 y 15 días) con bloqueo estricto amenazan la habilitación del personal (40 trabajadores) ante Chevron. La solución tecnológica debe priorizar en su primera fase la Gestión Documental del SIG y SST, articulando el flujo de información verídica y oportuna sin intermediaciones innecesarias.

\subsection{Definición del problema bajo criterios de diseño de ingeniería}
Esta sección define el problema siguiendo el estándar de Diseño en Ingeniería del resultado de aprendizaje~2 de ABET (\emph{Outcome 2})~\cite{abet}: aplicar el diseño de ingeniería para producir soluciones que respondan a necesidades especificadas, considerando la salud pública, la seguridad y el bienestar, así como factores globales, culturales, sociales, ambientales y económicos. Por eso avanza en dos pasos: primero se identifican las necesidades de la organización y después se formula el problema con sus restricciones y consideraciones de diseño. La formulación distingue las necesidades, los objetivos y las restricciones de una propuesta preliminar de solución. Sus alternativas y su funcionamiento se desarrollan en el Discovery de producto y en la ingeniería de procesos.

% ------------------------------------------------------------------------------
\subsubsection{Identificación de necesidades}
\label{sec:necesidades}
% ------------------------------------------------------------------------------

\paragraph{Método y estado.}
Se utiliza la secuencia de elicitación, análisis, especificación y validación de requisitos del SWEBOK~v3~\cite{swebok}. El equipo declara un diagnóstico documental y validaciones de reglas empresariales en octubre de 2026. Este documento especifica necesidades y requisitos; las pruebas de producto y sus resultados permanecen pendientes. Una necesidad expresa la capacidad de negocio y su beneficio; su traducción verificable se encuentra en la matriz RF y los RNF.

Los síntomas principales son reportes informales, conciliación tardía de consumos y horas, dossier armado manualmente y dificultad para recuperar documentos vigentes. El catálogo de la Tabla~\ref{tab:nec-catalogo} los relaciona con capacidades concretas. Los interesados, su poder/interés y sus responsabilidades se presentan una sola vez en la Tabla~\ref{tab:gobernanza-raci}.


El catálogo se presenta al inicio del diagnóstico (Tabla~\ref{tab:nec-catalogo}) para definir los códigos antes de utilizarlos en Porter y en la cadena de valor. Esta sección explica su elaboración y prioridad.

\paragraph{Prioridad y alcance.}
El catálogo incluye 11 necesidades M, 8 S y 1 C. La Fase 1 prioriza SIG/SST y el soporte operativo del dossier; compras, cotizaciones y tablero ampliado se completan en Fase 2. Esta criticidad no equivale a esfuerzo: la distribución de puntos y capacidad se presenta en la Tabla~\ref{tab:pb-moscow}.

\subsubsection{Formulación del problema}
\label{sec:problema}
{\small
\begin{longtable}{@{}L{3.0cm}L{12.1cm}@{}}
\caption{Problema, objetivo, alcance y verificación.}\label{tab:nec-problema}\\
\toprule\rowcolor{docHead}\textbf{Elemento} & \textbf{Definición}\\\midrule\endfirsthead
\toprule\rowcolor{docHead}\textbf{Elemento} & \textbf{Definición}\\\midrule\endhead
\midrule\multicolumn{2}{r@{}}{\footnotesize Continúa en la página siguiente}\\\endfoot
\bottomrule\endlastfoot
Situación & Los reportes, costos operativos y soportes de calidad se consolidan mediante papel, WhatsApp, Excel y Dropbox en frentes distribuidos.\\
Problema & La información llega fragmentada o tarde para controlar desviaciones, comprobar habilitaciones y preparar el expediente contractual.\\
Objetivo & Conectar cortes, horas, consumos y documentos vigentes para reducir reprocesos y apoyar decisiones durante la ejecución.\\
Solución propuesta & Plataforma web operativa/documental en Windows, con funciones Android limitadas a inspección/SST y sincronización.\\
Frontera & Sin contabilidad, liquidación de nómina, pagos ni facturación tributaria. Exportación operativa a finanzas; clientes sin cuentas en el sistema.\\
Éxito propuesto & Dossier radicado en el plazo contractual; captura oportuna y adopción de al menos 90\%; menos tiempo de preparación y ausencia de asignaciones no habilitadas. Ver BSC.\\
\end{longtable}
}

\subsubsection{Consideraciones de diseño y reglas comunes}
El diseño considera los factores del Outcome 2 de ABET~\cite{abet}: seguridad mediante habilitación y evidencia técnica; factores globales por estándares y contratos; culturales por captura a cargo de residentes; sociales por privacidad y uso responsable de indicadores; ambientales por reducción de papel y desperdicio; económicos por costos operativos y oportunidad documental.

{\small
\begin{longtable}{@{}L{3.4cm}L{11.7cm}@{}}
\caption{Reglas comunes de negocio y restricciones del proyecto.}\label{tab:reglas-comunes}\\
\toprule\rowcolor{docHead}\textbf{Tema} & \textbf{Regla de referencia}\\\midrule\endfirsthead
\toprule\rowcolor{docHead}\textbf{Tema} & \textbf{Regla de referencia}\\\midrule\endhead
\midrule\multicolumn{2}{r@{}}{\footnotesize Continúa en la página siguiente}\\\endfoot
\bottomrule\endlastfoot
Cortes y dispositivos & Semanal en proyectos y mensual en contratos marco, desde PC Windows por residente/supervisor; los operarios no reportan en el sistema. Android solo para inspecciones/SST fuera de áreas clasificadas.\\
Revisión técnica & Yonny registra visado u observaciones sin impedir nuevos cortes ni seguimiento. Los consolidados finales distinguen versiones visadas.\\
Consumibles y compras & Custodia del residente; suministro desde Bogotá sin devolución ordinaria de sobrantes. Merma técnica presupuestada del 5\%, distinta de la desviación de costos. Caja menor de \$3--5 millones COP; autonomía por compra hasta \$250.000 COP. Montos mayores requieren autorización central; Olga genera orden y Gerencia autoriza desembolso.\\
Tarifas y APU & Tarifas pactadas en contratos marco; nuevos APU para ítems fuera del marco. Plantilla institucional, revisión de Yonny e imprevistos del 7--10\%.\\
Dossier y cobro & Seis grupos contractuales; API 1104 se aplica a las evidencias de soldadura pertinentes, no impone universalmente seis carpetas~\cite{api1104}. Referencia: radicación en máximo 10 días calendario desde aceptación del acta de facturación. Pago a 30 días según el contrato, sin garantía del software.\\
SST y privacidad & Certificaciones y conceptos vigentes para la tarea; alertas a 30/15 días y sin dispensas administrativas. SST gestiona conceptos y restricciones; obra consulta habilitación y vigencia. Historias clínicas y diagnósticos fuera de la plataforma~\cite{res1843}.\\
Firmas & Roles técnicos aprueban sus soportes; Gerencia/Representante Legal formalizan la representación contractual. La imagen de firma exige autorización autenticada, consentimiento, versión y evidencia de integridad~\cite{firma-electronica}.\\
Migración y retención & Últimos tres años operativos, formatos vigentes, WPS/PQR permanentes y registros de conservación obligatoria. La selección no autoriza borrados; aplicar retenciones por tipo y plazo legal de SST donde corresponda~\cite{decreto1072}.\\
Conectividad y alertas & Escenario de 6 h sin red en Buenaventura; RNF-07 prevé margen de almacenamiento móvil. La alerta de tiempo mayor que avance es preventiva: debe contrastarse con programación y no demuestra por sí sola sobrecosto.\\
\end{longtable}
}

\subsection{Modelo de Alineación Estratégica (SAM)}
\label{sec:sam}

% --- Definición de Colores del Diagrama SAM ---
\colorlet{samPurple}{docFill}
\colorlet{samBlue}{docFill}
\colorlet{samOrange}{docFill}
\colorlet{samGreen}{docFill}
\colorlet{samRedArrow}{docAccent}
\colorlet{samBlackArrow}{docAccent}

\begin{figure}[!htbp]
\centering
\begin{tikzpicture}[
    box/.style={
        rectangle,
        rounded corners=3pt,
        draw=black!55,
        line width=0.6pt,
        text width=6.6cm,
        minimum height=4.2cm,
        text=black,
        inner sep=8pt,
        font=\sffamily
    },
    redArrow/.style={
        {Stealth[length=2.3mm,width=1.6mm]}-{Stealth[length=2.3mm,width=1.6mm]},
        color=samRedArrow,
        line width=0.9pt,
        shorten <=5pt,
        shorten >=5pt
    }
]

% Nodos principales (Cuadrantes del SAM)
\node[box, fill=samPurple] (bs) at (0, 6.2) {
    {\centering\color{primary}\textbf{\large Estrategia de Negocio}\par}
    \vspace{3pt}\hrule height 0.5pt\vspace{6pt}
    {\footnotesize\raggedright
    $\bullet$ Montajes y obras \textit{Oil \& Gas} (Chevron $\sim 80\%$).\par\vspace{3pt}
    $\bullet$ Rigor en normas API 1104, ASME, ASTM, RUC e ISO.\par\vspace{3pt}
    $\bullet$ Mano de obra costo crítico; nómina $\sim 40$ personas.\par\vspace{3pt}
    $\bullet$ Tarifas marco fijas y margen límite no negociable.\par}
};

\node[box, fill=samBlue] (its) at (8.6, 6.2) {
    {\centering\color{primary}\textbf{\large Estrategia de TI}\par}
    \vspace{3pt}\hrule height 0.5pt\vspace{6pt}
    {\footnotesize\raggedright
    $\bullet$ Fases: F1 SIG y SST; F2 Compras y cotizaciones.\par\vspace{3pt}
    $\bullet$ Reportes estructurados de horas y materiales a finanzas.\par\vspace{3pt}
    $\bullet$ Control documental, metadatos y bloqueo estricto SST.\par\vspace{3pt}
    $\bullet$ Sin usuarios externos; firma manuscrita autorizada.\par}
};

\node[box, fill=samOrange] (bo) at (0, 0) {
    {\centering\color{primary}\textbf{\large Operación del Negocio}\par}
    \vspace{3pt}\hrule height 0.5pt\vspace{6pt}
    {\footnotesize\raggedright
    $\bullet$ Sede Bogotá y frentes en 5 facilidades del país.\par\vspace{3pt}
    $\bullet$ Yonny (Proyectos), Edward/Johana (SIG/SST), Olga (Compras).\par\vspace{3pt}
    $\bullet$ Residente custodia consumibles (merma 5\%) y caja (\$250k).\par\vspace{3pt}
    $\bullet$ Reporte semanal/mes simplificado; Dilan en dossier.\par}
};

\node[box, fill=samGreen] (ito) at (8.6, 0) {
    {\centering\color{primary}\textbf{\large Operación de TI}\par}
    \vspace{3pt}\hrule height 0.5pt\vspace{6pt}
    {\footnotesize\raggedright
    $\bullet$ PC Windows oficina principal; Android puntual (inspección/SST).\par\vspace{3pt}
    $\bullet$ Áreas clasificadas: sin celulares no certificados; datos corporativos ($\le 6$h).\par\vspace{3pt}
    $\bullet$ Migración Dropbox 3 años + WPS/PQR; depuración manual.\par\vspace{3pt}
    $\bullet$ Bloqueo estricto SST (sin excepción); alertas 30d/15d.\par}
};

% Flechas de alineación e integración
\draw[redArrow] (bs) -- (its);
\draw[redArrow] (bs) -- (bo);
\draw[redArrow] (its) -- (ito);
\draw[redArrow] (bo) -- (ito);

\end{tikzpicture}
\caption{Modelo de Alineación Estratégica (SAM) para Loard Construcciones S.A.S.}
\label{fig:sam_model}
\end{figure}

A continuación, se detalla la caracterización de cada una de las cuatro secciones del modelo adaptadas a la realidad estratégica, operacional y tecnológica confirmada por la compañía:

\subsubsection{Lectura y alineación del modelo}
La estrategia de negocio se apoya en confiabilidad técnica y cumplimiento para una cartera concentrada en Chevron. La estrategia de TI propone trazabilidad operativa/documental y despliegue gradual. La operación del negocio conserva los responsables, cortes y controles de obra; la operación de TI aporta repositorio, permisos, sincronización y soporte~\cite{chan-reich-2007}.

La relación causal es: mejores datos de obra permiten detectar desviaciones; documentos vigentes y habilitaciones reducen reprocesos y riesgos; un expediente completo facilita la radicación. Estos efectos se comprobarán con los indicadores del BSC, sin asumir que digitalizar garantiza rentabilidad, seguridad o pago.


\subsection{Gobernanza y marco de decisión}
\label{sec:gobernanza}
La propuesta distingue patrocinio y dirección de la gestión del proyecto, tomando COBIT y TOGAF como referencias de adaptación, no como implantaciones completas~\cite{cobit,togaf}. Se prevé una revisión quincenal de 45 minutos para demostración, aceptación y prioridades, y refinamiento semanal de 30 minutos para resolver datos, interfaces y riesgos.

\textbf{Lectura de RACI:} R ejecuta; A aprueba en su ámbito; C es consultado; e I es informado.

{\small
\begin{longtable}{@{}L{3.2cm}L{1.2cm}L{1.2cm}L{1.1cm}L{7.2cm}@{}}
\caption{Interesados y responsabilidades de gobernanza.}\label{tab:gobernanza-raci}\\
\toprule\rowcolor{docHead}\textbf{Interesado} & \textbf{Poder} & \textbf{Interés} & \textbf{RACI} & \textbf{Responsabilidad}\\\midrule\endfirsthead
\toprule\rowcolor{docHead}\textbf{Interesado} & \textbf{Poder} & \textbf{Interés} & \textbf{RACI} & \textbf{Responsabilidad}\\\midrule\endhead
\midrule\multicolumn{5}{r@{}}{\footnotesize Continúa en la página siguiente}\\\endfoot
\bottomrule\endlastfoot
Gerencia / Representante Legal & Alto & Alto & A & Patrocinar, aprobar alcance y recursos, autorizar desembolsos y formalización contractual.\\
Dirección de Proyectos / Yonny & Alto & Alto & A/R & Validar cotización y avance; resolver observaciones sin bloquear nuevos registros.\\
SIG/SST / Edward y Johana & Alto & Alto & A/R & Definir controles documentales y de habilitación; SST gestiona conceptos, no historias clínicas.\\
Administración / Olga & Medio & Alto & R/C & Generar órdenes y custodiar soportes administrativos autorizados del personal.\\
Calidad / Dilan & Medio & Alto & R/C & Estructurar libro de tubería y expediente contractual; validar integridad de soportes.\\
Residentes / supervisores & Medio & Alto & R & Capturar cortes, horas y consumos, custodiar recursos y probar los flujos de campo.\\
Cuadrillas / operarios & Bajo & Medio & I & Ejecutar actividades y recibir información de habilitación; no diligenciar el sistema.\\
Clientes externos & Alto & Medio & C/I & Definir condiciones contractuales y recibir entregables por canales externos, sin cuentas directas.\\
Finanzas / laboratorios & Medio & Medio & C/I & Recibir exportaciones operativas o suministrar reportes de ensayos; fuera del aplicativo como usuarios externos.\\
\end{longtable}
}

La tabla asigna responsabilidades por interesado; las aprobaciones de cada actividad deben acordarse en el piloto. Cambios de alcance, criterios y fase se registrarán con responsable y decisión; ninguna aprobación administrativa podrá dispensar requisitos de seguridad.


\subsection{Balanced Scorecard: objetivos e indicadores KPI y KRI}
El Balanced Scorecard (BSC), o cuadro de mando integral, relaciona objetivos e indicadores. Los KPI son indicadores clave de desempeño y los KRI, indicadores clave de riesgo. Permiten contrastar el valor esperado de la propuesta. Sus metas son preliminares; se documentarán fuente, responsable, frecuencia y línea base antes del piloto. Una meta no constituye un resultado alcanzado.

{\small
\begin{longtable}{@{}L{2.5cm}L{7.0cm}L{5.2cm}@{}}
\caption{Objetivos e indicadores propuestos por perspectiva del BSC.}\label{tab:resumen_kpi_kri}\\
\toprule\rowcolor{docHead}\textbf{Perspectiva} & \textbf{KPI y meta preliminar} & \textbf{KRI y control}\\\midrule\endfirsthead
\toprule\rowcolor{docHead}\textbf{Perspectiva} & \textbf{KPI y meta preliminar} & \textbf{KRI y control}\\\midrule\endhead
\midrule\multicolumn{3}{r@{}}{\footnotesize Continúa en la página siguiente}\\\endfoot
\bottomrule\endlastfoot
Financiera & Desviación de costos operativos de mano de obra e insumos frente al presupuesto: objetivo inferior al 5\%, considerando la merma ya presupuestada. & Diferencia de tiempo y avance; aviso preventivo que se contrasta con programación y costos, sin inferir pérdida automáticamente.\\
Cliente & Radicación del dossier en máximo 10 días calendario desde aceptación del acta. Reducción propuesta de al menos 20\% del tiempo activo de preparación. & Entregas incompletas, vencidas u observadas; no conformidades por soportes obsoletos. Objetivo: cero uso de versiones obsoletas.\\
Procesos internos & Al menos 90\% de cortes registrados oportunamente según su frecuencia. Formalización interna: objetivo menor a 24 h, con línea base pendiente. & Asignaciones sin requisitos vigentes: objetivo cero; alertas a 30/15 días y bloqueo de asignación. Visado no bloqueante para nuevos cortes.\\
Aprendizaje y crecimiento & Al menos 90\% de residentes/supervisores registran un corte sin ayuda al primer mes de uso. & Errores de captura y accesos indebidos a conceptos ocupacionales; evaluar usabilidad y permisos, sin almacenar diagnósticos.\\
\end{longtable}
}

Se medirá el costo por corte contra la misma base aprobada; la oportunidad como cortes registrados a tiempo sobre cortes esperados; y la adopción mediante una tarea observada. El tiempo activo suma horas de preparación, mientras el plazo contractual cuenta días hasta la constancia de radicación. Calidad revisa tiempos y devoluciones; Dirección de Proyectos valida avance y costos; SST verifica habilitaciones. La periodicidad se alinea con los cortes y se consolida mensualmente.


\FloatBarrier
\section{Discovery de producto}
A partir del diagnóstico y de los objetivos del cuadro de mando, el OST organiza oportunidades, alternativas y pruebas de supuestos. Los recorridos de cliente y residente complementan ese análisis al mostrar dónde aparecen las fricciones y qué roles intervienen. Ambos insumos orientan el diseño del proceso propuesto, sin acreditar todavía una solución validada.

\subsection{Opportunity Solution Tree (OST)}
\label{sec:ost}

\subsubsection{Propósito y resultado esperado}
El OST (árbol de oportunidades y soluciones) conecta resultado, oportunidades, alternativas y pruebas de supuestos~\cite{torres-ost}. Se enfoca en reducir al menos 20\% el tiempo activo de preparación del dossier frente a una línea base comparable, manteniendo integridad y plazo contractual. No se incluyen esperas de aceptación externa en el tiempo activo.
\[
\text{Reducción}=\frac{T_{\mathrm{base}}-T_{\mathrm{piloto}}}{T_{\mathrm{base}}}\times100\%
\]
Las comparaciones usarán expedientes semejantes en juntas, ensayos y cantidad de soportes. El plazo de radicación se mide aparte según la Tabla~\ref{tab:reglas-comunes}.

\subsubsection{Oportunidades identificadas}
\begin{enumerate}
\item \textbf{O1: información completa de campo.} El residente necesita cortes y evidencias asociados al proyecto, con revisión no bloqueante (NEC-01, NEC-02 y NEC-04).
\item \textbf{O2: recuperación de evidencia vigente.} Calidad necesita localizar la versión correcta sin comparar múltiples copias (NEC-12--14).
\item \textbf{O3: expediente completo sin transcripción repetida.} Dilan necesita integrar libro de tubería, cartilla y ensayos conforme a la lista contractual (NEC-05).
\end{enumerate}
Son oportunidades derivadas del diagnóstico; su validación de producto y las pruebas siguientes permanecen pendientes.


\subsubsection{Árbol de oportunidades y soluciones}

La Figura~\ref{fig:ost-loard} presenta dos alternativas por
oportunidad. Las pruebas del último nivel buscan evaluar los
supuestos de cada alternativa antes de decidir su implementación.

\begin{figure}[!htbp]
\centering
\resizebox{\textwidth}{!}{%
\begin{tikzpicture}[
    font=\sffamily\fontsize{8.5}{10}\selectfont,
    resultado/.style={
        rectangle,
        rounded corners=3pt,
        draw=primary,
        fill=docHead,
        text width=14.6cm,
        minimum height=1.5cm,
        align=center,
        inner sep=5pt
    },
    oportunidad/.style={
        rectangle,
        rounded corners=3pt,
        draw=black!60,
        fill=docHead,
        text width=4.5cm,
        minimum height=1.7cm,
        align=center,
        inner sep=4pt
    },
    solucion/.style={
        rectangle,
        rounded corners=3pt,
        draw=secondary,
        fill=docFill,
        text width=2.15cm,
        minimum height=2.5cm,
        align=center,
        inner sep=4pt,
        font=\sffamily\footnotesize
    },
    prueba/.style={
        rectangle,
        rounded corners=3pt,
        draw=black!55,
        fill=docFill,
        text width=2.15cm,
        minimum height=2.5cm,
        align=center,
        inner sep=4pt,
        font=\sffamily\footnotesize
    },
    enlace/.style={
        -{Stealth[length=2mm]},
        draw=black!70,
        thick
    }
]

\node[resultado] (r) at (7.7,0) {
    \textbf{Resultado esperado}\\
    Reducir al menos un 20\% el tiempo de preparación
    del dossier de calidad\\
    \textit{Meta preliminar; línea base pendiente}
};

\node[oportunidad] (o1) at (2.5,-3.3) {
    \textbf{O1. Información de campo}\\
    Contar con reportes y evidencias completos,
    asociados al proyecto
};

\node[oportunidad] (o2) at (7.7,-3.3) {
    \textbf{O2. Evidencia vigente}\\
    Encontrar rápidamente el soporte correcto
    y su versión aprobada
};

\node[oportunidad] (o3) at (12.9,-3.3) {
    \textbf{O3. Dossier completo}\\
    Reunir ensayos y entregables
    sin repetir transcripciones
};

\node[solucion] (s1) at (1.2,-6.7) {
    \textbf{S1}\\Formulario en PC\\con validación de\\campos del corte
};

\node[solucion] (s2) at (3.8,-6.7) {
    \textbf{S2}\\Plantilla única\\de corte y revisión\\no bloqueante
};

\node[solucion] (s3) at (6.4,-6.7) {
    \textbf{S3}\\Repositorio con\\metadatos y\\control de versiones
};

\node[solucion] (s4) at (9,-6.7) {
    \textbf{S4}\\Índice de soportes\\y carpeta de\\versiones aprobadas
};

\node[solucion] (s5) at (11.6,-6.7) {
    \textbf{S5}\\Generar dossier\\desde registros\\estructurados
};

\node[solucion] (s6) at (14.2,-6.7) {
    \textbf{S6}\\Lista contractual\\y carga directa\\de soportes
};

\node[prueba] (p1) at (1.2,-10.2) {
    \textbf{P1}\\Registrar un corte\\completo en PC\\sin ayuda
};

\node[prueba] (p2) at (3.8,-10.2) {
    \textbf{P2}\\Probar plantilla\\y visado sin detener\\nuevos registros
};

\node[prueba] (p3) at (6.4,-10.2) {
    \textbf{P3}\\Recuperar la versión\\vigente entre\\varias copias
};

\node[prueba] (p4) at (9,-10.2) {
    \textbf{P4}\\Comparar búsqueda\\por índice y\\por carpetas
};

\node[prueba] (p5) at (11.6,-10.2) {
    \textbf{P5}\\Generar un dossier\\y revisar integridad\\y tiempo activo
};

\node[prueba] (p6) at (14.2,-10.2) {
    \textbf{P6}\\Detectar faltantes\\y comprobar\\aprobaciones
};

\draw[enlace] (r.south) -- (7.7,-1.7) -| (o1.north);
\draw[enlace] (r.south) -- (o2.north);
\draw[enlace] (r.south) -- (7.7,-1.7) -| (o3.north);

\draw[enlace] (o1.south) -- (2.5,-4.6) -| (s1.north);
\draw[enlace] (o1.south) -- (2.5,-4.6) -| (s2.north);
\draw[enlace] (o2.south) -- (7.7,-4.6) -| (s3.north);
\draw[enlace] (o2.south) -- (7.7,-4.6) -| (s4.north);
\draw[enlace] (o3.south) -- (12.9,-4.6) -| (s5.north);
\draw[enlace] (o3.south) -- (12.9,-4.6) -| (s6.north);

\draw[enlace] (s1.south) -- (p1.north);
\draw[enlace] (s2.south) -- (p2.north);
\draw[enlace] (s3.south) -- (p3.north);
\draw[enlace] (s4.south) -- (p4.north);
\draw[enlace] (s5.south) -- (p5.north);
\draw[enlace] (s6.south) -- (p6.north);

\end{tikzpicture}%
}
\caption{OST preliminar para Loard Construcciones S.A.S.:
resultado, oportunidades, alternativas y pruebas de supuestos.
Las conexiones representan relaciones de exploración,
no una secuencia BPMN.}
\label{fig:ost-loard}
\end{figure}

\subsubsection{Pruebas de supuestos y decisión de producto}
Las pruebas utilizarán datos ficticios o autorizados y anonimizados. Se compararán tiempo, errores, integridad, usabilidad y esfuerzo de mantenimiento; no se presentan resultados ejecutados.

{\small
\begin{longtable}{@{}L{1.5cm}L{3.4cm}L{5.4cm}L{4.3cm}@{}}
\caption{Pruebas propuestas para las alternativas del OST.}\label{tab:ost-pruebas}\\
\toprule\rowcolor{docHead}\textbf{Prueba} & \textbf{Supuesto} & \textbf{Actividad} & \textbf{Evidencia}\\\midrule\endfirsthead
\toprule\rowcolor{docHead}\textbf{Prueba} & \textbf{Supuesto} & \textbf{Actividad} & \textbf{Evidencia}\\\midrule\endhead
\midrule\multicolumn{4}{r@{}}{\footnotesize Continúa en la página siguiente}\\\endfoot
\bottomrule\endlastfoot
P1 / S1 & El formulario en PC permite cortes completos sin ayuda continua. & Completar un corte simulado con actividades, horas y consumos. & Tiempo, completitud, errores y ayuda requerida.\\
P2 / S2 & La plantilla única y la revisión reducen devoluciones. & Usar una plantilla semanal/mensual y devolver observaciones sin impedir nuevos registros. & Campos faltantes, tiempo y continuidad del registro.\\
P3 / S3 & Metadatos y versiones facilitan la recuperación. & Localizar soportes vigentes entre varias versiones en un prototipo. & Selección correcta y tiempo de búsqueda.\\
P4 / S4 & Un índice sencillo puede mejorar la organización actual. & Comparar las mismas búsquedas con índice y carpetas. & Tiempo, errores y esfuerzo de actualización.\\
P5 / S5 & Los registros estructurados permiten compilar el dossier. & Generar un expediente de muestra con juntas, planos y PDF de END. & Integridad, correspondencia y horas activas.\\
P6 / S6 & La lista contractual detecta faltantes antes de radicar. & Revisar un expediente con omisiones conocidas y verificar aprobaciones. & Faltantes detectados y trazabilidad; aceptación y pago son externos.\\
\end{longtable}
}

S1/S3/S5 exploran una solución integrada; S2/S4/S6 permiten contrastarla con mejoras de menor complejidad o transición. La prioridad empresarial es SIG/SST y soporte del dossier; la selección de producto dependerá de pruebas y capacidad, no solo de preferencia tecnológica. La contingencia móvil de inspección/SST se verifica por separado mediante HU-03 y RNF-07.


\subsection{Customer Journey Map y Definición de Roles de Usuario}
El cliente petrolero recibe propuestas y soportes por correo/canales contractuales, sin cuenta en el sistema. El residente captura información y consulta el avance en PC; la movilidad se limita a inspección/SST. Las experiencias de las Tablas~\ref{tab:viaje-cliente} y~\ref{tab:viaje-residente} son valoraciones cualitativas del diagnóstico, no resultados de una encuesta.

{\small
\begin{longtable}{@{}L{1.9cm}L{3.6cm}L{1.0cm}L{4.4cm}L{3.5cm}@{}}
\caption{Recorrido del cliente petrolero.}\label{tab:viaje-cliente}\\
\toprule\rowcolor{docHead}\textbf{Etapa} & \textbf{Acción del actor} & \textbf{Exper.} & \textbf{Fricción} & \textbf{Oportunidad}\\\midrule\endfirsthead
\toprule\rowcolor{docHead}\textbf{Etapa} & \textbf{Acción del actor} & \textbf{Exper.} & \textbf{Fricción} & \textbf{Oportunidad}\\\midrule\endhead
\midrule\multicolumn{5}{r@{}}{\footnotesize Continúa en la página siguiente}\\\endfoot
\bottomrule\endlastfoot
Solicitud & Define alcance y condiciones del servicio. & Media & Esperar precios e información técnica consistentes. & Catálogo y APU históricos (NEC-06).\\
Oferta & Evalúa propuesta y adjudica. & Media & Adaptación manual de cotizaciones fuera de marco. & Plantilla y revisión técnica (NEC-06, 07).\\
Habilitación & Comprueba requisitos de ingreso. & Baja & Certificados vencidos o evidencia dispersa. & Verificación y alertas SST (NEC-15, 18).\\
Ejecución & Revisa cortes y supervisa cumplimiento. & Media & Avance y rendimientos comunicados de forma informal. & Cortes estructurados (NEC-01, 04).\\
Calidad & Revisa dossier, planos y ensayos. & Baja & Faltantes y transcripción tardía. & Expediente contractual trazable (NEC-05).\\
Pago & Recibe facturación del contratista y tramita pago según contrato. & Media & Demoras de soportes y formalización. & Aprobaciones y constancia de radicación (NEC-16).\\
Posventa & Formula observaciones y audita. & Media & Difícil reconstrucción de versiones y decisiones. & Eventos y lecciones (NEC-11, 17).\\
\end{longtable}
}

{\small
\begin{longtable}{@{}L{1.9cm}L{3.6cm}L{1.0cm}L{4.4cm}L{3.5cm}@{}}
\caption{Recorrido del ingeniero residente.}\label{tab:viaje-residente}\\
\toprule\rowcolor{docHead}\textbf{Etapa} & \textbf{Acción del actor} & \textbf{Exper.} & \textbf{Fricción} & \textbf{Oportunidad}\\\midrule\endfirsthead
\toprule\rowcolor{docHead}\textbf{Etapa} & \textbf{Acción del actor} & \textbf{Exper.} & \textbf{Fricción} & \textbf{Oportunidad}\\\midrule\endhead
\midrule\multicolumn{5}{r@{}}{\footnotesize Continúa en la página siguiente}\\\endfoot
\bottomrule\endlastfoot
Recepción & Revisa contrato y formatos aplicables. & Baja & Riesgo de usar versión obsoleta. & Búsqueda de documentos vigentes (NEC-12, 14).\\
Planeación & Asigna personal y equipos. & Baja & Difícil verificar ubicación y requisitos para la tarea. & Habilitación y bitácora (NEC-03, 15, 18).\\
Abastecimiento & Recibe insumos y solicita compras. & Baja & Reposición urgente y autorizaciones dispersas. & Custodia, saldos y flujo por monto (NEC-02, 08).\\
Reporte & Registra avances y atiende observaciones. & Media & Seguimiento por teléfono/WhatsApp y doble digitación. & Corte simplificado y visado no bloqueante (NEC-01, 04).\\
Cierre & Reúne soportes para calidad. & Baja & Compilación manual contra el plazo contractual. & Dossier y aprobación trazables (NEC-05, 16).\\
\end{longtable}
}

Los momentos de mayor fricción son habilitación y entrega técnica. La causa transversal es la desconexión entre obra y oficina. El diseño debe simplificar el registro del residente y asegurar evidencia vigente, evitando trasladarle más tareas administrativas sin beneficio observable.


\FloatBarrier
\section{Ingeniería de procesos (BPMN 2.0)}
El modelo As-Is representa el proceso actual y sus cuellos de botella. A continuación, el To-Be plantea cambios en captura, revisión, consolidación y gestión de evidencias. Su comparación permite concretar los requerimientos de la sección~\ref{sec:requerimientos}.

\subsection{Diagrama As-Is: Mapa de Calor y Cuellos de Botella Actuales}
\label{sec:asis-loard}

\subsubsection{Alcance y convenciones}
El As-Is sigue la información desde los apuntes de cuadrilla hasta la radicación del dossier. Los operarios aportan notas informales; el residente consolida cortes semanales/mensuales, administración concilia y calidad prepara los soportes. Es un modelo documental del equipo, pendiente de recorrido detallado con responsables; no acredita tiempos medidos.

Se usa un pool de Loard y cuatro carriles internos. Círculos: inicio/fin; rectángulos: tareas; rombos X: alternativas exclusivas; flechas continuas: secuencia~\cite{omg-bpmn}. La intensidad azul añade criticidad cualitativa, también identificada por código en la matriz; no cambia la semántica BPMN.


\begin{figure}[p]
\centering
\begin{tikzpicture}[
 font=\sffamily\fontsize{8.5}{10}\selectfont,
 tarea/.style={rectangle,rounded corners=3pt,draw=black!70,fill=docFill,text width=3.1cm,minimum height=1.65cm,align=center,inner sep=3pt},
 alta/.style={tarea,fill=docHigh,draw=docAccent},
 media/.style={tarea,fill=docFill,draw=docLine},
 xor/.style={diamond,draw=black!80,fill=white,minimum size=8mm,inner sep=0pt,font=\bfseries\small},
 flujo/.style={-{Stealth[length=2mm,width=1.4mm]},thick,draw=black!80},
 evento/.style={circle,draw=black!80,fill=white,minimum size=7mm,inner sep=0pt},
 texto/.style={align=center,font=\sffamily\fontsize{8}{9.5}\selectfont,text width=3.3cm}]
\fill[gray!3] (0,.45) rectangle (15.6,-18);
\draw[thick,docLine] (0,1.25) rectangle (15.6,-18);
\fill[docHead] (.01,.46) rectangle (15.59,1.24);
\draw[thick] (0,.45) -- (15.6,.45);
\node[font=\sffamily\bfseries\small] at (7.8,.85) {Loard Construcciones S.A.S. -- Proceso actual (As-Is)};
\foreach \x in {3.9,7.8,11.7} {\draw[gray!65] (\x,.45) -- (\x,-18);}
\draw[gray!65] (0,-1.05) -- (15.6,-1.05);
\node[texto,font=\sffamily\bfseries\fontsize{8.5}{10}\selectfont] at (1.95,-.3) {Cuadrilla / operario\\Apuntes en frente};
\node[texto,font=\sffamily\bfseries\fontsize{8.5}{10}\selectfont] at (5.85,-.3) {Ingeniero residente /\\Supervisor};
\node[texto,font=\sffamily\bfseries\fontsize{8.5}{10}\selectfont] at (9.75,-.3) {Administración\\del proyecto};
\node[texto,font=\sffamily\bfseries\fontsize{8.5}{10}\selectfont] at (13.65,-.3) {Coordinación SIG /\\Calidad y responsables\\de firma};
\node[evento] (inicio) at (1.95,-1.75) {};
\node[texto,text width=1.5cm,anchor=west] at (2.4,-1.75) {Fin de jornada\\en frente};
\node[alta] (reporte) at (1.95,-3.2) {\textbf{CB-01}\\Entregar apuntes\\informales en papel\\al supervisor en frente};
\node[media] (revision) at (5.85,-4.6) {\textbf{CB-02}\\Recopilar apuntes,\\revisar y transmitir por\\WhatsApp / hojas cálculo};
\node[xor] (g1) at (5.85,-6.75) {X};
\node[texto,anchor=south] at (5.85,-6.25) {¿Información suficiente?};
\node[media] (corregir) at (1.95,-6.75) {Aclarar o completar\\apuntes en frente};
\node[alta] (consolidar) at (9.75,-8.2) {\textbf{CB-03}\\Consolidar reportes\\del corte en hojas\\de cálculo};
\node[alta] (costos) at (9.75,-10.55) {\textbf{CB-04}\\Conciliar consumos\\y costos acumulados};
\node[alta] (dossier) at (13.65,-10.55) {\textbf{CB-05}\\Transcribir ensayos y\\armar dossier del corte};
\node[xor] (g2) at (13.65,-12.95) {X};
\node[texto,anchor=south] at (13.65,-12.45) {¿Soportes completos\\y consistentes?};
\node[alta] (buscar) at (9.75,-12.95) {\textbf{CB-06}\\Buscar / corregir\\soportes en carpetas\\y copias};
\node[media] (firmar) at (13.65,-15.1) {\textbf{CB-07}\\Imprimir, obtener\\firmas y escanear actas};
\node[tarea,minimum height=1.3cm] (radicar) at (9.75,-16.3) {Radicar actas\\y dossier del corte};
\node[evento,line width=1.5pt] (fin) at (9.75,-17.5) {};
\node[texto,text width=3cm,anchor=west] at (10.2,-17.5) {Actas y soportes\\radicados};
% Íconos de tipificación BPMN 2.0 en tareas As-Is
\draw[flujo] (inicio.south) -- (reporte.north);
\draw[flujo] (reporte.east) -- (3.75,-3.2) |- (revision.west);
\draw[flujo] (revision.south) -- (g1.north);
\draw[flujo] (g1.west) -- node[above,font=\scriptsize]{No} (corregir.east);
\draw[flujo] (corregir.west) -- (.15,-6.75) -- (.15,-3.2) -- (reporte.west);
\draw[flujo] (g1.south) -- (5.85,-8.2) -- node[above,font=\scriptsize]{Sí} (consolidar.west);
\draw[flujo] (consolidar.south) -- (costos.north);
% Conector CB-04 -> CB-05: horizontal, con espacio libre entre las tareas.
\draw[flujo] (costos.east) -- (dossier.west);
\draw[flujo] (dossier.south) -- (g2.north);
\draw[flujo] (g2.west) -- node[above,font=\scriptsize]{No} (buscar.east);
% Retorno CB-06 -> CB-05: recorre el margen izquierdo y el hueco entre filas.
\draw[flujo] (buscar.west) -- (8,-12.95) -- (8,-9.35) -- (13.65,-9.35) -- (dossier.north);
\draw[flujo] (g2.south) -- node[right,font=\scriptsize]{Sí} (firmar.north);
\draw[flujo] (firmar.south) -- (13.65,-16.3) -- (radicar.east);
\draw[flujo] (radicar.south) -- (fin.north);
\end{tikzpicture}
\caption{Modelo BPMN del flujo actual. Azul medio: criticidad alta; azul muy claro: media; fondo neutro: contexto. La criticidad también se identifica por código en la matriz. Los códigos CB se explican en la Tabla~\ref{tab:cuellos-loard}. Modelo documental y procedimental fundamentado en análisis de artefactos y reglas operativas confirmadas.}
\label{fig:asis-loard}
\end{figure}

\subsubsection{Cuellos de botella y línea base}
{\small
\begin{longtable}{@{}L{1.6cm}L{9.8cm}L{3.0cm}@{}}
\caption{Criticidad y relación con las necesidades.}\label{tab:cuellos-loard}\\
\toprule\rowcolor{docHead}\textbf{Código} & \textbf{Problema y efecto} & \textbf{Nivel / NEC}\\\midrule\endfirsthead
\toprule\rowcolor{docHead}\textbf{Código} & \textbf{Problema y efecto} & \textbf{Nivel / NEC}\\\midrule\endhead
\midrule\multicolumn{3}{r@{}}{\footnotesize Continúa en la página siguiente}\\\endfoot
\bottomrule\endlastfoot
CB-01 & Apuntes y mensajes dispersos: recuperar e interpretar antes de consolidar. & Alta / 01, 02\\
CB-02 & Revisión sin estructura: aclaraciones y devoluciones. & Media / 01, 04\\
CB-03 & Consolidación manual: desfase entre avance y seguimiento. & Alta / 04, 09\\
CB-04 & Conciliación tardía: menor posibilidad de corregir desviaciones. & Alta / 02, 09\\
CB-05 & Transcripción y dossier tardío: retraso de radicación. & Alta / 05\\
CB-06 & Copias y carpetas: dificultad para recuperar versión vigente. & Alta / 12--14\\
CB-07 & Circulación física de firmas: espera de formalización. & Media / 16\\
\end{longtable}
}

La línea base deberá registrar horas activas de preparación, devoluciones y días hasta radicación. Las reglas de periodicidad fueron aportadas por el equipo; la jerarquía de criticidad y el flujo detallado son decisiones de modelado, no mediciones de duración.


\subsection{Diagrama To-Be: Proceso Optimizado con el Sistema de Información}
La propuesta centraliza el corte y su evidencia manteniendo los responsables de la Tabla~\ref{tab:gobernanza-raci}. Al guardar, el seguimiento se actualiza en trámite y la Dirección de Proyectos recibe la revisión. Las observaciones vuelven sobre el mismo corte sin detener nuevos registros; el expediente final requiere visado y revisión de calidad antes de la autorización de firmas y radicación.

\subsubsection{Proceso propuesto}
La Figura~\ref{fig:bpmn_tobe} muestra el alcance global de ambas fases. Las compuertas + dividen y reúnen revisión y preparación documental. El dossier se conforma progresivamente y se radica al cierre, no en cada corte. La Fase 1 ofrece consulta básica de registros y estados; el tablero ampliado se completa en Fase 2.


\begin{figure}[p]
\centering
\begin{tikzpicture}[
 font=\sffamily\fontsize{8.5}{10}\selectfont,
 tarea/.style={rectangle,rounded corners=3pt,draw=docAccent,fill=docFill,text width=3.1cm,minimum height=1.65cm,align=center,inner sep=3pt},
 automatica/.style={tarea,fill=docHead},
 compuerta/.style={diamond,draw=docInk,fill=white,minimum size=8mm,inner sep=0pt,font=\bfseries\small},
 flujo/.style={-{Stealth[length=2mm,width=1.4mm]},thick,draw=docInk},
 evento/.style={circle,draw=docInk,fill=white,minimum size=7mm,inner sep=0pt},
 texto/.style={align=center,font=\sffamily\fontsize{8}{9.5}\selectfont,text width=3.3cm}]
\draw[thick,docLine] (0,1.25) rectangle (15.6,-18.1);
\fill[docHead] (.01,.46) rectangle (15.59,1.24);
\draw[docLine] (0,.45) -- (15.6,.45);
\node[font=\sffamily\bfseries\small,text=docInk] at (7.8,.85) {Loard Construcciones S.A.S. -- Proceso propuesto (To-Be)};
\foreach \x in {3.9,7.8,11.7} {\draw[docLine] (\x,.45) -- (\x,-18.1);}
\draw[docLine] (0,-1.05) -- (15.6,-1.05);
\node[texto,font=\sffamily\bfseries\fontsize{8.5}{10}\selectfont] at (1.95,-.3) {Residente /\\Supervisor};
\node[texto,font=\sffamily\bfseries\fontsize{8.5}{10}\selectfont] at (5.85,-.3) {Dirección de Proyectos\\(Yonny)};
\node[texto,font=\sffamily\bfseries\fontsize{8.5}{10}\selectfont] at (9.75,-.3) {Sistema de Información\\ERP/SIG operativo};
\node[texto,font=\sffamily\bfseries\fontsize{8.5}{10}\selectfont] at (13.65,-.3) {Calidad, administración\\y representación legal};
\node[evento] (inicio) at (1.95,-1.65) {};
\node[texto,text width=1.3cm,anchor=west] at (2.4,-1.65) {Corte / cierre\\de obra};
\node[tarea] (captura) at (1.95,-3.1) {Registrar corte\\semanal / mensual,\\consumos y horas\\\textbf{NEC-01, 02, 04}};
\node[automatica] (registrar) at (9.75,-3.1) {Guardar versión\\del corte en trámite\\\textbf{NEC-01, 12, 14}};
\node[compuerta] (division) at (9.75,-4.6) {+};
\node[automatica] (tablero) at (9.75,-6.3) {Actualizar seguimiento\\con datos en trámite\\y alertas operativas\\\textbf{NEC-04, 09}};
\node[tarea] (revisar) at (5.85,-6.3) {Revisar corte y\\registrar visado u\\observaciones\\\textbf{NEC-04}};
\node[compuerta] (decision) at (5.85,-8.65) {X};
\node[texto,anchor=south] at (5.85,-8.15) {¿Conforme?};
\node[tarea] (corregir) at (1.95,-8.65) {Subsanar observaciones\\y guardar nueva versión\\del mismo corte};
\node[automatica] (dossier) at (9.75,-8.65) {Preparar dossier\\y verificar soportes\\\textbf{NEC-05, 12}};
\node[compuerta] (union) at (9.75,-10.75) {+};
\node[tarea] (calidad) at (13.65,-12.2) {Revisar expediente\\final y soportes\\contractuales\\\textbf{NEC-05}};
\node[tarea] (firmar) at (13.65,-14.6) {Autorizar firmas de\\representación sobre\\la versión aprobada\\\textbf{NEC-16}};
\node[tarea,minimum height=1.3cm] (radicar) at (9.75,-16.35) {Radicar dossier y\\registrar constancia};
\node[evento,line width=1.5pt] (fin) at (9.75,-17.65) {};
\node[texto,text width=3cm,anchor=west] at (10.2,-17.65) {Dossier radicado\\Plazo contractual: 10 días};
\draw[flujo] (inicio.south) -- (captura.north);
\draw[flujo] (captura.east) -- (registrar.west);
\draw[flujo] (registrar.south) -- (division.north);
\draw[flujo] (division.south) -- (tablero.north);
\draw[flujo] (division.west) -- (5.85,-4.6) -- (revisar.north);
\draw[flujo] (tablero.south) -- (dossier.north);
\draw[flujo] (revisar.south) -- (decision.north);
\draw[flujo] (decision.west) -- node[above,font=\scriptsize]{No} (corregir.east);
\draw[flujo] (corregir.north) -- (1.95,-4.85) -- (4.05,-4.85) -- (4.05,-6.3) -- (revisar.west);
\draw[flujo] (decision.east) -- node[above,font=\scriptsize]{Sí} (7.98,-8.65) -- (7.98,-10.75) -- (union.west);
\draw[flujo] (dossier.south) -- (union.north);
\draw[flujo] (union.east) -- (13.65,-10.75) -- (calidad.north);
\draw[flujo] (calidad.south) -- (firmar.north);
\draw[flujo] (firmar.south) -- (13.65,-16.35) -- (radicar.east);
\draw[flujo] (radicar.south) -- (fin.north);
\end{tikzpicture}
\caption{To-Be global: tareas de usuario en azul muy claro y tareas automáticas en azul medio. Las compuertas + dividen y reúnen las ramas de revisión y preparación. La revisión no bloquea nuevos cortes ni su seguimiento; la radicación final requiere el expediente revisado.}
\label{fig:bpmn_tobe}
\end{figure}

\subsubsection{Matriz de Impacto y Transformación Operativa (As-Is vs. To-Be)}

La Tabla~\ref{tab:asis_tobe} resume los cambios en la arquitectura operativa y su impacto directo en los indicadores clave de rendimiento (KPI Targets) del negocio:

\begin{table}[!htbp]
\centering
\small
\renewcommand{\arraystretch}{1.25}
\begin{tabularx}{\textwidth}{>{\raggedright\arraybackslash}L{2.6cm} X X >{\raggedright\arraybackslash}L{3.2cm}}
\toprule\rowcolor{docHead}
\textbf{Dimensión} & \textbf{Estado Actual (As-Is)} & \textbf{Estado Optimizado (To-Be)} & \textbf{Impacto Esperado / Target} \\
\midrule
\textbf{Captura de Campo} & Notas en papel y reportes informales por WhatsApp; transcripción manual. & Captura por Residente/Supervisor en PC Windows (móvil offline solo SST en áreas no clasificadas) (\textbf{NEC-01}). & Al menos 90\% de cortes semanales/mensuales registrados oportunamente durante el piloto. \\
\addlinespace
\textbf{Control de Consumibles y Mano de Obra} & Despacho desde Bogotá sin control de sitio; reconciliación tardía al cierre; mano de obra sin seguimiento vs cotización. & Registro por frente contra cotización base; alerta si \% tiempo $>$ \% avance; 5\% chatarra técnica; compras $\le \$250$k COP (\textbf{NEC-02, NEC-06}). & Control al registrar o sincronizar la información del costo operativo de mano de obra ($\sim 40$ trabajadores) y desvíos de insumos $< 5\%$. \\
\addlinespace
\textbf{Gestión Documental} & Repositorio Dropbox con duplicados, carpetas personales y versiones manuales (\texttt{v1}, \texttt{final}). & Migración curada (últimos 3 años + WPS/PQR permanentes); codificación ISO 9001 y control de bloqueo (\textbf{NEC-12, NEC-13}). & 0 no conformidades en auditorías por uso de versión obsoleta. \\
\addlinespace
\textbf{Dossier de Calidad} & Transcripción manual de ensayos NDT; armado artesanal tardío que demora cobro a clientes (Chevron). & Cartilla de soldadura sobre planos, metadatos NDT PDF y 6 carpetas contractuales consolidándose en paralelo (\textbf{NEC-05}). & Entrega de dossier en $\le 10$ días calendario desde la aceptación del acta de facturación; pago a 30 días. \\
\addlinespace
\textbf{Cumplimiento SST y Privacidad} & Control manual en Excel; riesgo de bloqueo en puerta o filtración de diagnósticos médicos. & Bloqueo estricto por no aptitud médica (sin dispensas); alertas a 30 y 15 días; conceptos ocupacionales restringidos por rol (\textbf{NEC-18, NEC-19}). & 0 operarios no habilitados en frente; Validación de accesos, habilitación y reserva de datos en las pruebas del piloto. \\
\addlinespace
\textbf{Flujo de Aprobación y Formalización} & Impresión en papel, circulación física para firma de Gerencia y escaneo. & Aprobación técnica no bloqueante de Yonny e inserción de imagen de firma manuscrita digitalizada (firmas contractuales de Gerencia) (\textbf{NEC-16}). & Meta preliminar: formalización interna en $<$24 h, con línea base pendiente. \\
\bottomrule
\end{tabularx}
\caption{Matriz Comparativa de Transformación Operativa: Estado As-Is vs. Estado To-Be.}
\label{tab:asis_tobe}
\end{table}

\paragraph{Controles y validación.}
El proceso aplica las reglas comunes de la Tabla~\ref{tab:reglas-comunes}: permisos por rol, integridad de eventos, habilitación por tarea, protección de conceptos ocupacionales y aprobación de versiones. Se comprobarán mediante los RF, los RNF y las historias de usuario de la sección~\ref{sec:requerimientos}. El software apoya el cumplimiento y la detección de errores; no sustituye controles físicos de planta, criterio médico ni aceptación contractual externa.


\FloatBarrier
\section{Especificación de requerimientos}
\label{sec:requerimientos}
El diagnóstico, el Discovery y el contraste As-Is/To-Be fundamentan esta especificación. Primero se presentan los requerimientos funcionales y su trazabilidad; después, los atributos de calidad y las restricciones; finalmente, las historias de usuario, sus criterios de aceptación y la planificación por fases. Los tres bloques describen un mismo alcance desde perspectivas complementarias.

Quedan fuera la integración contable, la facturación, los pagos, la nómina, los balances y las conciliaciones tributarias; los costos especificados son registros operativos de obra. El sistema generará reportes estructurados de horas y materiales para exportación al área financiera externa. El alcance se divide en dos fases: Fase 1 (SIG, SST, captura operativa y dossier) y Fase 2 (compras, cotizaciones y tablero ampliado).

\subsection{Requerimientos funcionales y trazabilidad}
\label{sec:rf-loard}
La Tabla~\ref{tab:pb-trazabilidad} consolida los 49 requisitos, su necesidad de origen, la historia que los cubre y la prioridad M/S/C. Cada fila distingue la obligación del sistema y una comprobación de aceptación propuesta. Estas comprobaciones concretan el comportamiento; las condiciones transversales de calidad se evalúan con los RNF. Los códigos HU de la matriz remiten a las historias desarrolladas en la sección~\ref{sec:pb-loard}. Los dominios son: M1 operación (RF-01--07), M2 inventario/equipos (08--13), M3 dossier (14--18), M4 costeo/cotización (19--25), M5 compras (26--29), M6 tablero (30--34), M7 documental (35--44) y M8 talento/SST (45--49).

{\small
\begin{longtable}{@{}L{1.6cm}L{8.05cm}L{2.0cm}L{1.5cm}L{1.2cm}@{}}
\caption{Requerimientos funcionales y trazabilidad NEC--RF--HU.}\label{tab:pb-trazabilidad}\\
\toprule\rowcolor{docHead}\textbf{ID} & \textbf{Requerimiento y aceptación} & \textbf{Origen} & \textbf{HU} & \textbf{Prior.}\\\midrule\endfirsthead
\toprule\rowcolor{docHead}\textbf{ID} & \textbf{Requerimiento y aceptación} & \textbf{Origen} & \textbf{HU} & \textbf{Prior.}\\\midrule\endhead
\midrule\multicolumn{5}{r@{}}{\footnotesize Continúa en la página siguiente}\\\endfoot
\bottomrule\endlastfoot
RF-01 & El sistema debe permitir registrar el corte operativo de avance (semanal en proyectos puntuales o mensual en contratos marco) desde computador Windows en campamento.\newline \textbf{Aceptación:} Registrar un corte de cada tipo de contrato; confirmar identificador, periodo y estado de visado pendiente sin impedir crear el siguiente corte. & NEC-01 & HU-01 & M\\
RF-02 & El sistema debe exigir que cada corte indique proyecto, frente de obra, periodo de corte, cantidades ejecutadas, horas de personal e insumos consumidos.\newline \textbf{Aceptación:} Omitir un dato obligatorio y comprobar su identificación y el rechazo del envío; completarlo y verificar la persistencia de los valores. & NEC-01 & HU-01 & M\\
RF-03 & El sistema debe permitir adjuntar fotografías y evidencias técnicas vinculadas al corte de avance.\newline \textbf{Aceptación:} Adjuntar y recuperar una evidencia desde su corte y actividad; probar un archivo no admitido y un reintento de carga fallida. & NEC-01 & HU-02 & S\\
RF-04 & El sistema debe consolidar los registros del corte en actas y reportes exportables en formatos estándar (PDF/Excel).\newline \textbf{Aceptación:} Exportar un periodo de muestra y contrastar cantidades, horas e insumos con sus cortes, diferenciando registros pendientes y visados. & NEC-01 & HU-05 & S\\
RF-05 & El sistema debe permitir registrar inspecciones y novedades SST sin conexión en dispositivos móviles Android y sincronizarlas automáticamente al recuperar señal.\newline \textbf{Aceptación:} Crear una inspección sin red, recuperar conexión y comprobar una sola copia persistida; aplicar la contingencia y reserva local de RNF-07. & NEC-20 & HU-03 & S\\
RF-06 & El sistema debe registrar el avance físico por actividad frente a la cotización base, emitiendo alerta si el \% de tiempo consumido supera el \% de avance.\newline \textbf{Aceptación:} Comparar una actividad con base conocida y verificar cantidad y porcentaje; probar un caso donde tiempo consumido supere el avance. & NEC-04 & HU-06 & M\\
RF-07 & El sistema debe monitorear las horas de mano de obra y rendimientos ejecutados frente a lo cotizado (costo crítico: nómina de $\sim 40$ trabajadores).\newline \textbf{Aceptación:} Usar horas y cantidades de una actividad de muestra y reconstruir el rendimiento y su desviación contra la base cotizada. & NEC-04 & HU-07 & M\\
RF-08 & El sistema debe registrar despachos centralizados desde Bogotá, consumos en sitio bajo custodia del residente y saldos de consumibles por frente.\newline \textbf{Aceptación:} Reconstruir el saldo de un frente mediante despachos, entradas y consumos confirmados, identificando soporte y responsable de cada movimiento. & NEC-02 & HU-11 & M\\
RF-09 & El sistema debe calcular el saldo de cada consumible a partir de movimientos confirmados y distinguir el consumo real de la merma técnica presupuestada del 5\%, sin duplicarla.\newline \textbf{Aceptación:} Conciliar saldo inicial, entradas, consumos y mermas registradas; comprobar que el 5\% presupuestado no se añade de nuevo ni se descuenta como consumo ficticio. & NEC-02 & HU-11 & M\\
RF-10 & El sistema debe alertar cuando un consumible crítico llegue a su punto de reorden considerando tiempos de transporte desde Bogotá.\newline \textbf{Aceptación:} Configurar un punto de reorden y registrar un movimiento que lo alcance; comprobar el aviso con ítem, saldo, frente y responsable. & NEC-02 & HU-13 & S\\
RF-11 & El sistema debe comparar el consumo real con el presupuesto de la cotización base, considerando una sola vez la merma técnica prevista del 5\%.\newline \textbf{Aceptación:} Comparar un ítem y unidad comunes contra su presupuesto; verificar que la merma prevista se contempla una sola vez. & NEC-02 & HU-12 & M\\
RF-12 & El sistema debe mantener la hoja de vida de cada herramienta o equipo con certificados de calibración vigentes e historial de mantenimiento.\newline \textbf{Aceptación:} Consultar un equipo y recuperar sus mantenimientos y certificados con fechas, estado y trazabilidad al soporte. & NEC-03 & HU-14 & S\\
RF-13 & El sistema debe registrar la asignación por frente y responsable, impidiendo reasignar recursos sin retorno formal.\newline \textbf{Aceptación:} Intentar asignar un equipo ya entregado a otro responsable; exigir retorno o transferencia formal antes de confirmar la nueva custodia. & NEC-03 & HU-14 & S\\
RF-14 & El sistema debe registrar ensayos no destructivos (NDT) y pruebas asociadas a actividades de soldadura e ingeniería.\newline \textbf{Aceptación:} Registrar un ensayo y recuperarlo desde el proyecto y la junta o actividad; verificar que no quede como soporte huérfano. & NEC-05 & HU-15 & S\\
RF-15 & El sistema debe permitir cargar reportes en PDF de terceros especializados e idealmente extraer metadatos clave para trazabilidad (capacidad deseable pero no excluyente).\newline \textbf{Aceptación:} Cargar un PDF y completar metadatos manualmente; la falta de extracción automática no debe impedir el registro válido. & NEC-05 & HU-15 & S\\
RF-16 & El sistema debe gestionar el libro de tubería (GT-FR64 / Weld Book) y asociar la cartilla de soldadura sobre planos con juntas, soldador, electrodo y ensayos.\newline \textbf{Aceptación:} Seleccionar una junta y reconstruir soldador, electrodo, plano y ensayos; contrastar esos vínculos con la exportación GT-FR64. & NEC-05 & HU-16 & M\\
RF-17 & El sistema debe generar el dossier de calidad API 1104 estructurado en 6 carpetas contractuales para entrega en $\le 10$ días calendario desde la aceptación del acta de facturación.\newline \textbf{Aceptación:} Generar el expediente por contrato y verificar sus seis grupos, versiones y fecha inicial del plazo; distinguir generación interna de radicación y aceptación externas. & NEC-05 & HU-17 & M\\
RF-18 & El sistema debe verificar la completitud de las 6 carpetas contractuales del dossier antes de autorizar su exportación para radicación.\newline \textbf{Aceptación:} Retirar un soporte obligatorio y comprobar el listado de faltantes y el bloqueo de exportación definitiva; completarlo y repetir la revisión. & NEC-05 & HU-18 & S\\
RF-19 & El sistema debe mantener tarifas fijas congeladas para contratos marco y un catálogo de APU para proyectos u obras adicionales.\newline \textbf{Aceptación:} Consultar una tarifa marco por contrato y vigencia y crear un APU para un ítem adicional sin alterar la tarifa pactada. & NEC-06 & HU-24 & M\\
RF-20 & El sistema debe permitir copiar un APU existente como base de uno nuevo preservando la trazabilidad de origen.\newline \textbf{Aceptación:} Copiar un APU y editar el nuevo; comprobar que conserva el identificador de origen y que el original permanece intacto. & NEC-06 & HU-24 & M\\
RF-21 & El sistema debe conservar el histórico de tarifas congeladas y cotizaciones con fecha, proyecto y cliente.\newline \textbf{Aceptación:} Recuperar dos versiones históricas y comprobar fecha, cliente y proyecto, sin sobrescritura de sus valores anteriores. & NEC-06 & HU-25 & M\\
RF-22 & El sistema debe permitir buscar tarifas, rendimientos y APU históricos por palabra clave, cliente, sede y año.\newline \textbf{Aceptación:} Buscar una referencia conocida por los filtros previstos; verificar coincidencias autorizadas y el mensaje cuando no existen resultados. & NEC-06 & HU-25 & M\\
RF-23 & El sistema debe generar cotizaciones desde una plantilla base institucional (GC-FR04) adaptada por residentes y visada por la Dirección de Proyectos.\newline \textbf{Aceptación:} Generar una oferta GC-FR04 con tarifas y APU seleccionados y registrar la revisión técnica y la versión correspondiente. & NEC-07 & HU-26 & S\\
RF-24 & El sistema debe registrar el ciclo de vida de cada cotización (borrador, enviada, adjudicada o descartada).\newline \textbf{Aceptación:} Recorrer los estados de una oferta y verificar autor, fecha e historial; una adjudicación no genera por sí sola facturación. & NEC-07 & HU-26 & S\\
RF-25 & El sistema debe parametrizar imprevistos entre 7\% y 10\%, protegiendo el margen comercial de ganancia como dato confidencial.\newline \textbf{Aceptación:} Probar imprevistos en ambos extremos del intervalo de 7--10\% y comprobar que usuarios sin permiso no consulten el margen. & NEC-07, NEC-10 & HU-26 & S\\
RF-26 & El sistema debe registrar compras locales del residente con caja menor ($\le \$250.000$ COP) y generar requisiciones centralizadas para montos superiores.\newline \textbf{Aceptación:} Probar montos de \$250.000 COP y superiores: el primero sigue caja menor y el segundo genera requisición centralizada con soporte. & NEC-08 & HU-08 & M\\
RF-27 & El sistema debe enrutar requisiciones $> \$250.000$ COP por flujo de aprobación en Bogotá (Dirección de Proyectos / Compras / Gerencia General).\newline \textbf{Aceptación:} Comprobar la revisión de Dirección de Proyectos, la orden gestionada por Olga y la autorización de desembolso de Gerencia, con sus decisiones registradas. & NEC-08 & HU-27 & M\\
RF-28 & El sistema debe asociar cuadros comparativos de cotizaciones de proveedores centralizados gestionados en Bogotá.\newline \textbf{Aceptación:} Vincular varias ofertas a una requisición y recuperar precios, plazos, proveedor, soporte y justificación de selección. & NEC-08 & HU-28 & S\\
RF-29 & El sistema debe monitorear el ciclo integral de la orden de compra desde su autorización en Bogotá hasta su recepción física en obra.\newline \textbf{Aceptación:} Seguir una compra desde la autorización hasta la recepción; probar una entrega incompleta y verificar que su pendiente siga visible. & NEC-08 & HU-29 & M\\
RF-30 & El sistema debe mostrar por proyecto el avance físico real vs base, consumo de materiales, costo crítico de mano de obra y generar reportes para contabilidad externa.\newline \textbf{Aceptación:} Contrastar el tablero y la exportación con los cortes fuente del mismo periodo; identificar fecha y estado de consolidación, sin ejecutar contabilidad. & NEC-09 & HU-30 & M\\
RF-31 & El sistema debe alertar cuando el porcentaje de tiempo transcurrido supere el porcentaje de avance físico ejecutado.\newline \textbf{Aceptación:} Utilizar porcentajes conocidos de tiempo y avance y comprobar cuándo aparece la alerta; identificar proyecto, corte y programación base. & NEC-04 & HU-31 & S\\
RF-32 & El sistema debe mostrar la concentración del valor contratado por cliente y periodo para analizar la dependencia comercial, tomando como antecedente la concentración histórica aproximada del 80\% en Chevron.\newline \textbf{Aceptación:} Calcular la participación por cliente sobre el valor contratado del periodo e identificar el origen; distinguirla de facturación y recaudo. & NEC-10 & HU-32 & S\\
RF-33 & El sistema debe emitir alertas automáticas ante desviaciones de horas y costos operativos de mano de obra respecto de la base aprobada.\newline \textbf{Aceptación:} Introducir horas y costos superiores a su base y comprobar el aviso con actividad, corte, presupuesto, observado y diferencia verificable. & NEC-09 & HU-33 & S\\
RF-34 & El sistema debe permitir registrar lecciones aprendidas de obra vinculadas a ítems de APU o tarifas marco.\newline \textbf{Aceptación:} Registrar una lección y recuperarla desde el APU o tarifa relacionada, conservando proyecto, autor, problema y recomendación. & NEC-11 & HU-10 & C\\
RF-35 & El sistema debe gestionar el repositorio documental migrado de forma curada (últimos 3 años + WPS/PQR permanentes y formatos vigentes).\newline \textbf{Aceptación:} Revisar el inventario de migración y comprobar la inclusión de documentos permanentes y de conservación obligatoria según RNF-21 y RNF-25. & NEC-12 & HU-19 & M\\
RF-36 & El sistema debe asegurar la disponibilidad exclusiva de versiones vigentes bajo ISO 9001 a usuarios operativos, archivando automáticamente las obsoletas.\newline \textbf{Aceptación:} Publicar una versión aprobada y comprobar que sustituye a la anterior en el acceso operativo, conservando la obsoleta en historial autorizado. & NEC-12 & HU-19 & M\\
RF-37 & El sistema debe implementar reserva de documentos (check-out/check-in) impidiendo la generación de copias en conflicto.\newline \textbf{Aceptación:} Abrir dos intentos de edición sobre el mismo documento: solo el titular de la reserva puede confirmar cambios; registrar liberación o cancelación. & NEC-13 & HU-20 & M\\
RF-38 & El sistema debe indexar metadatos normalizados por documento (proceso SIG, tipo, cliente, sede, año, versión y estado).\newline \textbf{Aceptación:} Guardar un documento con sus metadatos y recuperarlo por cada campo aplicable; identificar campos faltantes antes de su publicación. & NEC-14 & HU-21 & M\\
RF-39 & El sistema debe ejecutar búsquedas ágiles por metadatos estructurados y contenido de texto.\newline \textbf{Aceptación:} Combinar filtros y texto de un PDF procesable; comprobar coincidencias, permisos y el objetivo de respuesta de RNF-01. & NEC-14 & HU-21 & M\\
RF-40 & El sistema debe aplicar control de acceso por roles (RBAC) sin usuarios directos ni portales para clientes externos o interventorías.\newline \textbf{Aceptación:} Probar acciones permitidas y denegadas con cuentas de distintos roles; verificar que la plataforma no permita altas de clientes externos. & NEC-15 & HU-34 & M\\
RF-41 & El sistema debe restringir por rol los conceptos y restricciones ocupacionales y los márgenes comerciales; no debe almacenar historias clínicas ni diagnósticos.\newline \textbf{Aceptación:} Consultar con SST, obra y un rol sin permiso; verificar la separación de conceptos y márgenes y la exclusión de información clínica. & NEC-15 & HU-34 & M\\
RF-42 & El sistema debe permitir insertar la imagen de firma manuscrita digitalizada a roles autorizados (técnicos y operativos), reservando con exclusividad las firmas de representación contractual al Gerente General y al Representante Legal.\newline \textbf{Aceptación:} Autorizar una firma sobre una versión identificada; verificar identidad, consentimiento, atribución y evidencia de integridad conforme a RNF-19. & NEC-16 & HU-22 & S\\
RF-43 & El sistema debe generar una pista de auditoría con integridad verificable (audit trail) registrando usuario, fecha UTC, IP y acción.\newline \textbf{Aceptación:} Ejecutar operaciones documentales y recuperar sus eventos con usuario, fecha UTC, IP y acción; comprobar integridad y protección contra modificación. & NEC-17 & HU-23 & M\\
RF-44 & El sistema debe emitir reportes de auditoría filtrables por documento, usuario y periodo para revisiones ISO 9001 y RUC.\newline \textbf{Aceptación:} Filtrar los mismos eventos por documento, usuario y periodo y contrastarlos con el reporte exportado, respetando permisos. & NEC-17 & HU-23 & M\\
RF-45 & El sistema debe registrar certificaciones críticas del personal (alturas, espacios confinados, soldadura) con soporte y fecha de expiración.\newline \textbf{Aceptación:} Registrar una certificación con soporte y fechas; rechazar fechas incoherentes e identificar como no acreditado un requisito sin evidencia. & NEC-18 & HU-35 & M\\
RF-46 & El sistema debe generar alertas preventivas escalonadas a 30 días y a 15 días calendario previos al vencimiento.\newline \textbf{Aceptación:} Simular las fechas de 30 y 15 días previas al vencimiento; comprobar avisos a SST y residente sin duplicaciones para el mismo hito. & NEC-18 & HU-35 & M\\
RF-47 & SST debe registrar el concepto ocupacional, su vigencia y las restricciones relevantes para la tarea; el residente consulta solo la habilitación y las condiciones operativas autorizadas. Las historias clínicas quedan con el prestador de salud.\newline \textbf{Aceptación:} Registrar un concepto con vigencia y restricciones para la tarea; verificar la vista limitada de obra y la ausencia de historias clínicas y diagnósticos. & NEC-18 & HU-36 & M\\
RF-48 & El sistema debe bloquear de forma estricta e inmediata la asignación en obra de personal No habilitado para la tarea o con certificaciones vencidas (cero tolerancia, sin dispensas).\newline \textbf{Aceptación:} Intentar asignar a una persona con requisito faltante, vencido o no habilitante; comprobar bloqueo sin opción de dispensa administrativa. & NEC-18 & HU-09 & M\\
RF-49 & El sistema debe registrar inducciones a plantas petroleras (por ejemplo Chevron) y capacitaciones normativas obligatorias.\newline \textbf{Aceptación:} Relacionar persona, planta, curso y vigencia con su evidencia; consultar cumplimiento e identificar registros vencidos o incompletos. & NEC-19 & HU-37 & S\\
\end{longtable}
}

HU-04 complementa esta matriz como control transversal del módulo M1: la Dirección de Proyectos revisa el corte, registra visado u observaciones y mantiene historial, sin bloquear nuevos registros del residente. La línea base conserva los 49 RF y este control transversal.

\subsection{Requerimientos no funcionales y restricciones técnicas y legales}
\label{sec:rnf}

\subsubsection{Convenciones}

Los requerimientos no funcionales (\texttt{RNF}) establecen atributos de calidad y condiciones verificables. Se organizan por temas de calidad, tomando ISO/IEC 25010:2023 como referencia~\cite{iso25010}. Los umbrales de rendimiento, carga, disponibilidad y usabilidad son objetivos propuestos para el piloto, no resultados de pruebas ya ejecutadas. La columna \emph{Origen} establece trazabilidad; la prioridad MoSCoW indica criticidad dentro de la fase respectiva.


\paragraph{Lectura de las verificaciones.}
El \textbf{p95} es el percentil 95: al menos el 95\% de las mediciones debe quedar por debajo del umbral indicado. El \textbf{RPO} es la pérdida máxima de datos tolerada, expresada como tiempo desde el último punto recuperable; el \textbf{RTO} es el tiempo objetivo para restablecer el servicio. Cada criterio se acompaña de un método propuesto para obtener evidencia. Antes de ejecutar las pruebas se fijarán versión, entorno, datos, responsable y resultado esperado; después se conservarán mediciones, incidencias y conclusión. Los métodos no constituyen resultados obtenidos.

\subsubsection{Rendimiento}

{\small
\begin{longtable}{@{}L{1.8cm} L{4.65cm} L{4.75cm} L{2.1cm} c@{}}
\caption{Requerimientos no funcionales de rendimiento.}\\
\toprule\rowcolor{docHead}
\textbf{ID} & \textbf{Requerimiento} & \textbf{Criterio de verificación} & \textbf{Origen} & \textbf{Prior.} \\
\midrule
\endfirsthead
\toprule\rowcolor{docHead}
\textbf{ID} & \textbf{Requerimiento} & \textbf{Criterio de verificación} & \textbf{Origen} & \textbf{Prior.} \\
\midrule
\endhead
\bottomrule
\endfoot
RNF-01 & El sistema debe responder con rapidez a las búsquedas documentales sobre el fondo curado. & El 95 \% de las búsquedas por metadatos responde en menos de 3 s sobre el acervo documental migrado y curado (últimos 3 años + WPS/PQR permanentes).\newline \textbf{Método:} Registrar conjunto de búsquedas, tamaño del fondo y entorno de prueba; calcular el percentil 95 con tiempos medidos de extremo a extremo. & RF-39 & M \\
RNF-02 & El guardado de un corte desde Windows debe responder oportunamente. & p95 inferior a 2 s con conexión estable, para un corte de 50 actividades sin contar la subida de anexos. La interfaz confirma persistencia en servidor.\newline \textbf{Método:} Medir desde el envío hasta la confirmación de persistencia; identificar red, equipo y conjunto de cortes utilizados. & RF-01 & M \\
RNF-03 & El sistema debe atender la concurrencia prevista. & En una prueba de 30 sesiones activas con el mismo conjunto de datos y carga definida, el p95 de búsquedas conserva RNF-01 y el de guardado RNF-02. El dimensionamiento es una hipótesis pendiente de prueba.\newline \textbf{Método:} Ejecutar sesiones concurrentes con la misma mezcla de operaciones y registrar errores además de tiempos de respuesta. & Transversal & S \\
RNF-04 & La generación del dossier de calidad API 1104 debe completarse en tiempo óptimo. & Un dossier estructurado en 6 carpetas contractuales con hasta 300 soportes se compila en menos de 5 min.\newline \textbf{Método:} Medir con 300 soportes representativos, declarando tamaños y formatos; verificar integridad y apertura del expediente generado. & RF-17 & S \\
\end{longtable}
}

\subsubsection{Disponibilidad y confiabilidad}

{\small
\begin{longtable}{@{}L{1.8cm} L{4.65cm} L{4.75cm} L{2.1cm} c@{}}
\caption{Requerimientos no funcionales de disponibilidad y confiabilidad.}\\
\toprule\rowcolor{docHead}
\textbf{ID} & \textbf{Requerimiento} & \textbf{Criterio de verificación} & \textbf{Origen} & \textbf{Prior.} \\
\midrule
\endfirsthead
\toprule\rowcolor{docHead}
\textbf{ID} & \textbf{Requerimiento} & \textbf{Criterio de verificación} & \textbf{Origen} & \textbf{Prior.} \\
\midrule
\endhead
\bottomrule
\endfoot
RNF-05 & El sistema debe estar disponible durante la jornada extendida de operación. & Disponibilidad mensual del 99,5 \% entre las 6:00 a.\,m. y las 8:00 p.\,m., con contingencia ante fallas de red externa.\newline \textbf{Método:} Calcular tiempo disponible sobre la ventana mensual comprometida y documentar incidentes, mantenimientos y exclusiones acordadas; distinguir caída del servicio de pérdida de conectividad local. & Transversal & S \\
RNF-06 & La información debe poder recuperarse ante contingencias de infraestructura. & Copias de seguridad diarias automatizadas con retención de 30 días, RPO $\le 24$ h y RTO $< 4$ h.\newline \textbf{Método:} Restaurar una copia en un entorno de prueba y medir pérdida de datos y tiempo de recuperación. La retención de respaldos no sustituye la conservación documental. & Transversal & M \\
RNF-07 & Las inspecciones móviles deben tolerar cortes de red sin pérdida de datos. & Conservar localmente y cifrados los registros pendientes durante al menos 72 h; probar un corte de 6 h y recuperación posterior sin duplicados. Es un margen técnico sobre el escenario empresarial de 6 h.\newline \textbf{Método:} Reabrir la aplicación con registros pendientes y simular reintentos tras recuperar red; contrastar cantidades y evidencias locales con el servidor. & RF-05, NEC-20 & S \\
RNF-08 & Los movimientos de inventario y caja menor deben ser transaccionalmente consistentes. & Una falla a mitad de una transacción no deja saldos inconsistentes. La prueba de fallo arroja 0 discrepancias.\newline \textbf{Método:} Interrumpir el guardado y repetir la solicitud; conciliar saldo inicial, movimientos confirmados y saldo final para comprobar atomicidad y ausencia de duplicados. & RF-08, RF-09 & M \\
\end{longtable}
}

\subsubsection{Usabilidad}

{\small
\begin{longtable}{@{}L{1.8cm} L{4.65cm} L{4.75cm} L{2.1cm} c@{}}
\caption{Requerimientos no funcionales de usabilidad.}\\
\toprule\rowcolor{docHead}
\textbf{ID} & \textbf{Requerimiento} & \textbf{Criterio de verificación} & \textbf{Origen} & \textbf{Prior.} \\
\midrule
\endfirsthead
\toprule\rowcolor{docHead}
\textbf{ID} & \textbf{Requerimiento} & \textbf{Criterio de verificación} & \textbf{Origen} & \textbf{Prior.} \\
\midrule
\endhead
\bottomrule
\endfoot
RNF-09 & El personal técnico de obra debe operar el sistema en PC Windows con agilidad. & Tras una inducción de 1 h, un Ingeniero Residente registra y emite un corte de obra semanal o mensual en menos de 10 min.\newline \textbf{Método:} Observar una tarea con datos de muestra tras la inducción y registrar duración, errores y ayudas solicitadas; documentar participantes y condiciones. & RF-01 & M \\
RNF-10 & La interfaz debe estar en español y respetar la nomenclatura técnica del sector. & El 100 \% de etiquetas, formularios y reportes en español, conservando códigos institucionales (GC-FR04, GT-FR64) y términos normalizados de API 1104.\newline \textbf{Método:} Revisar las pantallas, validaciones y exportaciones del alcance entregado contra un glosario aprobado por el equipo. & NEC-14 & S \\
RNF-11 & La interfaz debe cumplir criterios básicos de accesibilidad. & Contraste y tamaño de texto conformes con WCAG 2.1 nivel AA.\newline \textbf{Método:} Verificar contraste y ampliación de texto en las pantallas del piloto, documentando criterios evaluados y fallos; una revisión parcial no acredita conformidad integral. & Transversal & C \\
RNF-12 & Los mensajes de error deben indicar la causa y proponer una acción correctiva. & El 100 \% de errores operacionales incluye la causa de la falla y la acción sugerida para subsanarla.\newline \textbf{Método:} Provocar errores de validación, permisos y conexión; comprobar que el mensaje identifica la causa y una acción sin exponer datos sensibles. & Transversal & S \\
\end{longtable}
}

\subsubsection{Seguridad}

{\small
\begin{longtable}{@{}L{1.8cm} L{4.65cm} L{4.75cm} L{2.1cm} c@{}}
\caption{Requerimientos no funcionales de seguridad.}\\
\toprule\rowcolor{docHead}
\textbf{ID} & \textbf{Requerimiento} & \textbf{Criterio de verificación} & \textbf{Origen} & \textbf{Prior.} \\
\midrule
\endfirsthead
\toprule\rowcolor{docHead}
\textbf{ID} & \textbf{Requerimiento} & \textbf{Criterio de verificación} & \textbf{Origen} & \textbf{Prior.} \\
\midrule
\endhead
\bottomrule
\endfoot
RNF-13 & El acceso debe protegerse con autenticación robusta y multifactor. & Política de complejidad de contraseñas y doble factor (2FA) obligatorio para Gerencia General, Representante Legal y Administrador.\newline \textbf{Método:} Probar inicio de sesión con y sin segundo factor en los perfiles obligatorios y documentar la configuración de la política de contraseñas. & RF-40 & M \\
RNF-14 & Los datos sensibles deben cifrarse en tránsito y en reposo. & TLS 1.3 en tránsito y AES-256 o mecanismo equivalente en reposo para conceptos ocupacionales y márgenes. Se excluyen historias clínicas y nómina.\newline \textbf{Método:} Inspeccionar conexiones y configuración de almacenamiento; verificar que respaldos y copias locales sensibles mantengan protección equivalente. & RF-41, RF-47 & M \\
RNF-15 & Los permisos deben seguir mínimo privilegio y separación de información ocupacional. & Pruebas por rol: SST consulta conceptos autorizados; obra solo habilitación y vigencia; ningún rol administrativo accede a historias clínicas. No hay portal externo.\newline \textbf{Método:} Usar una matriz de permisos y cuentas de prueba para intentar lectura, descarga y modificación por interfaz y solicitudes directas. & RF-40, RF-41 & M \\
RNF-16 & Los eventos de auditoría deben preservarse con integridad verificable. & La aplicación impide editar o borrar entradas; almacena eventos append-only, verifica integridad y alerta intentos de alteración. Se prueban permisos, restauración y detección de modificaciones.\newline \textbf{Método:} Intentar modificar eventos con permisos de aplicación y alterar una copia de prueba; comprobar detección y preservación al restaurar. & RF-43 & M \\
RNF-17 & Las sesiones inactivas deben cerrarse automáticamente. & Cierre automático tras 30 min de inactividad en equipos Windows y 15 min en dispositivos móviles.\newline \textbf{Método:} Simular inactividad en ambos dispositivos y comprobar que volver a una pantalla no permita operar sin renovar la autenticación. & RF-40 & S \\
RNF-18 & El tratamiento de datos debe ajustarse a la Ley 1581 y a la Resolución 1843 de 2025. & Documentar finalidad, base de tratamiento, autorizaciones cuando procedan y atención de derechos; verificar exclusión de historias clínicas y trazabilidad de acceso a conceptos ocupacionales.\newline \textbf{Método:} Revisar formatos, finalidades, permisos y evidencias de atención de solicitudes. La revisión técnica no reemplaza la validación jurídica del responsable. & RF-41, RF-47 & M \\
\end{longtable}
}

\subsubsection{Cumplimiento normativo}

{\small
\begin{longtable}{@{}L{1.8cm} L{4.65cm} L{4.75cm} L{2.1cm} c@{}}
\caption{Requerimientos no funcionales de cumplimiento normativo.}\\
\toprule\rowcolor{docHead}
\textbf{ID} & \textbf{Requerimiento} & \textbf{Criterio de verificación} & \textbf{Origen} & \textbf{Prior.} \\
\midrule
\endfirsthead
\toprule\rowcolor{docHead}
\textbf{ID} & \textbf{Requerimiento} & \textbf{Criterio de verificación} & \textbf{Origen} & \textbf{Prior.} \\
\midrule
\endhead
\bottomrule
\endfoot
RNF-19 & La formalización debe registrar identidad, consentimiento e integridad de cada aprobación. & La imagen de firma se inserta solo por autorización autenticada del titular sobre una versión concreta, con sello temporal y hash. La aceptación contractual del mecanismo debe conservarse como evidencia.\newline \textbf{Método:} Revisar el registro de aprobación y modificar una copia del documento; comprobar que cambie su hash y no se reutilice el consentimiento anterior. & RF-42 & S \\
RNF-20 & El control documental debe apoyar ISO 9001:2015, numeral 7.5. & En las pruebas se verifican identificación, revisión, aprobación, recuperación de versiones vigentes y archivo de obsoletas. La conformidad de auditoría requiere evaluación externa.\newline \textbf{Método:} Recorrer un documento desde borrador hasta aprobación, consulta y sustitución, conservando las evidencias de cada estado. & RF-35, RF-36 & M \\
RNF-21 & La conservación documental debe aplicar una tabla de retención por tipo de registro. & Conservar al menos 20 años desde el fin de la relación laboral los registros señalados en el art. 2.2.4.6.13 del Decreto 1072~\cite{decreto1072}. Los demás siguen su plazo aplicable. No eliminar por antigüedad de migración.\newline \textbf{Método:} Probar registros de distinta clase y fecha; comprobar que una regla de migración no permita eliminar un documento sujeto a conservación. & NEC-18, NEC-19 & M \\
\end{longtable}
}

\subsubsection{Compatibilidad e interoperabilidad}

{\small
\begin{longtable}{@{}L{1.8cm} L{4.65cm} L{4.75cm} L{2.1cm} c@{}}
\caption{Requerimientos no funcionales de compatibilidad e interoperabilidad.}\\
\toprule\rowcolor{docHead}
\textbf{ID} & \textbf{Requerimiento} & \textbf{Criterio de verificación} & \textbf{Origen} & \textbf{Prior.} \\
\midrule
\endfirsthead
\toprule\rowcolor{docHead}
\textbf{ID} & \textbf{Requerimiento} & \textbf{Criterio de verificación} & \textbf{Origen} & \textbf{Prior.} \\
\midrule
\endhead
\bottomrule
\endfoot
RNF-22 & El sistema debe operar en computadores portátiles Windows como entorno primario. & Optimizado para navegadores modernos (Chrome, Edge) en Windows 10/11; aplicativo Android reservado para inspecciones y SST.\newline \textbf{Método:} Completar registro, consulta, aprobación y exportación en la matriz de navegadores y equipos Windows acordada para el piloto. & Transversal & M \\
RNF-23 & La aplicación móvil debe respetar la política de áreas clasificadas. & Solo se usa en áreas no clasificadas autorizadas. El software advierte la restricción, pero no sustituye la certificación del equipo, permisos de ingreso ni controles físicos de planta.\newline \textbf{Método:} Revisar la advertencia y el procedimiento operativo de uso móvil; no atribuir al software certificación ATEX ni control físico del acceso. & RF-05 & S \\
RNF-24 & El sistema debe importar y exportar formatos técnicos de ingeniería. & Importación de reportes NDT en PDF, matrices Excel y planos; exportación de dossier API 1104 en PDF unificado y consolidados en Excel.\newline \textbf{Método:} Usar archivos representativos de cada formato y cotejar campos, vínculos y legibilidad tras importar o exportar; un plano puede ser soporte sin edición CAD. & RF-17, RF-19 & M \\
RNF-25 & La migración desde Dropbox debe ser selectiva y conservar obligaciones de retención. & Migrar últimos 3 años operativos, WPS/PQR, formatos vigentes y todos los registros de conservación obligatoria o contractual. Comparar sumas de verificación; no eliminar originales sin autorización documentada.\newline \textbf{Método:} Comparar inventario origen-destino, metadatos y sumas de verificación; registrar exclusiones y autorizaciones antes de intervenir originales. & NEC-12, NEC-14 & M \\
RNF-26 & El sistema debe proveer exportación estructurada para el área financiera externa. & Generación de reportes estructurados de horas hombre y materiales en Excel/CSV para alimentación contable externa, sin incorporar módulos de contabilidad o facturación interna.\newline \textbf{Método:} Abrir Excel/CSV en el destino acordado y conciliar horas, unidades, materiales, proyecto y periodo con los registros fuente. & NEC-09 & S \\
\end{longtable}
}

\subsubsection{Mantenibilidad y escalabilidad}

{\small
\begin{longtable}{@{}L{1.8cm} L{4.65cm} L{4.75cm} L{2.1cm} c@{}}
\caption{Requerimientos no funcionales de mantenibilidad y escalabilidad.}\\
\toprule\rowcolor{docHead}
\textbf{ID} & \textbf{Requerimiento} & \textbf{Criterio de verificación} & \textbf{Origen} & \textbf{Prior.} \\
\midrule
\endfirsthead
\toprule\rowcolor{docHead}
\textbf{ID} & \textbf{Requerimiento} & \textbf{Criterio de verificación} & \textbf{Origen} & \textbf{Prior.} \\
\midrule
\endhead
\bottomrule
\endfoot
RNF-27 & La arquitectura debe ser modular por fases y dominios. & Desacoplamiento entre Fase 1 (SIG/SST y soporte operativo) y Fase 2 (compras, cotizaciones y tablero ampliado). Cobertura de pruebas automatizadas $\ge 70\%$ en la lógica de negocio.\newline \textbf{Método:} Ejecutar la suite y documentar la herramienta, el alcance y la métrica de cobertura; verificar que cambios de Fase 2 no rompan flujos de Fase 1. & Transversal & S \\
RNF-28 & El sistema debe escalar elásticamente según la dinámica de contratos. & Soporta duplicar el volumen documental o de proyectos simultáneos sin requerir rediseño de arquitectura.\newline \textbf{Método:} Declarar volumen base y volumen duplicado, configuración y recursos; repetir operaciones representativas y contrastar los objetivos de rendimiento aplicables. & Transversal & S \\
RNF-29 & El sistema debe contar con documentación técnica y manuales por rol. & Manual técnico, manual de usuario por rol (Residente, Yonny, Olga, Dilan, Edward, Gerencia) y guía de despliegue.\newline \textbf{Método:} Comprobar que un usuario de cada rol pueda seguir su guía y que el equipo técnico reproduzca el despliegue con la versión entregada. & Transversal & S \\
RNF-30 & El sistema debe usar exclusivamente software con licencia legítima o código abierto. & Inventario formal de componentes con licencia verificada y cero software sin licenciamiento comercial o abierto legítimo.\newline \textbf{Método:} Contrastar el inventario de dependencias con sus licencias y comprobantes aplicables, incluyendo herramientas y componentes redistribuidos. & Cadena de valor & M \\
\end{longtable}
}

\subsubsection{Restricciones técnicas y legales}

Las restricciones son condiciones que el sistema debe cumplir y que el equipo no puede negociar. La Tabla~\ref{tab:restricciones} parte de las restricciones de la formulación del problema (sección~\ref{sec:problema}) y precisa qué exige cada una y con qué RNF o historia se atiende.

{\small
\begin{longtable}{@{}L{1.8cm} L{1.8cm} L{3.3cm} L{5.0cm} L{2.8cm}@{}}
\caption{Restricciones técnicas y legales del sistema.}\label{tab:restricciones}\\
\toprule\rowcolor{docHead}
\textbf{ID} & \textbf{Tipo} & \textbf{Norma o condición} & \textbf{Qué exige al sistema} & \textbf{Cómo se atiende} \\
\midrule
\endfirsthead
\toprule\rowcolor{docHead}
\textbf{ID} & \textbf{Tipo} & \textbf{Norma o condición} & \textbf{Qué exige al sistema} & \textbf{Cómo se atiende} \\
\midrule
\endhead
\midrule
\multicolumn{5}{r@{}}{\footnotesize\emph{Continúa en la página siguiente}} \\
\endfoot
\bottomrule
\endlastfoot
RES-01 & Legal & Ley 1581~\cite{ley1581} y Resolución 1843 de 2025~\cite{res1843} & Gestionar conceptos y restricciones ocupacionales bajo mínimo privilegio. Las historias clínicas permanecen con el prestador médico; el cargo de coordinación SST no habilita acceso automático a diagnósticos. & RNF-14, RNF-15, RNF-18; HU-34, HU-36. \\
\addlinespace
RES-02 & Legal & Decreto 1072 de 2015, SG-SST~\cite{decreto1072} & Conservar los registros del SG-SST que señala el art. 2.2.4.6.13 durante al menos 20 años desde que termina la relación laboral del trabajador. Preservación electrónica con integridad verificable y respaldada. & RNF-06 y RNF-21; HU-35 y HU-37. \\
\addlinespace
RES-03 & Contractual / legal & Firmas y representación autorizada & Registrar autenticación, consentimiento, versión y evidencia de integridad; la imagen de firma por sí sola no acredita estos controles. Los roles de representación se fijan según el mandato y contrato aplicable~\cite{firma-electronica}. & RNF-19; HU-22. \\
\addlinespace
RES-04 & Legal & Resoluciones 4272 de 2021 (alturas) y 0491 de 2020 (espacios confinados) & Bloqueo estricto del sistema: cero tolerancia e impedimento automático de asignación a obra si el personal tiene examen \textit{No habilitado para la tarea} o certificaciones vencidas (sin excepciones ni dispensas). Alertas preventivas a 30d y 15d. & RF-45, RF-46 y RF-48; HU-35 y HU-09. \\
\addlinespace
RES-05 & Normativa / contractual & ISO 9001:2015, RUC y requisitos técnicos de soldadura aplicables & Organizar las seis carpetas contractuales del caso y sus evidencias; radicar el dossier en un máximo de 10 días calendario desde la aceptación del acta de facturación. Aplicar API 1104 a las evidencias de soldadura que correspondan. & RNF-16, RNF-20, RNF-24; HU-15 a HU-18. \\
\addlinespace
RES-06 & Técnica & Conectividad deficiente en sedes remotas (Buenaventura) & Tolerancia a contingencias excepcionales de pérdida de red de hasta 6 horas continuas en Buenaventura. La aplicación móvil Android para inspecciones y SST debe operar en modo offline sin pérdida de información. & RNF-07; HU-03; \nec{20}. \\
\addlinespace
RES-07 & Técnica / Seguridad & Parque tecnológico y seguridad industrial ATEX & El sistema opera primariamente desde computadores portátiles Windows en campamentos y oficinas. Queda estrictamente prohibido el uso de celulares dentro de áreas industriales clasificadas con riesgo de explosión por falta de certificación ATEX. & RNF-09, RNF-22 y RNF-23; HU-01. \\
\addlinespace
RES-08 & Técnica & Repositorio histórico en Dropbox & Migración selectiva y curada de los últimos 3 años de proyectos operativos más el fondo permanente de WPS, PQR y formatos vigentes del SIG, depurando carpetas personales y versiones obsoletas. & RNF-25; HU-19. \\
\addlinespace
RES-09 & Alcance / Financiero & Delimitación financiera y despliegue por fases & El sistema no procesa contabilidad, facturación ni nómina; genera reportes estructurados de horas hombre y materiales para el área contable externa. Se desarrolla en Fase 1 (SIG/SST y soporte operativo) y Fase 2 (compras, cotizaciones y tablero ampliado). & WON'T en backlog; RNF-26 y RNF-27; HU-30. \\
\end{longtable}
}

\subsection{Product Backlog e historias de usuario}

\label{sec:pb-loard}
El backlog se construye a partir del catálogo de necesidades (\nec{01} a \nec{20}), los requerimientos funcionales de la sección~\ref{sec:rf-loard}, las condiciones de calidad de la sección~\ref{sec:rnf} y los procesos As-Is y To-Be. Las historias se organizan por rol y épica. Un rol describe una responsabilidad, así que una misma persona puede asumir varios si la empresa lo autoriza.

El backlog tiene 37 historias: 36 cubren los 49 requerimientos funcionales y una (HU-04) representa el visado de la Dirección de Proyectos que aparece en el To-Be. Sus criterios de aceptación son una propuesta que se validará con los responsables de cada proceso.

\begingroup
\setlength{\tabcolsep}{4pt}
\renewcommand{\arraystretch}{1.12}

% Modelo uniforme: ID, título, prioridad/esfuerzo, módulo, rol, proceso, resultado, cuerpo y trazabilidad.
\newcommand{\huCampo}[1]{\cellcolor{docHead}\textbf{#1}}
\newcommand{\huBanda}[1]{%
  \multicolumn{2}{|>{\columncolor{docHead}\centering\arraybackslash}p{\dimexpr\linewidth-2\tabcolsep-2\arrayrulewidth\relax}|}{\textbf{#1}}\\\hline}
\newcommand{\huFichaModelo}[9]{%
  \par\addvspace{0.5em}
  {\small
  \begin{longtable}{|L{4.2cm}|L{\dimexpr\linewidth-4.2cm-4\tabcolsep-3\arrayrulewidth\relax}|}
    \hline
    \huCampo{Título de la \mbox{#1}} & \textbf{#2}\\\hline
    \huCampo{Módulo} & #4\\\hline
    \endfirsthead
    \hline
    \multicolumn{2}{|L{\dimexpr\linewidth-2\tabcolsep-2\arrayrulewidth\relax}|}{\textbf{#1: #2} (continuación)}\\\hline
    \endhead
    \multicolumn{2}{|r|}{\footnotesize Continúa en la página siguiente}\\\hline\endfoot
    \endlastfoot
    \huBanda{Descripción de la historia de Usuario}
    \huCampo{Cómo (Rol)} & #5\\\hline
    \huCampo{Quiero (Proceso)} & #6\\\hline
    \huCampo{Para (Resultado)} & #7\\\hline
    #8
  \end{longtable}
  \textbf{Prioridad y esfuerzo:} #3.\par
  #9\par}
}

\subsubsection{Historias de usuario en fichas individuales}
\label{sec:hu-fichas}
Dos historias funcionan como controles transversales: HU-34 controla el acceso y HU-23 conserva los eventos documentales. Cada ficha indica además sus dependencias específicas. Los proyectos, frentes, actividades, personas y contratos se cargan como datos de referencia antes de registrar operaciones, acordándose la carga inicial con los responsables del piloto.

Las 37 historias utilizan la misma estructura del modelo del curso~\cite{garcia-formato-hu}: título, módulo, descripción con rol/proceso/resultado, criterios de aceptación, descripción del contexto, dado/cuando/entonces, criterios para verificación de funcionalidad y manejo de errores. Se conservan los dos escenarios de HU-01, HU-04, HU-17, HU-30 y HU-35; las demás incluyen un escenario principal y validaciones de error específicas. La prioridad, el esfuerzo, el origen y las dependencias se indican como información complementaria a cada ficha. Los escenarios y verificaciones son especificaciones propuestas, no pruebas ejecutadas.

\paragraph{Historias de usuario: Ingeniero residente / supervisor de campo.}\mbox{}\par
\textbf{Épica: Registro operativo de cortes y avance de obra.}\par

\huFichaModelo
{HU-01}
{Registrar el corte operativo de avance}
{MUST (5 SP)}
{M1. Registro operativo, avance y visado técnico}
{Ingeniero residente / supervisor de campo}
{registrar el avance por corte (semanal en proyectos puntuales o mensual en contratos marco) desde mi computador portátil Windows en campamento}
{consolidar cantidades ejecutadas, horas de mano de obra e insumos frente a la cotización base sin intermediación de operarios}
{
\huBanda{Criterios de Aceptación}
\multicolumn{2}{|L{\dimexpr\linewidth-2\tabcolsep-2\arrayrulewidth\relax}|}{\begin{enumerate}[nosep,leftmargin=*]
\item Dado que el reporte contiene proyecto, frente, tipo de contrato, fecha de corte, cantidades ejecutadas por actividad, horas de mano de obra e insumos utilizados; cuando el residente lo envía para visado técnico, entonces el sistema le asigna identificador único, lo deja en estado \emph{visado pendiente} y notifica al Director de Proyectos (Yonny), permitiendo al residente continuar con nuevos registros de obra sin bloquearse.
\item Dado que el corte no presenta cantidades ejecutadas; cuando el residente intenta registrar el corte, entonces el sistema exige consignar obligatoriamente la causa de la parada como novedad justificativa antes de procesar el envío.
\end{enumerate}}\\\hline
\huBanda{Descripción del contexto}
\multicolumn{2}{|L{\dimexpr\linewidth-2\tabcolsep-2\arrayrulewidth\relax}|}{\textbf{Escenario 1: registro de corte con avance}. Se cumple la fecha de corte (semanal o mensual) y el residente consolida los frentes asignados.}\\\hline
\huCampo{Dado} & que el reporte contiene proyecto, frente, tipo de contrato, fecha de corte, cantidades ejecutadas por actividad, horas de mano de obra e insumos utilizados,\\\hline
\huCampo{Cuando} & el residente lo envía para visado técnico,\\\hline
\huCampo{Entonces} & el sistema le asigna identificador único, lo deja en estado \emph{visado pendiente} y notifica al Director de Proyectos (Yonny), permitiendo al residente continuar con nuevos registros de obra sin bloquearse.\\\hline
\huBanda{Descripción del contexto}
\multicolumn{2}{|L{\dimexpr\linewidth-2\tabcolsep-2\arrayrulewidth\relax}|}{\textbf{Escenario 2: corte sin avance físico}. El frente estuvo detenido por lluvia, paradas de planta o suspensión del cliente.}\\\hline
\huCampo{Dado} & que el corte no presenta cantidades ejecutadas,\\\hline
\huCampo{Cuando} & el residente intenta registrar el corte,\\\hline
\huCampo{Entonces} & el sistema exige consignar obligatoriamente la causa de la parada como novedad justificativa antes de procesar el envío.\\\hline
\huCampo{Criterios para Verificación de funcionalidad} & \begin{enumerate}[nosep,leftmargin=*]
\item El formulario opera en entorno PC Windows y no requiere reportes diarios estrictos (no se usa formato GT-FR48 como reporte diario en el sistema).
\item Se permite guardar borradores de trabajo localmente antes del envío definitivo.
\item El residente visualiza el estado de los cortes emitidos y las observaciones formuladas.
\end{enumerate}\\\hline
\huCampo{Manejo de errores} & \begin{enumerate}[nosep,leftmargin=*]
\item Falta un dato obligatorio: el sistema resalta los campos pendientes e impide el despacho.
\item Interrupción en el envío: el sistema preserva la sesión y evita duplicar registros.
\end{enumerate}\\\hline
}
{\textbf{Origen:} RF-01, RF-02; NEC-01. \textbf{Tarea del To-Be:} capturar corte semanal/mensual, insumos y horas (CB-01).\par
\textbf{Dependencias:} catálogo de proyectos, frentes y contratos.}

\huFichaModelo
{HU-02}
{Adjuntar fotografías y evidencias del corte}
{SHOULD (3 SP)}
{M1. Registro operativo, avance y visado técnico}
{Ingeniero residente / supervisor de campo}
{adjuntar fotografías y respaldos técnicos al informe de corte}
{respaldar documentalmente las actividades ejecutadas eliminando la dispersión en chats informales de WhatsApp}
{
\huBanda{Criterios de Aceptación}
\multicolumn{2}{|L{\dimexpr\linewidth-2\tabcolsep-2\arrayrulewidth\relax}|}{\begin{enumerate}[nosep,leftmargin=*]
\item Cada archivo gráfico queda indexado al proyecto, corte y actividad, registrando autor, fecha y descripción.
\item El residente puede previsualizar las imágenes cargadas y añadir notas técnicas aclaratorias.
\item Si ocurre una falla en la transferencia, el sistema identifica los archivos faltantes permitiendo reintento selectivo.
\end{enumerate}}\\\hline
\huBanda{Descripción del contexto}
\multicolumn{2}{|L{\dimexpr\linewidth-2\tabcolsep-2\arrayrulewidth\relax}|}{El residente reúne evidencias para respaldar el avance de un corte.}\\\hline
\huCampo{Dado} & existe un corte del proyecto y una actividad identificada, y el residente tiene permiso para adjuntar soportes\\\hline
\huCampo{Cuando} & selecciona fotografías o archivos admitidos, añade su descripción y confirma la carga\\\hline
\huCampo{Entonces} & las evidencias quedan relacionadas con proyecto, corte y actividad, con autor y fecha, y pueden previsualizarse\\\hline
\huCampo{Criterios para Verificación de funcionalidad} & \begin{enumerate}[nosep,leftmargin=*]
\item Adjuntar dos evidencias y recuperarlas desde el corte y la actividad correspondientes.
\item Comprobar autor, fecha, descripción y previsualización en los archivos persistidos.
\end{enumerate}\\\hline
\huCampo{Manejo de errores} & \begin{enumerate}[nosep,leftmargin=*]
\item Formato no admitido: informar el archivo rechazado y permitir reemplazarlo sin perder los demás datos.
\item Transferencia interrumpida: identificar archivos pendientes y permitir reintento selectivo sin duplicar los ya guardados.
\end{enumerate}\\\hline
}
{\textbf{Origen:} RF-03; NEC-01.\par
\textbf{Dependencias:} HU-01.}

\huFichaModelo
{HU-03}
{Capturar inspecciones y SST sin conexión en móvil}
{SHOULD (8 SP)}
{M1. Registro operativo, avance y visado técnico}
{Ingeniero residente / supervisor de campo}
{capturar inspecciones o novedades SST desde el aplicativo móvil Android sin conexión y sincronizarlas al recuperar señal}
{asegurar la continuidad operativa en sedes con conectividad inestable (como contingencias excepcionales en Buenaventura de hasta 6 horas continuas), respetando la prohibición de uso de móviles en zonas clasificadas ATEX}
{
\huBanda{Criterios de Aceptación}
\multicolumn{2}{|L{\dimexpr\linewidth-2\tabcolsep-2\arrayrulewidth\relax}|}{\begin{enumerate}[nosep,leftmargin=*]
\item La aplicación móvil Android permite diligenciar inspecciones fuera de línea, marcando el registro como \emph{pendiente de sincronización}.
\item Al restablecerse la conectividad, el sistema sincroniza automáticamente sin duplicar registros ni evidencias.
\item El sistema emite advertencia de seguridad recordando la prohibición de operar dispositivos móviles en áreas de proceso con atmósfera potencialmente explosiva (falta de certificación ATEX).
\end{enumerate}}\\\hline
\huBanda{Descripción del contexto}
\multicolumn{2}{|L{\dimexpr\linewidth-2\tabcolsep-2\arrayrulewidth\relax}|}{Una inspección autorizada se realiza fuera de un área clasificada durante una contingencia de conectividad.}\\\hline
\huCampo{Dado} & el usuario autorizado dispone de Android y no hay conexión de red\\\hline
\huCampo{Cuando} & registra una inspección y posteriormente recupera conexión\\\hline
\huCampo{Entonces} & el registro se conserva cifrado como pendiente y se sincroniza sin duplicación, confirmando su persistencia en servidor\\\hline
\huCampo{Criterios para Verificación de funcionalidad} & \begin{enumerate}[nosep,leftmargin=*]
\item Simular una pérdida de red de seis horas y recuperar todos los registros y evidencias, conforme a RNF-07.
\item Repetir la sincronización y contrastar identificadores para verificar que cada registro exista una sola vez.
\end{enumerate}\\\hline
\huCampo{Manejo de errores} & \begin{enumerate}[nosep,leftmargin=*]
\item Sincronización fallida: conservar cifrados los pendientes e informar su estado para reintento.
\item Datos obligatorios incompletos: conservar el borrador e indicar qué debe completarse; mostrar la advertencia de uso móvil solo en áreas autorizadas.
\end{enumerate}\\\hline
}
{\textbf{Origen:} RF-05; NEC-20.\par
\textbf{Dependencias:} HU-01.}

\paragraph{Historias de usuario: Director de Proyectos (Yonny).}\mbox{}\par
\textbf{Épica: Visado técnico de avances y control de proyecto.}\par

\huFichaModelo
{HU-04}
{Revisar y visar los reportes de corte}
{MUST (3 SP)}
{M1. Registro operativo, avance y visado técnico}
{Director de Proyectos (Yonny)}
{revisar y visar los informes de corte registrados por los residentes}
{validar técnicamente cantidades e insumos frente a la cotización base, aplicando la regla de negocio no bloqueante para no detener la operación en campo}
{
\huBanda{Criterios de Aceptación}
\multicolumn{2}{|L{\dimexpr\linewidth-2\tabcolsep-2\arrayrulewidth\relax}|}{\begin{enumerate}[nosep,leftmargin=*]
\item Dado que las cantidades e insumos reportados coinciden con el avance real y la programación de obra; cuando el Director de Proyectos aprueba el corte, entonces el registro queda marcado como \emph{visado técnicamente}; el tablero conserva también los cortes en trámite, pero los consolidados definitivos usan versiones visadas.
\item Dado que el director formula reparos técnicos; cuando devuelve el corte registrando la observación correspondiente, entonces el residente recibe la notificación para subsanar, pero la regla de negocio no bloqueante permite que el residente continúe registrando nuevos eventos de obra sin interrumpir la operación.
\end{enumerate}}\\\hline
\huBanda{Descripción del contexto}
\multicolumn{2}{|L{\dimexpr\linewidth-2\tabcolsep-2\arrayrulewidth\relax}|}{\textbf{Escenario 1: visado conforme}. Existen cortes emitidos por los residentes en los proyectos a cargo de la dirección.}\\\hline
\huCampo{Dado} & que las cantidades e insumos reportados coinciden con el avance real y la programación de obra,\\\hline
\huCampo{Cuando} & el Director de Proyectos aprueba el corte,\\\hline
\huCampo{Entonces} & el registro queda marcado como \emph{visado técnicamente}; el tablero conserva también los cortes en trámite, pero los consolidados definitivos usan versiones visadas.\\\hline
\huBanda{Descripción del contexto}
\multicolumn{2}{|L{\dimexpr\linewidth-2\tabcolsep-2\arrayrulewidth\relax}|}{\textbf{Escenario 2: visado con observaciones}. Se detecta una discrepancia en las cantidades ejecutadas o insumos imputados.}\\\hline
\huCampo{Dado} & que el director formula reparos técnicos,\\\hline
\huCampo{Cuando} & devuelve el corte registrando la observación correspondiente,\\\hline
\huCampo{Entonces} & el residente recibe la notificación para subsanar, pero la regla de negocio no bloqueante permite que el residente continúe registrando nuevos eventos de obra sin interrumpir la operación.\\\hline
\huCampo{Criterios para Verificación de funcionalidad} & \begin{enumerate}[nosep,leftmargin=*]
\item La revisión no bloquea la creación de nuevos registros operativos por parte de los frentes de obra.
\item Se mantiene historial auditable de observaciones y ajustes efectuados sobre cada versión de corte.
\item Los tableros gerenciales distinguen claramente entre cortes con visado definitivo y cortes en trámite.
\end{enumerate}\\\hline
\huCampo{Manejo de errores} & \begin{enumerate}[nosep,leftmargin=*]
\item Intento de formular reparo sin texto de observación: el sistema exige ingresar la justificación.
\item Acceso sin credenciales de dirección: el sistema bloquea la acción de visado.
\end{enumerate}\\\hline
}
{\textbf{Origen:} NEC-01, NEC-04; regla de negocio confirmada. \textbf{Tarea del To-Be:} revisar y visar avance vs. cotización base (CB-02).\par
\textbf{Dependencias:} HU-01.}

\paragraph{Historias de usuario: Ingeniero residente / supervisor de campo.}\mbox{}\par
\textbf{Épica: Seguimiento operativo y control de insumos.}\par

\huFichaModelo
{HU-05}
{Consolidar el informe de avance por corte}
{SHOULD (5 SP)}
{M1. Registro operativo, avance y visado técnico}
{Ingeniero residente / supervisor de campo}
{generar el informe consolidado por corte (semanal o mensual) a partir de los datos registrados}
{agilizar la emisión de actas de avance físico para cobro y control de obra sin reprocesos manuales}
{
\huBanda{Criterios de Aceptación}
\multicolumn{2}{|L{\dimexpr\linewidth-2\tabcolsep-2\arrayrulewidth\relax}|}{\begin{enumerate}[nosep,leftmargin=*]
\item El residente selecciona proyecto, frente, contrato y periodo de corte.
\item El informe agrupa cantidades por actividad y unidad, sumando exclusivamente unidades métricas homogéneas.
\item Permite exportar el consolidado a PDF y Excel con los soportes de trazabilidad requeridos por el cliente.
\end{enumerate}}\\\hline
\huBanda{Descripción del contexto}
\multicolumn{2}{|L{\dimexpr\linewidth-2\tabcolsep-2\arrayrulewidth\relax}|}{El residente prepara el informe de avance de un periodo semanal o mensual.}\\\hline
\huCampo{Dado} & existen cortes relacionados con proyecto, frente, contrato y periodo, con su estado de visado\\\hline
\huCampo{Cuando} & selecciona el periodo y solicita consolidar y exportar el informe\\\hline
\huCampo{Entonces} & obtiene cantidades agrupadas por actividad y unidad homogénea, con horas e insumos y distinción entre datos pendientes y visados\\\hline
\huCampo{Criterios para Verificación de funcionalidad} & \begin{enumerate}[nosep,leftmargin=*]
\item Contrastar los totales de una muestra con los cortes de origen sin sumar unidades incompatibles.
\item Abrir las exportaciones PDF y Excel y comprobar filtros, periodo, estado y correspondencia de datos.
\end{enumerate}\\\hline
\huCampo{Manejo de errores} & \begin{enumerate}[nosep,leftmargin=*]
\item Periodo sin registros: informar que no hay información para consolidar.
\item Fallo de exportación: conservar los filtros y permitir repetir la generación sin alterar los cortes.
\end{enumerate}\\\hline
}
{\textbf{Origen:} RF-04; NEC-01.\par
\textbf{Dependencias:} HU-01 y HU-04.}

\huFichaModelo
{HU-06}
{Consultar el avance físico por actividad}
{MUST (5 SP)}
{M1. Registro operativo, avance y visado técnico}
{Ingeniero residente / supervisor de campo}
{comparar las cantidades ejecutadas con las presupuestadas en la cotización base}
{identificar diferencias de avance durante la obra y detectar desviaciones tempranas}
{
\huBanda{Criterios de Aceptación}
\multicolumn{2}{|L{\dimexpr\linewidth-2\tabcolsep-2\arrayrulewidth\relax}|}{\begin{enumerate}[nosep,leftmargin=*]
\item Se cargan las cantidades presupuestadas y unidades de cada ítem de cotización del proyecto.
\item El avance ejecutado se calcula automáticamente a partir de los cortes registrados, mostrando el porcentaje de avance físico.
\item Si el \% de tiempo transcurrido supera el \% de avance ejecutado, el sistema muestra una alerta preventiva identificada por texto e indicador visual.
\end{enumerate}}\\\hline
\huBanda{Descripción del contexto}
\multicolumn{2}{|L{\dimexpr\linewidth-2\tabcolsep-2\arrayrulewidth\relax}|}{El residente revisa la ejecución física frente a la base aprobada del proyecto.}\\\hline
\huCampo{Dado} & hay cantidades y unidades cotizadas por actividad y cortes con avance registrado\\\hline
\huCampo{Cuando} & consulta la comparación del avance y del tiempo transcurrido\\\hline
\huCampo{Entonces} & se muestran cantidades base y ejecutadas, porcentaje de avance y alerta preventiva cuando el porcentaje de tiempo supera el avance\\\hline
\huCampo{Criterios para Verificación de funcionalidad} & \begin{enumerate}[nosep,leftmargin=*]
\item Calcular manualmente una actividad y contrastar cantidad ejecutada, base y porcentaje mostrado.
\item Probar porcentajes de tiempo inferiores y superiores al avance y verificar la condición de la alerta.
\end{enumerate}\\\hline
\huCampo{Manejo de errores} & \begin{enumerate}[nosep,leftmargin=*]
\item Base faltante o cantidad base nula: señalar que el porcentaje no puede calcularse sin una referencia válida.
\item Unidades incompatibles: identificar los ítems que requieren corrección antes de compararlos.
\end{enumerate}\\\hline
}
{\textbf{Origen:} RF-06; NEC-04.\par
\textbf{Dependencias:} HU-01 y HU-04.}

\huFichaModelo
{HU-07}
{Monitorear mano de obra y rendimientos}
{MUST (5 SP)}
{M1. Registro operativo, avance y visado técnico}
{Ingeniero residente / supervisor de campo}
{comparar las horas hombre y rendimientos de cuadrilla ejecutados frente a lo presupuestado en la cotización}
{controlar de forma prioritaria el costo más crítico del proyecto (nómina de $\sim 40$ trabajadores) antes de afectar el margen}
{
\huBanda{Criterios de Aceptación}
\multicolumn{2}{|L{\dimexpr\linewidth-2\tabcolsep-2\arrayrulewidth\relax}|}{\begin{enumerate}[nosep,leftmargin=*]
\item Registra cantidades ejecutadas y horas hombre laboradas por actividad y frente.
\item El sistema calcula el rendimiento unitario y lo contrasta con la tasa de rendimiento estipulada en la cotización o APU.
\item Si el costo o consumo de horas hombre excede el presupuesto base, se genera alerta inmediata en la interfaz del residente.
\end{enumerate}}\\\hline
\huBanda{Descripción del contexto}
\multicolumn{2}{|L{\dimexpr\linewidth-2\tabcolsep-2\arrayrulewidth\relax}|}{El residente necesita interpretar el consumo de horas de una cuadrilla durante un corte.}\\\hline
\huCampo{Dado} & hay horas y cantidades ejecutadas por actividad y una referencia cotizada de rendimiento\\\hline
\huCampo{Cuando} & consulta el rendimiento y las desviaciones de mano de obra\\\hline
\huCampo{Entonces} & el sistema compara horas y rendimiento real con la base de la actividad y muestra las desviaciones de costo operativo\\\hline
\huCampo{Criterios para Verificación de funcionalidad} & \begin{enumerate}[nosep,leftmargin=*]
\item Reconstruir el rendimiento unitario con horas y cantidades de una muestra y compararlo con el resultado.
\item Verificar que la alerta identifique actividad, corte y datos de la base sin liquidar nómina.
\end{enumerate}\\\hline
\huCampo{Manejo de errores} & \begin{enumerate}[nosep,leftmargin=*]
\item Cantidad ejecutada nula: no dividir por cero; mostrar horas registradas y explicar por qué no hay rendimiento calculable.
\item Referencia cotizada o tarifa operativa faltante: señalar la limitación y evitar presentar un costo no sustentado.
\end{enumerate}\\\hline
}
{\textbf{Origen:} RF-07; NEC-04.\par
\textbf{Dependencias:} HU-01, HU-04 y referencia de HU-24/HU-25.}

\textbf{Épica: Abastecimiento y compras en obra.}\par

\huFichaModelo
{HU-08}
{Crear requisiciones y compras locales con caja menor}
{MUST (3 SP)}
{M5. Compras, caja menor y abastecimiento}
{Ingeniero residente / supervisor de campo}
{gestionar compras de insumos locales en sitio y solicitar insumos mayores a Bogotá}
{ejercer mi autonomía de compra local hasta \$250.000 COP con caja menor (\$3M-\$5M COP) o canalizar requisiciones superiores a \$250.000 COP hacia la administración central en Bogotá}
{
\huBanda{Criterios de Aceptación}
\multicolumn{2}{|L{\dimexpr\linewidth-2\tabcolsep-2\arrayrulewidth\relax}|}{\begin{enumerate}[nosep,leftmargin=*]
\item Para compras locales $\le \$250.000$ COP, el residente registra el gasto con soporte de factura/recibo con cargo a su caja menor sin requerir visto bueno previo de Bogotá.
\item Para solicitudes $> \$250.000$ COP, el sistema genera una requisición formal enrutada a Bogotá (Olga / Yonny) con descripción, justificación técnica y urgencia.
\item El sistema descuenta del saldo de caja menor las compras locales y mantiene el saldo visible al registrar o sincronizar la información.
\item El residente puede consultar el estado de tramitación de sus requisiciones enviadas a Bogotá.
\end{enumerate}}\\\hline
\huBanda{Descripción del contexto}
\multicolumn{2}{|L{\dimexpr\linewidth-2\tabcolsep-2\arrayrulewidth\relax}|}{Un frente de obra necesita comprar un insumo y determinar si utiliza caja menor o gestión central.}\\\hline
\huCampo{Dado} & el residente dispone de proyecto, frente y caja menor identificados y conoce el monto de la solicitud\\\hline
\huCampo{Cuando} & registra concepto, monto, soporte y justificación de la necesidad\\\hline
\huCampo{Entonces} & las compras de hasta \$250.000 COP se registran en caja menor y las superiores generan requisición centralizada, con saldo o estado consultable\\\hline
\huCampo{Criterios para Verificación de funcionalidad} & \begin{enumerate}[nosep,leftmargin=*]
\item Probar una compra de \$250.000 COP y otra superior y comprobar la ruta aplicada.
\item Conciliar el saldo de caja menor con los movimientos confirmados y verificar los datos de la requisición enviada.
\end{enumerate}\\\hline
\huCampo{Manejo de errores} & \begin{enumerate}[nosep,leftmargin=*]
\item Soporte o dato obligatorio faltante: identificarlo antes de confirmar el movimiento.
\item Movimiento incompatible con el saldo disponible: impedir su confirmación y solicitar revisión; un reintento no debe duplicar el gasto.
\end{enumerate}\\\hline
}
{\textbf{Origen:} RF-26; NEC-08.\par
\textbf{Dependencias:} Catálogo de proyectos, frentes y caja menor.}

\textbf{Épica: Asignación segura y cero tolerancia en SST.}\par

\huFichaModelo
{HU-09}
{Bloqueo estricto por no aptitud médica o certificaciones vencidas}
{MUST (5 SP)}
{M8. Talento humano, habilitación y SST}
{Ingeniero residente / supervisor de campo}
{recibir el bloqueo automático al intentar asignar personal no habilitado para la tarea médicamente o con certificaciones vencidas a cualquier actividad de obra}
{controlar la habilitación de tareas y reducir el riesgo de asignaciones no autorizadas por clientes petroleros, visualizando únicamente el estado habilitación por tarea sin acceder a diagnósticos clínicos}
{
\huBanda{Criterios de Aceptación}
\multicolumn{2}{|L{\dimexpr\linewidth-2\tabcolsep-2\arrayrulewidth\relax}|}{\begin{enumerate}[nosep,leftmargin=*]
\item Al intentar asignar un trabajador a un frente o tarea crítica (alturas, confinados, soldadura), el sistema verifica al registrar o sincronizar la información su concepto médico y certificaciones.
\item Si el concepto es \emph{No habilitado para la tarea} o una certificación obligatoria está vencida, el sistema bloquea terminantemente la asignación (cero tolerancia, no existen excepciones ni dispensas).
\item La interfaz muestra únicamente la causal normativa (\emph{Concepto médico no habilitado para la tarea} o \emph{Curso de alturas vencido}) y su fecha, sin revelar información sobre patologías o diagnósticos clínicos (Ley 1581 / Res. 1843 de 2025).
\end{enumerate}}\\\hline
\huBanda{Descripción del contexto}
\multicolumn{2}{|L{\dimexpr\linewidth-2\tabcolsep-2\arrayrulewidth\relax}|}{El residente intenta asignar a un trabajador a una tarea que exige requisitos específicos.}\\\hline
\huCampo{Dado} & la tarea identifica certificaciones y concepto ocupacional exigidos y existe información del trabajador\\\hline
\huCampo{Cuando} & solicita confirmar la asignación\\\hline
\huCampo{Entonces} & el sistema verifica requisitos y vigencia; solo permite la asignación habilitada y bloquea faltantes, vencimientos o conceptos no habilitantes, sin dispensas\\\hline
\huCampo{Criterios para Verificación de funcionalidad} & \begin{enumerate}[nosep,leftmargin=*]
\item Probar personas habilitadas, con certificado vencido y con concepto no habilitante, y verificar el resultado de cada asignación.
\item Comprobar que el residente vea causa operativa y vigencia, sin diagnósticos ni historia clínica.
\end{enumerate}\\\hline
\huCampo{Manejo de errores} & \begin{enumerate}[nosep,leftmargin=*]
\item Requisito faltante o pendiente: no interpretar la ausencia de información como habilitación.
\item Intento de dispensa o acceso a datos clínicos: denegar la operación y conservar la trazabilidad aplicable.
\end{enumerate}\\\hline
}
{\textbf{Origen:} RF-48; NEC-18.\par
\textbf{Dependencias:} HU-35 y HU-36.}

\textbf{Épica: Aprendizaje entre proyectos.}\par

\huFichaModelo
{HU-10}
{Registrar lecciones aprendidas}
{COULD (2 SP)}
{M6. Seguimiento gerencial y aprendizaje}
{Ingeniero residente / supervisor de campo}
{registrar fallas técnicas, desviaciones de rendimiento y oportunidades de mejora observadas en obra}
{incorporar la experiencia adquirida al preparar nuevas cotizaciones y ajustes metodológicos}
{
\huBanda{Criterios de Aceptación}
\multicolumn{2}{|L{\dimexpr\linewidth-2\tabcolsep-2\arrayrulewidth\relax}|}{\begin{enumerate}[nosep,leftmargin=*]
\item Cada registro identifica proyecto, actividad, situación ocurrida, causa identificada, acción correctiva y responsable.
\item Se permite vincular la lección aprendida con ítems específicos de APU o tarifas marco.
\item El módulo comercial puede consultar el repositorio histórico al momento de formular nuevas ofertas.
\end{enumerate}}\\\hline
\huBanda{Descripción del contexto}
\multicolumn{2}{|L{\dimexpr\linewidth-2\tabcolsep-2\arrayrulewidth\relax}|}{Tras una actividad de obra, el residente documenta un hallazgo útil para futuras ofertas.}\\\hline
\huCampo{Dado} & el proyecto y la actividad están identificados y existe un APU o tarifa que puede relacionarse\\\hline
\huCampo{Cuando} & registra situación, causa identificada, acción correctiva y recomendación\\\hline
\huCampo{Entonces} & la lección queda vinculada al proyecto y puede consultarse desde el contexto comercial sin modificar tarifas históricas\\\hline
\huCampo{Criterios para Verificación de funcionalidad} & \begin{enumerate}[nosep,leftmargin=*]
\item Guardar una lección y recuperarla desde el proyecto y el APU o tarifa relacionada.
\item Verificar que la consulta comercial conserve autor, situación y recomendación, y no sobrescriba la referencia.
\end{enumerate}\\\hline
\huCampo{Manejo de errores} & \begin{enumerate}[nosep,leftmargin=*]
\item Información insuficiente para interpretar el hallazgo: indicar los campos pendientes.
\item Referencia inexistente: solicitar una relación válida antes de confirmar el vínculo.
\end{enumerate}\\\hline
}
{\textbf{Origen:} RF-34; NEC-11.\par
\textbf{Dependencias:} HU-24.}

\paragraph{Historias de usuario: Custodia en sitio y administración de materiales.}\mbox{}\par
\textbf{Épica: Control de consumibles e inventario.}\par

\huFichaModelo
{HU-11}
{Registrar movimientos y control de consumibles}
{MUST (5 SP)}
{M2. Consumibles, herramientas y equipos}
{Ingeniero residente (en sitio) / Coordinación administrativa (Olga en Bogotá)}
{registrar despachos centralizados desde Bogotá, consumos en sitio y saldos de materiales}
{custodiar los consumibles en obra ante la ausencia de almacenista en sitio y considerar el 5\% de chatarra técnica}
{
\huBanda{Criterios de Aceptación}
\multicolumn{2}{|L{\dimexpr\linewidth-2\tabcolsep-2\arrayrulewidth\relax}|}{\begin{enumerate}[nosep,leftmargin=*]
\item Los consumibles estándar (electrodos, discos de corte, EPP) se despachan centralizadamente desde Bogotá y el residente confirma la recepción en campamento.
\item El residente custodia y descuenta los consumos por frente, distinguiendo las mermas reales registradas del 5\% ya previsto en el presupuesto, sin duplicarlo.
\item No se pueden realizar salidas por cantidades superiores al saldo disponible.
\item Cada movimiento queda asentado con integridad verificable, registrando autor, fecha y frente beneficiario.
\end{enumerate}}\\\hline
\huBanda{Descripción del contexto}
\multicolumn{2}{|L{\dimexpr\linewidth-2\tabcolsep-2\arrayrulewidth\relax}|}{Un despacho de consumibles llega desde Bogotá y posteriormente se utiliza en el frente.}\\\hline
\huCampo{Dado} & hay materiales, unidades, frente, responsable y saldo inicial identificados\\\hline
\huCampo{Cuando} & el residente confirma entradas y registra consumos o mermas reales con su soporte\\\hline
\huCampo{Entonces} & se actualiza el saldo con movimientos trazables, diferenciando consumos y mermas y sin duplicar la merma presupuestada\\\hline
\huCampo{Criterios para Verificación de funcionalidad} & \begin{enumerate}[nosep,leftmargin=*]
\item Conciliar saldo inicial, entradas, consumos y mermas con el saldo final de una muestra.
\item Intentar una salida superior al saldo y comprobar su rechazo y la conservación del saldo anterior.
\end{enumerate}\\\hline
\huCampo{Manejo de errores} & \begin{enumerate}[nosep,leftmargin=*]
\item Cantidad no válida o superior a existencias: rechazar el movimiento e indicar el motivo.
\item Interrupción o envío repetido: no dejar una operación parcial ni duplicar el movimiento confirmado.
\end{enumerate}\\\hline
}
{\textbf{Origen:} RF-08, RF-09; NEC-02.\par
\textbf{Dependencias:} Catálogo de materiales y proyectos.}

\huFichaModelo
{HU-12}
{Comparar consumo real frente a cotización base}
{MUST (3 SP)}
{M2. Consumibles, herramientas y equipos}
{Ingeniero residente (en sitio) / Coordinación administrativa (Olga en Bogotá)}
{comparar el consumo acumulado de insumos con la base cotizada incluyendo chatarra técnica}
{detectar desviaciones de materiales oportunamente antes del cierre del proyecto}
{
\huBanda{Criterios de Aceptación}
\multicolumn{2}{|L{\dimexpr\linewidth-2\tabcolsep-2\arrayrulewidth\relax}|}{\begin{enumerate}[nosep,leftmargin=*]
\item Se compara el volumen consumido real frente al presupuesto de la cotización base, contemplando una sola vez la merma técnica prevista del 5\%.
\item El sistema desagrega desviaciones por tipo de consumible (soldadura, corte, combustibles, elementos de fijación).
\item Se generan avisos automáticos cuando el ritmo de consumo excede la proyección física ejecutada.
\end{enumerate}}\\\hline
\huBanda{Descripción del contexto}
\multicolumn{2}{|L{\dimexpr\linewidth-2\tabcolsep-2\arrayrulewidth\relax}|}{El residente revisa si el consumo de un frente se aparta de la cotización base.}\\\hline
\huCampo{Dado} & existen consumos confirmados y cantidades presupuestadas en unidades comparables\\\hline
\huCampo{Cuando} & consulta el consumo acumulado por frente e ítem\\\hline
\huCampo{Entonces} & se muestran consumo, presupuesto y desviación, aplicando una sola vez la merma técnica prevista del 5\%\\\hline
\huCampo{Criterios para Verificación de funcionalidad} & \begin{enumerate}[nosep,leftmargin=*]
\item Contrastar un ítem con cálculo manual y comprobar que el 5\% no se contabilice dos veces.
\item Recuperar los movimientos de origen de la desviación por tipo de insumo y frente.
\end{enumerate}\\\hline
\huCampo{Manejo de errores} & \begin{enumerate}[nosep,leftmargin=*]
\item Presupuesto inexistente: informar que no hay base suficiente para determinar desviación.
\item Unidad o periodo inconsistente: identificar la referencia que requiere revisión antes de consolidar.
\end{enumerate}\\\hline
}
{\textbf{Origen:} RF-11; NEC-02.\par
\textbf{Dependencias:} HU-11.}

\huFichaModelo
{HU-13}
{Recibir avisos de reorden de consumibles}
{SHOULD (3 SP)}
{M2. Consumibles, herramientas y equipos}
{Ingeniero residente (en sitio) / Coordinación administrativa (Olga en Bogotá)}
{configurar umbrales mínimos de stock para consumibles críticos en obra}
{anticipar solicitudes de reabastecimiento a Bogotá antes del agotamiento físico}
{
\huBanda{Criterios de Aceptación}
\multicolumn{2}{|L{\dimexpr\linewidth-2\tabcolsep-2\arrayrulewidth\relax}|}{\begin{enumerate}[nosep,leftmargin=*]
\item Permite fijar un stock de seguridad por consumible según el tiempo de transporte desde Bogotá a la sede (hasta 48 h para frentes lejanos).
\item Al alcanzarse el umbral de reorden, el sistema notifica al residente y a compras en Bogotá (Olga).
\item El aviso se cierra una vez confirmada la recepción del nuevo despacho.
\end{enumerate}}\\\hline
\huBanda{Descripción del contexto}
\multicolumn{2}{|L{\dimexpr\linewidth-2\tabcolsep-2\arrayrulewidth\relax}|}{La obra necesita anticipar la reposición de un consumible abastecido desde Bogotá.}\\\hline
\huCampo{Dado} & el material crítico tiene un punto de reorden y un responsable definidos\\\hline
\huCampo{Cuando} & un movimiento confirmado deja el saldo en el punto de reorden o por debajo\\\hline
\huCampo{Entonces} & se genera un aviso con material, frente y saldo y se permite consultar el seguimiento de reposición\\\hline
\huCampo{Criterios para Verificación de funcionalidad} & \begin{enumerate}[nosep,leftmargin=*]
\item Registrar saldos por encima y por debajo del punto definido y verificar cuándo se produce el aviso.
\item Comprobar destinatario, ítem y frente y que el aviso no autorice automáticamente una compra.
\end{enumerate}\\\hline
\huCampo{Manejo de errores} & \begin{enumerate}[nosep,leftmargin=*]
\item Punto de reorden sin configurar: señalar que el material requiere parametrización.
\item Fallo en la notificación: conservar el aviso pendiente e identificar el destinatario para revisión.
\end{enumerate}\\\hline
}
{\textbf{Origen:} RF-10; NEC-03.\par
\textbf{Dependencias:} HU-11.}

\paragraph{Historias de usuario: Encargado de herramientas y equipos.}\mbox{}\par
\textbf{Épica: Trazabilidad de herramientas y maquinaria.}\par

\huFichaModelo
{HU-14}
{Mantener la hoja de vida y movimientos de equipos}
{SHOULD (8 SP)}
{M2. Consumibles, herramientas y equipos}
{Encargado de herramientas y equipos}
{consultar la hoja de vida, estado de calibración y asignación de cada equipo o herramienta}
{controlar la disponibilidad operativa y asegurar el cumplimiento de calibración requerido para el dossier}
{
\huBanda{Criterios de Aceptación}
\multicolumn{2}{|L{\dimexpr\linewidth-2\tabcolsep-2\arrayrulewidth\relax}|}{\begin{enumerate}[nosep,leftmargin=*]
\item Cada equipo cuenta con identificador único, hoja de vida, certificados de calibración vigentes e historial de mantenimientos.
\item Se registran asignaciones por frente y responsable, impidiendo reasignar un equipo que figure en uso sin su devolución previa.
\item Los certificados de calibración quedan disponibles para su vinculación directa a la carpeta 4 del dossier de calidad.
\end{enumerate}}\\\hline
\huBanda{Descripción del contexto}
\multicolumn{2}{|L{\dimexpr\linewidth-2\tabcolsep-2\arrayrulewidth\relax}|}{Un responsable consulta un equipo antes de asignarlo o incorporar su certificado al dossier.}\\\hline
\huCampo{Dado} & el equipo tiene identificador, historial técnico y estado de custodia\\\hline
\huCampo{Cuando} & consulta su hoja de vida o solicita una nueva asignación\\\hline
\huCampo{Entonces} & recupera mantenimiento y calibración y solo confirma otra custodia después de un retorno o transferencia formal\\\hline
\huCampo{Criterios para Verificación de funcionalidad} & \begin{enumerate}[nosep,leftmargin=*]
\item Consultar certificados e historial de un equipo y verificar su identificación y fechas.
\item Intentar reasignarlo mientras está en uso y comprobar el control de retorno o transferencia.
\end{enumerate}\\\hline
\huCampo{Manejo de errores} & \begin{enumerate}[nosep,leftmargin=*]
\item Certificado o dato técnico faltante: mostrar la ausencia sin presentar una calibración como vigente.
\item Equipo ya asignado sin cierre de custodia: rechazar la reasignación e identificar el movimiento pendiente.
\end{enumerate}\\\hline
}
{\textbf{Origen:} RF-12, RF-13; NEC-03.\par
\textbf{Dependencias:} Catálogo de herramientas y equipos.}
\paragraph{Historias de usuario: Dossier de calidad (Dilan) y Coordinación SIG (Edward y Johana).}\mbox{}\par
\textbf{Épica: Estructuración del dossier de calidad bajo API 1104.}\par

\huFichaModelo
{HU-15}
{Registrar ensayos NDT y reportes técnicos de terceros}
{SHOULD (5 SP)}
{M3. Calidad, libro de tubería y dossier}
{Ingeniero residente con experiencia en calidad (Dilan)}
{registrar ensayos no destructivos (tintas penetrantes, ultrasonido, radiografía, pruebas hidrostáticas) a partir de reportes PDF emitidos por laboratorios terceros}
{almacenar los informes PDF y opcionalmente extraer metadatos clave para indexar los soportes directamente a la actividad y junta soldada correspondiente}
{
\huBanda{Criterios de Aceptación}
\multicolumn{2}{|L{\dimexpr\linewidth-2\tabcolsep-2\arrayrulewidth\relax}|}{\begin{enumerate}[nosep,leftmargin=*]
\item Permite cargar el reporte en PDF del tercero especializado y capturar metadatos esenciales (tipo de ensayo, junta/elemento, fecha, técnico calificado, criterio de aceptación y resultado).
\item El sistema asocia el ensayo al proyecto, hito y plano respectivo, evitando transcribir manualmente tablas extensas.
\item Si el reporte carece de resultado concluyente o firma técnica, queda catalogado como \emph{ensayo observado} e impide el cierre de la junta.
\end{enumerate}}\\\hline
\huBanda{Descripción del contexto}
\multicolumn{2}{|L{\dimexpr\linewidth-2\tabcolsep-2\arrayrulewidth\relax}|}{Calidad recibe el PDF de un ensayo emitido por un laboratorio externo.}\\\hline
\huCampo{Dado} & se conocen el proyecto, la junta o elemento y los datos de identificación del reporte\\\hline
\huCampo{Cuando} & carga el PDF y registra o revisa sus metadatos técnicos\\\hline
\huCampo{Entonces} & el ensayo queda asociado a su actividad y junta; un reporte sin conclusión o firma técnica se identifica como observado para revisión\\\hline
\huCampo{Criterios para Verificación de funcionalidad} & \begin{enumerate}[nosep,leftmargin=*]
\item Recuperar el PDF desde proyecto y junta y contrastar tipo de ensayo, fecha y resultado con el reporte.
\item Comprobar que la captura manual funcione aunque la extracción automática de metadatos no esté disponible.
\end{enumerate}\\\hline
\huCampo{Manejo de errores} & \begin{enumerate}[nosep,leftmargin=*]
\item PDF no legible o carga fallida: informar el archivo afectado y permitir reemplazo o reintento.
\item Resultado no concluyente o firma técnica faltante: marcar el reporte como observado, sin asumir conformidad de la junta.
\end{enumerate}\\\hline
}
{\textbf{Origen:} RF-14, RF-15; NEC-05.\par
\textbf{Dependencias:} Ninguna.}

\huFichaModelo
{HU-16}
{Gestionar el libro de tubería y cartilla de soldadura}
{MUST (8 SP)}
{M3. Calidad, libro de tubería y dossier}
{Ingeniero residente con experiencia en calidad (Dilan)}
{estructurar el Libro de Tubería (\textit{Weld Book}, formato GT-FR64) vinculando la cartilla de soldadura elaborada sobre planos}
{garantizar la trazabilidad integral exigida por clientes petroleros entre junta, soldador calificado, electrodo y ensayos NDT}
{
\huBanda{Criterios de Aceptación}
\multicolumn{2}{|L{\dimexpr\linewidth-2\tabcolsep-2\arrayrulewidth\relax}|}{\begin{enumerate}[nosep,leftmargin=*]
\item Cada junta soldada posee identificación única referenciada a la línea de tubería y al isométrico o plano de diseño.
\item Asocia el código del soldador (con calificación WPS/PQR vigente), lote de soldadura/electrodo y reportes NDT asociados.
\item Permite subir la cartilla de soldadura dibujada sobre planos como documento gráfico de soporte para trazabilidad visual.
\item La consulta de una junta resalta de inmediato cualquier soporte pendiente o no conforme según la norma API 1104.
\end{enumerate}}\\\hline
\huBanda{Descripción del contexto}
\multicolumn{2}{|L{\dimexpr\linewidth-2\tabcolsep-2\arrayrulewidth\relax}|}{Calidad organiza la trazabilidad de soldadura en el libro de tubería.}\\\hline
\huCampo{Dado} & hay juntas identificadas, planos, soldadores y reportes técnicos que pueden relacionarse\\\hline
\huCampo{Cuando} & registra las relaciones entre junta, soldador, electrodo y ensayos y adjunta la cartilla sobre planos\\\hline
\huCampo{Entonces} & se puede reconstruir la evidencia por junta y consultar pendientes y soportes al exportar GT-FR64\\\hline
\huCampo{Criterios para Verificación de funcionalidad} & \begin{enumerate}[nosep,leftmargin=*]
\item Seleccionar una junta y recuperar plano, soldador, electrodo y ensayos asociados.
\item Contrastar una exportación GT-FR64 con sus registros y comprobar la identificación de soportes pendientes.
\end{enumerate}\\\hline
\huCampo{Manejo de errores} & \begin{enumerate}[nosep,leftmargin=*]
\item Junta duplicada o referencia inexistente: impedir la relación inconsistente e indicar el dato a corregir.
\item Ensayo pendiente o no conforme: mostrar la condición sin completar ni aprobar evidencia ficticia.
\end{enumerate}\\\hline
}
{\textbf{Origen:} RF-16; NEC-05.\par
\textbf{Dependencias:} HU-15.}

\huFichaModelo
{HU-17}
{Generar el dossier de calidad API 1104}
{MUST (13 SP)}
{M3. Calidad, libro de tubería y dossier}
{Ingeniero residente con experiencia en calidad (Dilan)}
{compilar automáticamente el paquete del dossier de calidad de un proyecto o corte de obra bajo la norma API 1104}
{entregar el expediente técnico completo para radicación en un máximo de 10 días calendario desde la aceptación del acta de facturación, apoyando el cumplimiento contractual sin garantizar la aceptación ni el pago del cliente}
{
\huBanda{Criterios de Aceptación}
\multicolumn{2}{|L{\dimexpr\linewidth-2\tabcolsep-2\arrayrulewidth\relax}|}{\begin{enumerate}[nosep,leftmargin=*]
\item Dado que las 6 categorías obligatorias de documentos están cargadas y validadas; cuando el residente de calidad ejecuta la compilación del dossier, entonces el sistema genera el expediente unificado estructurado en las 6 carpetas contractuales definidas para este caso: 1. Documentación contractual y legal (contrato, orden de servicio, pólizas), 2. Actas de corte de obra, 3. Planos de ingeniería (mecánicos, eléctricos, de diseño y planos finales as-built), 4. Certificados de calidad y conformidad de materiales y equipos (incluyendo Libro de Tubería GT-FR64, WPS/PQR y ensayos NDT), 5. Informes técnicos semanales e informe final de obra elaborados por el residente, 6. Memorias de cálculo; con tabla de contenido hipervinculada y control de versiones.
\item Dado que existen documentos obligatorios pendientes en alguna de las 6 carpetas; cuando el residente de calidad solicita la generación, entonces el sistema emite únicamente un informe de consistencia preliminar resaltando los faltantes que impiden la radicación formal ante el cliente.
\end{enumerate}}\\\hline
\huBanda{Descripción del contexto}
\multicolumn{2}{|L{\dimexpr\linewidth-2\tabcolsep-2\arrayrulewidth\relax}|}{\textbf{Escenario 1: expediente completo y listo para radicación}. El proyecto o corte cuenta con todas las juntas inspeccionadas y soportes aprobados.}\\\hline
\huCampo{Dado} & que las 6 categorías obligatorias de documentos están cargadas y validadas,\\\hline
\huCampo{Cuando} & el residente de calidad ejecuta la compilación del dossier,\\\hline
\huCampo{Entonces} & el sistema genera el expediente unificado estructurado en las 6 carpetas contractuales definidas para este caso: 1. Documentación contractual y legal (contrato, orden de servicio, pólizas), 2. Actas de corte de obra, 3. Planos de ingeniería (mecánicos, eléctricos, de diseño y planos finales as-built), 4. Certificados de calidad y conformidad de materiales y equipos (incluyendo Libro de Tubería GT-FR64, WPS/PQR y ensayos NDT), 5. Informes técnicos semanales e informe final de obra elaborados por el residente, 6. Memorias de cálculo; con tabla de contenido hipervinculada y control de versiones.\\\hline
\huBanda{Descripción del contexto}
\multicolumn{2}{|L{\dimexpr\linewidth-2\tabcolsep-2\arrayrulewidth\relax}|}{\textbf{Escenario 2: expediente con soportes faltantes}. Falta un certificado de material o el reporte final de un ensayo NDT.}\\\hline
\huCampo{Dado} & que existen documentos obligatorios pendientes en alguna de las 6 carpetas,\\\hline
\huCampo{Cuando} & el residente de calidad solicita la generación,\\\hline
\huCampo{Entonces} & el sistema emite únicamente un informe de consistencia preliminar resaltando los faltantes que impiden la radicación formal ante el cliente.\\\hline
\huCampo{Criterios para Verificación de funcionalidad} & \begin{enumerate}[nosep,leftmargin=*]
\item La estructura del dossier cumple estrictamente los lineamientos de la norma API 1104 y los 6 componentes contractuales obligatorios.
\item La compilación se completa en menos de 5 minutos para proyectos estándar con hasta 300 soportes.
\item Se genera registro de versión con integridad verificable con estampa de tiempo UTC.
\end{enumerate}\\\hline
\huCampo{Manejo de errores} & \begin{enumerate}[nosep,leftmargin=*]
\item Archivo dañado o ilegible: se marca con advertencia y se ofrece sustitución.
\item Plazo de 10 días desde la aceptación del acta de facturación: el sistema calcula y muestra los días restantes antes del límite de radicación.
\end{enumerate}\\\hline
}
{\textbf{Origen:} RF-17; NEC-05. \textbf{Tarea del To-Be:} consolidar dossier y controlar versiones (CB-05).\par
\textbf{Dependencias:} HU-04, HU-14, HU-15, HU-16, HU-19 y HU-21.}

\huFichaModelo
{HU-18}
{Comprobar las 6 categorías mandatorias del dossier}
{SHOULD (5 SP)}
{M3. Calidad, libro de tubería y dossier}
{Ingeniero residente con experiencia en calidad (Dilan)}
{verificar la completitud de las 6 carpetas contractuales antes de la exportación final}
{asegurar que no existan rechazos o glosas de interventoría que retrasen el pago a 30 días}
{
\huBanda{Criterios de Aceptación}
\multicolumn{2}{|L{\dimexpr\linewidth-2\tabcolsep-2\arrayrulewidth\relax}|}{\begin{enumerate}[nosep,leftmargin=*]
\item La matriz de chequeo valida presencia de los 6 grupos documentales contractuales: contractual/pólizas, actas de corte, planos as-built, certificados de materiales/equipos (con Weld Book GT-FR64 y NDT), informes técnicos e informe final, y memorias de cálculo.
\item El sistema emite semáforo de completitud por carpeta y porcentaje global del dossier.
\item Bloquea el marcado de \emph{Dossier Aprobado para Radicación} si alguna carpeta mandatoria tiene soportes ausentes.
\end{enumerate}}\\\hline
\huBanda{Descripción del contexto}
\multicolumn{2}{|L{\dimexpr\linewidth-2\tabcolsep-2\arrayrulewidth\relax}|}{Calidad realiza la revisión de integridad previa a la entrega definitiva del dossier.}\\\hline
\huCampo{Dado} & existe un expediente del proyecto con su lista contractual de seis grupos y documentos asociados\\\hline
\huCampo{Cuando} & ejecuta la comprobación de integridad\\\hline
\huCampo{Entonces} & obtiene faltantes y versiones por revisar y no se autoriza la exportación para radicación definitiva hasta completar los obligatorios\\\hline
\huCampo{Criterios para Verificación de funcionalidad} & \begin{enumerate}[nosep,leftmargin=*]
\item Retirar un soporte obligatorio de un expediente de muestra y verificar su detección.
\item Incorporar el soporte, repetir la revisión y comprobar el cambio de estado sin perder el historial de versiones.
\end{enumerate}\\\hline
\huCampo{Manejo de errores} & \begin{enumerate}[nosep,leftmargin=*]
\item Soporte faltante o versión no aprobada: identificar el grupo y el documento que impiden la entrega.
\item Error durante la comprobación: no presentar el expediente como completo; informar que la revisión debe repetirse.
\end{enumerate}\\\hline
}
{\textbf{Origen:} RF-18; NEC-05.\par
\textbf{Dependencias:} HU-17.}

\textbf{Épica: Control y archivo documental del SIG.}\par

\huFichaModelo
{HU-19}
{Publicar y consultar versiones vigentes del SIG}
{MUST (5 SP)}
{M7. Gestión documental, permisos y auditoría}
{Coordinadores de Gestión (Edward y Johana)}
{publicar versiones vigentes de procedimientos y formatos controlados del SIG}
{evitar el uso de formatos obsoletos en campo y respaldar las auditorías bajo ISO 9001:2015}
{
\huBanda{Criterios de Aceptación}
\multicolumn{2}{|L{\dimexpr\linewidth-2\tabcolsep-2\arrayrulewidth\relax}|}{\begin{enumerate}[nosep,leftmargin=*]
\item Cada versión incorpora código formal ISO 9001, versión numérica, fecha de aprobación, autor y archivo fuente.
\item Al liberar una versión nueva, la anterior pasa automáticamente al archivo histórico como obsoleta.
\item Los residentes y usuarios operativos solo pueden descargar la versión oficial vigente.
\end{enumerate}}\\\hline
\huBanda{Descripción del contexto}
\multicolumn{2}{|L{\dimexpr\linewidth-2\tabcolsep-2\arrayrulewidth\relax}|}{SIG necesita publicar una versión documental aprobada sin perder el historial.}\\\hline
\huCampo{Dado} & el documento tiene código, versión y revisión y el usuario está autorizado para publicarlo\\\hline
\huCampo{Cuando} & confirma la publicación de una nueva versión\\\hline
\huCampo{Entonces} & esa versión queda disponible para uso operativo y la anterior se conserva como obsoleta en el historial autorizado\\\hline
\huCampo{Criterios para Verificación de funcionalidad} & \begin{enumerate}[nosep,leftmargin=*]
\item Publicar una versión y comprobar que las consultas operativas recuperen la vigente.
\item Consultar el historial con un rol autorizado y verificar conservación y reglas de retención.
\end{enumerate}\\\hline
\huCampo{Manejo de errores} & \begin{enumerate}[nosep,leftmargin=*]
\item Documento sin aprobación o metadatos obligatorios: impedir publicación y señalar faltantes.
\item Intento de eliminar registros sujetos a conservación: denegarlo aunque no pertenezcan al tramo operativo de migración.
\end{enumerate}\\\hline
}
{\textbf{Origen:} RF-35, RF-36; NEC-12.\par
\textbf{Dependencias:} Ninguna.}

\huFichaModelo
{HU-20}
{Reservar documentos durante la edición (Check-out/Check-in)}
{MUST (3 SP)}
{M7. Gestión documental, permisos y auditoría}
{Coordinación SIG / calidad}
{bloquear temporalmente un documento del sistema durante su modificación}
{reducir la aparición de copias duplicadas en conflicto como ocurría en el repositorio histórico de Dropbox}
{
\huBanda{Criterios de Aceptación}
\multicolumn{2}{|L{\dimexpr\linewidth-2\tabcolsep-2\arrayrulewidth\relax}|}{\begin{enumerate}[nosep,leftmargin=*]
\item Al iniciar edición de un documento controlado, el sistema asigna estado de reserva (\textit{check-out}) con usuario y hora.
\item Otros usuarios pueden visualizar o consultar la versión vigente en modo solo lectura.
\item La reserva se libera al confirmar la nueva versión (\textit{check-in}) o al cancelar formalmente la sesión de edición.
\end{enumerate}}\\\hline
\huBanda{Descripción del contexto}
\multicolumn{2}{|L{\dimexpr\linewidth-2\tabcolsep-2\arrayrulewidth\relax}|}{Dos usuarios necesitan consultar o editar un documento controlado.}\\\hline
\huCampo{Dado} & existe una versión vigente y un usuario autorizado solicita su edición\\\hline
\huCampo{Cuando} & reserva el documento y posteriormente confirma o cancela los cambios\\\hline
\huCampo{Entonces} & se identifica al titular de la reserva, otros usuarios consultan en lectura y la liberación conserva versión e historial\\\hline
\huCampo{Criterios para Verificación de funcionalidad} & \begin{enumerate}[nosep,leftmargin=*]
\item Intentar dos reservas simultáneas y comprobar que solo una permita editar.
\item Confirmar y cancelar una edición de prueba y verificar liberación, usuario y versión resultante.
\end{enumerate}\\\hline
\huCampo{Manejo de errores} & \begin{enumerate}[nosep,leftmargin=*]
\item Documento reservado por otra persona: impedir sobrescritura e indicar al titular de la reserva.
\item Edición interrumpida: conservar el estado y exigir recuperación o liberación autorizada y trazable.
\end{enumerate}\\\hline
}
{\textbf{Origen:} RF-37; NEC-13.\par
\textbf{Dependencias:} HU-19.}

\huFichaModelo
{HU-21}
{Clasificar y buscar soportes mediante metadatos}
{MUST (5 SP)}
{M7. Gestión documental, permisos y auditoría}
{Coordinación SIG / calidad}
{clasificar la documentación con metadatos estructurados y ejecutar búsquedas rápidas}
{localizar evidencias técnicas y operativas sin navegar estructuras caóticas de carpetas}
{
\huBanda{Criterios de Aceptación}
\multicolumn{2}{|L{\dimexpr\linewidth-2\tabcolsep-2\arrayrulewidth\relax}|}{\begin{enumerate}[nosep,leftmargin=*]
\item Los documentos se indexan con: proceso SIG, tipo documental, cliente, sede, año, versión y estado.
\item Motor de búsqueda por combinación de metadatos y búsqueda de texto en documentos PDF procesables.
\item Respeta estrictamente los perfiles de acceso RBAC, impidiendo visualizar metadatos de documentos restringidos.
\end{enumerate}}\\\hline
\huBanda{Descripción del contexto}
\multicolumn{2}{|L{\dimexpr\linewidth-2\tabcolsep-2\arrayrulewidth\relax}|}{Un usuario necesita localizar un soporte sin recorrer manualmente carpetas.}\\\hline
\huCampo{Dado} & los documentos tienen metadatos y contenido procesable y el usuario posee permisos de consulta\\\hline
\huCampo{Cuando} & combina filtros o busca términos del contenido indexado\\\hline
\huCampo{Entonces} & recupera coincidencias autorizadas con su versión y estado, sin exponer metadatos restringidos\\\hline
\huCampo{Criterios para Verificación de funcionalidad} & \begin{enumerate}[nosep,leftmargin=*]
\item Buscar documentos conocidos combinando proceso, tipo, cliente, sede, año, versión y estado.
\item Repetir la búsqueda con roles distintos y verificar permisos y la respuesta objetivo de RNF-01.
\end{enumerate}\\\hline
\huCampo{Manejo de errores} & \begin{enumerate}[nosep,leftmargin=*]
\item Sin coincidencias: informar el resultado y permitir ajustar filtros.
\item Archivo sin texto procesable o índice pendiente: distinguir la limitación sin atribuir contenido que no se haya indexado.
\end{enumerate}\\\hline
}
{\textbf{Origen:} RF-38, RF-39; NEC-14.\par
\textbf{Dependencias:} HU-19.}

\textbf{Épica: Legalización y formalización documental.}\par

\huFichaModelo
{HU-22}
{Insertar firma manuscrita autorizada en actas y certificados}
{SHOULD (3 SP)}
{M7. Gestión documental, permisos y auditoría}
{Roles autorizados / Gerencia}
{insertar mi imagen de firma manuscrita digitalizada autorizada en actas de avance, protocolos y certificados del dossier según mi nivel de atribución}
{formalizar los entregables técnicos y contractuales con identidad, consentimiento e integridad verificables, de acuerdo con el mecanismo aceptado en cada contrato}
{
\huBanda{Criterios de Aceptación}
\multicolumn{2}{|L{\dimexpr\linewidth-2\tabcolsep-2\arrayrulewidth\relax}|}{\begin{enumerate}[nosep,leftmargin=*]
\item El sistema permite a los roles autorizados (Residentes de obra, Dilan en calidad, Edward en SST, Yonny en dirección técnica y Gerencia) cargar y estampar su imagen de firma manuscrita digitalizada en documentos de su competencia operativa o técnica, previa autenticación y consentimiento del titular sobre una versión identificada.
\item La inserción de firmas con efectos legales y contractuales vinculantes (propuestas comerciales, actas de cierre contractual y aprobación final del Dossier de Calidad ante la operadora) está estrictamente reservada para la Gerencia General y el Representante Legal con confirmación de credenciales.
\item El documento firmado adquiere condición de solo lectura con integridad verificable, registrando estampa de tiempo UTC, usuario autenticado y hash criptográfico de integridad en la pista de auditoría.
\item La aceptación contractual del mecanismo debe conservarse como evidencia; la imagen de firma por sí sola no acredita identidad, consentimiento ni validez de la aprobación. Véase RNF-19.
\end{enumerate}}\\\hline
\huBanda{Descripción del contexto}
\multicolumn{2}{|L{\dimexpr\linewidth-2\tabcolsep-2\arrayrulewidth\relax}|}{Un titular necesita formalizar un documento dentro de sus atribuciones.}\\\hline
\huCampo{Dado} & hay una versión identificada y el titular está autenticado y autorizado para el tipo de firma\\\hline
\huCampo{Cuando} & expresa su consentimiento e inserta su imagen de firma sobre esa versión\\\hline
\huCampo{Entonces} & se conserva evidencia de identidad, consentimiento, tiempo e integridad y se respeta la reserva de representación contractual\\\hline
\huCampo{Criterios para Verificación de funcionalidad} & \begin{enumerate}[nosep,leftmargin=*]
\item Comprobar autorización para un documento operativo y denegación de representación contractual a un rol sin atribución.
\item Modificar una copia y verificar que su hash no corresponda a la versión autorizada, según RNF-19.
\end{enumerate}\\\hline
\huCampo{Manejo de errores} & \begin{enumerate}[nosep,leftmargin=*]
\item Titular no autenticado o sin consentimiento: no insertar la firma.
\item Documento modificado o atribución insuficiente: detener la aprobación y solicitar revisión sobre la versión correcta.
\end{enumerate}\\\hline
}
{\textbf{Origen:} RF-42; NEC-16.\par
\textbf{Dependencias:} HU-17, HU-19 y HU-34.}

\textbf{Épica: Pista de auditoría con integridad verificable.}\par

\huFichaModelo
{HU-23}
{Consultar el registro de auditoría (Audit Trail)}
{MUST (5 SP)}
{M7. Gestión documental, permisos y auditoría}
{Coordinación SIG / calidad}
{auditar el historial cronológico, con integridad verificable, de operaciones sobre documentos y registros}
{sustentar auditorías de certificación ISO 9001, auditorías RUC y revisiones contractuales de clientes}
{
\huBanda{Criterios de Aceptación}
\multicolumn{2}{|L{\dimexpr\linewidth-2\tabcolsep-2\arrayrulewidth\relax}|}{\begin{enumerate}[nosep,leftmargin=*]
\item Registra de manera automática y con integridad verificable cada acción (creación, lectura, edición, descarga, firma y visado) con estampa UTC, usuario e IP.
\item La aplicación impide a sus usuarios, incluidos administradores, modificar o borrar entradas; la integridad verificable permite detectar alteraciones, según RNF-16.
\item Permite exportar reportes de auditoría estructurados por rango de fechas, documento o usuario.
\end{enumerate}}\\\hline
\huBanda{Descripción del contexto}
\multicolumn{2}{|L{\dimexpr\linewidth-2\tabcolsep-2\arrayrulewidth\relax}|}{SIG o calidad necesita reconstruir qué ocurrió con un documento.}\\\hline
\huCampo{Dado} & existen eventos registrados con usuario, fecha UTC, IP y acción y el consultante tiene permiso\\\hline
\huCampo{Cuando} & filtra la bitácora por documento, usuario o periodo y solicita su exportación\\\hline
\huCampo{Entonces} & recupera la secuencia de eventos autorizada con integridad verificable, sin poder editar ni borrar entradas desde la aplicación\\\hline
\huCampo{Criterios para Verificación de funcionalidad} & \begin{enumerate}[nosep,leftmargin=*]
\item Ejecutar operaciones de muestra y contrastar sus eventos y filtros con el reporte exportado.
\item Intentar modificar una entrada y alterar una copia de prueba para verificar protección y detección según RNF-16.
\end{enumerate}\\\hline
\huCampo{Manejo de errores} & \begin{enumerate}[nosep,leftmargin=*]
\item Consulta sin permiso: denegar el acceso a eventos restringidos.
\item Fallo de integridad: señalar la anomalía y preservar la evidencia para revisión, sin corregirla silenciosamente.
\end{enumerate}\\\hline
}
{\textbf{Origen:} RF-43, RF-44; NEC-17.\par
\textbf{Dependencias:} Ninguna.}

\paragraph{Historias de usuario: Área comercial y estimación de costos.}\mbox{}\par
\textbf{Épica: APU, tarifas marco y presupuestos.}\par

\huFichaModelo
{HU-24}
{Mantener tarifas fijas de contratos marco y catálogo de APU}
{MUST (5 SP)}
{M4. Tarifas, APU y cotizaciones}
{Encargado comercial}
{administrar las tarifas fijas congeladas de contratos marco y el catálogo de APU para proyectos especiales}
{cotizar con consistencia técnica y respetar los acuerdos de precios pactados con operadoras petroleras}
{
\huBanda{Criterios de Aceptación}
\multicolumn{2}{|L{\dimexpr\linewidth-2\tabcolsep-2\arrayrulewidth\relax}|}{\begin{enumerate}[nosep,leftmargin=*]
\item Mantiene el catálogo de tarifas fijas congeladas para contratos marco vigentes (por ejemplo Chevron).
\item Permite estructurar análisis de precios unitarios nuevos únicamente para proyectos u obras adicionales no contempladas en tarifas marco.
\item Desagrega cada APU en mano de obra, materiales, equipos, consumibles y rendimientos normalizados.
\end{enumerate}}\\\hline
\huBanda{Descripción del contexto}
\multicolumn{2}{|L{\dimexpr\linewidth-2\tabcolsep-2\arrayrulewidth\relax}|}{Comercial prepara una referencia de precio para una nueva intervención.}\\\hline
\huCampo{Dado} & hay tarifas marco identificadas por contrato y un catálogo de APU accesible según rol\\\hline
\huCampo{Cuando} & consulta una tarifa o crea un nuevo APU, eventualmente copiando uno existente\\\hline
\huCampo{Entonces} & conserva la tarifa pactada, desagrega el APU y registra su origen sin modificar la referencia histórica\\\hline
\huCampo{Criterios para Verificación de funcionalidad} & \begin{enumerate}[nosep,leftmargin=*]
\item Crear un APU con mano de obra, materiales, equipos y rendimientos y verificar su desglose.
\item Copiar y editar un APU y comprobar que el original y las tarifas marco permanezcan intactos.
\end{enumerate}\\\hline
\huCampo{Manejo de errores} & \begin{enumerate}[nosep,leftmargin=*]
\item Ítem sin unidad o dato necesario: indicar el campo pendiente antes de confirmar.
\item Intento de sobrescribir una tarifa congelada sin control de vigencia: impedir la operación y exigir la gestión autorizada de la referencia.
\end{enumerate}\\\hline
}
{\textbf{Origen:} RF-19, RF-20; NEC-06.\par
\textbf{Dependencias:} Ninguna.}

\huFichaModelo
{HU-25}
{Consultar históricos de tarifas congeladas y cotizaciones}
{MUST (3 SP)}
{M4. Tarifas, APU y cotizaciones}
{Encargado comercial}
{consultar el histórico de tarifas, costos de insumos y rendimientos de obras anteriores}
{sustentar nuevas propuestas comerciales sobre bases históricas sólidas y auditables}
{
\huBanda{Criterios de Aceptación}
\multicolumn{2}{|L{\dimexpr\linewidth-2\tabcolsep-2\arrayrulewidth\relax}|}{\begin{enumerate}[nosep,leftmargin=*]
\item Permite filtrar históricos por cliente, sede, tipo de intervención (mecánica, civil, piping) y vigencia.
\item Muestra la evolución de rendimientos de cuadrilla para calibrar ofertas.
\item Las consultas respetan los perfiles de confidencialidad comercial asignados.
\end{enumerate}}\\\hline
\huBanda{Descripción del contexto}
\multicolumn{2}{|L{\dimexpr\linewidth-2\tabcolsep-2\arrayrulewidth\relax}|}{Comercial revisa antecedentes para sustentar una oferta nueva.}\\\hline
\huCampo{Dado} & hay tarifas, cotizaciones y rendimientos históricos con cliente, proyecto y fecha\\\hline
\huCampo{Cuando} & consulta y filtra las referencias por los atributos disponibles\\\hline
\huCampo{Entonces} & recupera antecedentes identificados por vigencia sin sobrescribir sus valores y respetando los permisos comerciales\\\hline
\huCampo{Criterios para Verificación de funcionalidad} & \begin{enumerate}[nosep,leftmargin=*]
\item Recuperar un antecedente conocido por palabra clave, cliente, sede y año.
\item Contrastar dos versiones y comprobar que la consulta o reutilización no cambie sus condiciones históricas.
\end{enumerate}\\\hline
\huCampo{Manejo de errores} & \begin{enumerate}[nosep,leftmargin=*]
\item Sin antecedentes para los filtros: informar la ausencia y permitir otra consulta.
\item Acceso a información reservada: denegarlo sin revelar sus valores en resultados o exportaciones.
\end{enumerate}\\\hline
}
{\textbf{Origen:} RF-21, RF-22; NEC-06.\par
\textbf{Dependencias:} HU-24.}

\huFichaModelo
{HU-26}
{Generar cotizaciones con plantilla base y control de imprevistos}
{SHOULD (8 SP)}
{M4. Tarifas, APU y cotizaciones}
{Encargado comercial (adaptada por Residente y aprobada por Yonny)}
{elaborar ofertas económicas a partir de la plantilla base institucional en formato GC-FR04}
{agilizar el ciclo de cotización configurando entre 7\% y 10\% de imprevistos y protegiendo el margen comercial confidencial}
{
\huBanda{Criterios de Aceptación}
\multicolumn{2}{|L{\dimexpr\linewidth-2\tabcolsep-2\arrayrulewidth\relax}|}{\begin{enumerate}[nosep,leftmargin=*]
\item El Ingeniero Residente puede pre-diligenciar y adaptar la plantilla base con las cantidades y alcances de la obra.
\item El Director de Proyectos (Yonny) revisa y visa técnicamente la propuesta antes de su despacho comercial.
\item Permite parametrizar el porcentaje de imprevistos dentro del rango habitual del 7\% al 10\%.
\item El margen de utilidad bruta/comercial se mantiene estrictamente restringido y protegido bajo perfil RBAC para que solo la Gerencia General tenga visibilidad.
\end{enumerate}}\\\hline
\huBanda{Descripción del contexto}
\multicolumn{2}{|L{\dimexpr\linewidth-2\tabcolsep-2\arrayrulewidth\relax}|}{El área comercial necesita elaborar una oferta a partir de referencias aprobadas.}\\\hline
\huCampo{Dado} & se dispone de plantilla GC-FR04, cantidades y tarifas o APU pertinentes\\\hline
\huCampo{Cuando} & adapta la oferta, configura imprevistos y la remite al visado técnico\\\hline
\huCampo{Entonces} & se genera una versión de la cotización con su estado e historial, manteniendo el margen restringido\\\hline
\huCampo{Criterios para Verificación de funcionalidad} & \begin{enumerate}[nosep,leftmargin=*]
\item Probar imprevistos del 7\% y 10\% y contrastar su aplicación en una oferta de muestra.
\item Recorrer borrador, enviada, adjudicada o descartada y verificar revisión técnica, historial y permisos sobre margen.
\end{enumerate}\\\hline
\huCampo{Manejo de errores} & \begin{enumerate}[nosep,leftmargin=*]
\item Porcentaje fuera del intervalo o referencia faltante: identificar la condición que debe corregirse.
\item Intento de consultar margen sin permiso: denegar la visualización sin impedir las tareas autorizadas del rol.
\end{enumerate}\\\hline
}
{\textbf{Origen:} RF-23, RF-24, RF-25; NEC-07, NEC-10.\par
\textbf{Dependencias:} HU-24 y HU-25.}

\paragraph{Historias de usuario: Aprobaciones y compras centralizadas.}\mbox{}\par
\textbf{Épica: Adquisiciones y cadena de suministro.}\par

\huFichaModelo
{HU-27}
{Aprobar requisiciones centralizadas según niveles de autorización}
{MUST (3 SP)}
{M5. Compras, caja menor y abastecimiento}
{Director de Proyectos (Yonny) / Coordinación administrativa (Olga) / Gerencia General}
{aprobar requisiciones de compras que superen la autonomía local de \$250.000 COP}
{garantizar un flujo estructurado donde el Residente solicita, Bogotá cotiza y Gerencia autoriza el desembolso}
{
\huBanda{Criterios de Aceptación}
\multicolumn{2}{|L{\dimexpr\linewidth-2\tabcolsep-2\arrayrulewidth\relax}|}{\begin{enumerate}[nosep,leftmargin=*]
\item Las compras locales $\le \$250.000$ COP no requieren este flujo (se ejecutan con caja menor en sitio).
\item Requisiciones $> \$250.000$ COP son asignadas a Bogotá: el Director de Proyectos (Yonny) o la Coordinación (Olga) verifican pertinencia técnica.
\item Una vez validada, Olga genera la Orden de Compra y la Gerencia General registra la autorización de desembolso; el giro se ejecuta por canales externos, fuera de la plataforma.
\item Toda aprobación o rechazo queda registrado con autor, fecha, monto y justificación.
\end{enumerate}}\\\hline
\huBanda{Descripción del contexto}
\multicolumn{2}{|L{\dimexpr\linewidth-2\tabcolsep-2\arrayrulewidth\relax}|}{Una requisición excede la autonomía de compra local y requiere aprobación en Bogotá.}\\\hline
\huCampo{Dado} & la solicitud supera \$250.000 COP e identifica proyecto, monto, justificación y solicitante\\\hline
\huCampo{Cuando} & los responsables revisan la pertinencia técnica, gestionan la orden y autorizan el desembolso\\\hline
\huCampo{Entonces} & se conserva cada decisión con autor, fecha y justificación, mientras el pago se realiza por canales externos\\\hline
\huCampo{Criterios para Verificación de funcionalidad} & \begin{enumerate}[nosep,leftmargin=*]
\item Comparar la ruta de una compra local de \$250.000 COP con la de una requisición superior.
\item Verificar la revisión técnica, la orden de Olga y la autorización de Gerencia y recuperar sus decisiones.
\end{enumerate}\\\hline
\huCampo{Manejo de errores} & \begin{enumerate}[nosep,leftmargin=*]
\item Autorización faltante: no presentar la solicitud como aprobada ni continuar una etapa que dependa de ella.
\item Rechazo sin justificación o monto inconsistente: solicitar corrección y conservar el estado anterior.
\end{enumerate}\\\hline
}
{\textbf{Origen:} RF-27; NEC-08.\par
\textbf{Dependencias:} HU-08.}

\huFichaModelo
{HU-28}
{Comparar cotizaciones de proveedores en Bogotá}
{SHOULD (3 SP)}
{M5. Compras, caja menor y abastecimiento}
{Coordinación administrativa y compras (Olga)}
{comparar cotizaciones de distribuidores y proveedores centralizados}
{seleccionar la alternativa más favorable en precio, calidad y plazo de entrega a frentes remotos}
{
\huBanda{Criterios de Aceptación}
\multicolumn{2}{|L{\dimexpr\linewidth-2\tabcolsep-2\arrayrulewidth\relax}|}{\begin{enumerate}[nosep,leftmargin=*]
\item Permite anexar al menos tres cotizaciones comerciales de proveedores vinculadas a la requisición.
\item Cuadro comparativo automático de precios, tiempos de despacho a sede y condiciones comerciales.
\item Registro de la selección justificada para sustentar la Orden de Compra ante Gerencia.
\end{enumerate}}\\\hline
\huBanda{Descripción del contexto}
\multicolumn{2}{|L{\dimexpr\linewidth-2\tabcolsep-2\arrayrulewidth\relax}|}{Compras compara alternativas para atender una requisición centralizada.}\\\hline
\huCampo{Dado} & la requisición tiene ofertas de proveedores con precio, condiciones y plazo\\\hline
\huCampo{Cuando} & consulta el cuadro comparativo y registra la alternativa seleccionada\\\hline
\huCampo{Entonces} & la selección queda justificada y vinculada a ofertas y soportes para su revisión por Gerencia\\\hline
\huCampo{Criterios para Verificación de funcionalidad} & \begin{enumerate}[nosep,leftmargin=*]
\item Anexar las cotizaciones de muestra y contrastar precios, proveedor, plazo y condiciones con sus soportes.
\item Recuperar el cuadro comparativo y la justificación de selección desde la requisición.
\end{enumerate}\\\hline
\huCampo{Manejo de errores} & \begin{enumerate}[nosep,leftmargin=*]
\item Datos de oferta incompletos: identificar la alternativa no comparable.
\item Selección sin justificación o soporte: exigir su registro antes de cerrar la comparación.
\end{enumerate}\\\hline
}
{\textbf{Origen:} RF-28; NEC-08.\par
\textbf{Dependencias:} HU-08 y HU-27.}

\huFichaModelo
{HU-29}
{Seguimiento al ciclo de abastecimiento y despacho}
{MUST (3 SP)}
{M5. Compras, caja menor y abastecimiento}
{Coordinación administrativa y compras (Olga)}
{monitorear el estado integral de la orden desde su emisión en Bogotá hasta la recepción física en obra}
{asegurar que los insumos despachados lleguen a tiempo y se incorporen al inventario en campamento}
{
\huBanda{Criterios de Aceptación}
\multicolumn{2}{|L{\dimexpr\linewidth-2\tabcolsep-2\arrayrulewidth\relax}|}{\begin{enumerate}[nosep,leftmargin=*]
\item Estados visibles de la orden: solicitada, cotizada, autorizada, despachada desde Bogotá y recibida en obra.
\item La confirmación de entrega en frente por parte del residente actualiza el saldo de inventario (HU-11).
\item Notificaciones automáticas de retrasos en despachos intermunicipales hacia sedes como Buenaventura o frentes petroleros.
\end{enumerate}}\\\hline
\huBanda{Descripción del contexto}
\multicolumn{2}{|L{\dimexpr\linewidth-2\tabcolsep-2\arrayrulewidth\relax}|}{Administración sigue una compra desde Bogotá hasta su recepción en obra.}\\\hline
\huCampo{Dado} & existen requisición, autorizaciones y orden con un frente de destino identificado\\\hline
\huCampo{Cuando} & registra despacho, recepción e incidencias del abastecimiento\\\hline
\huCampo{Entonces} & las etapas y pendientes permanecen visibles y la recepción confirmada se relaciona con el movimiento de inventario\\\hline
\huCampo{Criterios para Verificación de funcionalidad} & \begin{enumerate}[nosep,leftmargin=*]
\item Seguir una orden de muestra desde autorización hasta recepción y comprobar responsables y soportes.
\item Registrar una recepción parcial y verificar saldo recibido y pendiente sin duplicar la entrada de inventario.
\end{enumerate}\\\hline
\huCampo{Manejo de errores} & \begin{enumerate}[nosep,leftmargin=*]
\item Recepción incompleta o discrepante: conservar el pendiente e identificar la incidencia.
\item Reintento de recepción ya confirmada: evitar duplicar la entrada o modificar el historial sin trazabilidad.
\end{enumerate}\\\hline
}
{\textbf{Origen:} RF-29; NEC-08.\par
\textbf{Dependencias:} HU-08, HU-11 y HU-27.}

\paragraph{Historias de usuario: Gerencia general y control directivo.}\mbox{}\par
\textbf{Épica: Tablero gerencial y monitoreo de costos críticos.}\par

\huFichaModelo
{HU-30}
{Consultar el tablero gerencial y reportes para finanzas}
{MUST (8 SP)}
{M6. Seguimiento gerencial y aprendizaje}
{Gerencia General}
{consultar el estado integral de los frentes de obra y generar reportes estructurados de horas y materiales para contabilidad}
{supervisar el avance físico real, la custodia de insumos y facilitar la información necesaria para el área financiera externa}
{
\huBanda{Criterios de Aceptación}
\multicolumn{2}{|L{\dimexpr\linewidth-2\tabcolsep-2\arrayrulewidth\relax}|}{\begin{enumerate}[nosep,leftmargin=*]
\item Dado que los cortes de obra han sido registrados por los residentes; cuando la gerencia visualiza el tablero de mando, entonces el sistema exhibe avance físico consolidado, consumos de insumos, horas hombre laboradas y porcentaje de tiempo consumido.
\item Dado que se selecciona un periodo de corte y proyecto; cuando la gerencia o administración ejecuta la exportación financiera, entonces el sistema emite un reporte estructurado en formato interoperable (Excel/CSV) con el desglose exacto de horas laboradas por trabajador y consumibles despachados, sin procesar facturación ni contabilidad internamente.
\end{enumerate}}\\\hline
\huBanda{Descripción del contexto}
\multicolumn{2}{|L{\dimexpr\linewidth-2\tabcolsep-2\arrayrulewidth\relax}|}{\textbf{Escenario 1: supervisión gerencial consolidada}. La Gerencia ingresa al sistema para revisar el avance de los proyectos en ejecución.}\\\hline
\huCampo{Dado} & que los cortes de obra han sido registrados por los residentes,\\\hline
\huCampo{Cuando} & la gerencia visualiza el tablero de mando,\\\hline
\huCampo{Entonces} & el sistema exhibe avance físico consolidado, consumos de insumos, horas hombre laboradas y porcentaje de tiempo consumido.\\\hline
\huBanda{Descripción del contexto}
\multicolumn{2}{|L{\dimexpr\linewidth-2\tabcolsep-2\arrayrulewidth\relax}|}{\textbf{Escenario 2: exportación de reporte estructurado para contabilidad}. El área contable externa requiere los datos de nómina operativa y consumos del corte.}\\\hline
\huCampo{Dado} & que se selecciona un periodo de corte y proyecto,\\\hline
\huCampo{Cuando} & la gerencia o administración ejecuta la exportación financiera,\\\hline
\huCampo{Entonces} & el sistema emite un reporte estructurado en formato interoperable (Excel/CSV) con el desglose exacto de horas laboradas por trabajador y consumibles despachados, sin procesar facturación ni contabilidad internamente.\\\hline
\huCampo{Criterios para Verificación de funcionalidad} & \begin{enumerate}[nosep,leftmargin=*]
\item El tablero distingue avance físico real vs. avance programado según cotización base.
\item Permite exportar reportes estructurados de insumos y mano de obra para alimentación de sistemas contables externos.
\item Los datos financieros y márgenes brutos se reservan con exclusividad a este perfil.
\end{enumerate}\\\hline
\huCampo{Manejo de errores} & \begin{enumerate}[nosep,leftmargin=*]
\item Discrepancia en registros: el tablero indica qué frentes presentan cortes pendientes de visado.
\item Falla de conexión a bases remotas: se muestran los últimos datos consolidados con fecha visible.
\end{enumerate}\\\hline
}
{\textbf{Origen:} RF-30; NEC-09. \textbf{Tarea del To-Be:} actualizar tablero gerencial (CB-03 y CB-04).\par
\textbf{Dependencias:} HU-01, HU-04, HU-06 y HU-11.}

\huFichaModelo
{HU-31}
{Alerta de avance físico vs. tiempo transcurrido}
{SHOULD (3 SP)}
{M6. Seguimiento gerencial y aprendizaje}
{Gerencia General}
{recibir una alerta urgente cuando el porcentaje de tiempo transcurrido supere el porcentaje de avance físico ejecutado}
{adoptar medidas correctivas inmediatas sobre frentes de obra atrasados y proteger la rentabilidad del proyecto}
{
\huBanda{Criterios de Aceptación}
\multicolumn{2}{|L{\dimexpr\linewidth-2\tabcolsep-2\arrayrulewidth\relax}|}{\begin{enumerate}[nosep,leftmargin=*]
\item El sistema calcula dinámicamente: $\% \text{ Tiempo Consumido} = \frac{\text{Días Transcurridos}}{\text{Plazo Contractual}} \times 100$.
\item Si $\% \text{ Tiempo Consumido} > \% \text{ Avance Físico}$, se dispara una alerta de alta prioridad resaltada en el tablero.
\item Identifica el proyecto, corte y programación de referencia; el aviso se contrasta con costos y planificación, sin atribuir automáticamente una pérdida económica.
\end{enumerate}}\\\hline
\huBanda{Descripción del contexto}
\multicolumn{2}{|L{\dimexpr\linewidth-2\tabcolsep-2\arrayrulewidth\relax}|}{Gerencia revisa si el ritmo de avance corresponde al tiempo contractual transcurrido.}\\\hline
\huCampo{Dado} & hay programación base y cortes que permiten calcular tiempo y avance del mismo proyecto\\\hline
\huCampo{Cuando} & consulta el tablero o se actualizan los datos del corte\\\hline
\huCampo{Entonces} & aparece una alerta si el porcentaje de tiempo supera al avance, con referencias que permitan revisar la situación\\\hline
\huCampo{Criterios para Verificación de funcionalidad} & \begin{enumerate}[nosep,leftmargin=*]
\item Calcular el porcentaje de días transcurridos sobre plazo contractual y contrastarlo con la pantalla.
\item Probar las condiciones de alerta e identificar proyecto, corte y fecha de la base utilizada.
\end{enumerate}\\\hline
\huCampo{Manejo de errores} & \begin{enumerate}[nosep,leftmargin=*]
\item Plazo base inválido o ausente: no calcular un porcentaje engañoso y solicitar revisión.
\item Cortes desactualizados: indicar su fecha y estado, sin presentar la alerta como prueba automática de pérdida económica.
\end{enumerate}\\\hline
}
{\textbf{Origen:} RF-31; NEC-04.\par
\textbf{Dependencias:} HU-06 y HU-30.}

\huFichaModelo
{HU-32}
{Monitorear concentración comercial de clientes}
{SHOULD (3 SP)}
{M6. Seguimiento gerencial y aprendizaje}
{Gerencia General}
{evaluar el grado de concentración del valor contratado por cliente en la cartera de la empresa}
{gestionar estratégicamente la alta dependencia comercial identificada con Chevron ($\sim 80\%$ de los ingresos)}
{
\huBanda{Criterios de Aceptación}
\multicolumn{2}{|L{\dimexpr\linewidth-2\tabcolsep-2\arrayrulewidth\relax}|}{\begin{enumerate}[nosep,leftmargin=*]
\item Gráfica de distribución del valor contratado por cliente y periodo, con trazabilidad a los contratos de origen.
\item Indicador de participación por cliente sobre el total contratado del mismo periodo, diferenciándolo de facturación y recaudo.
\item La consulta identifica el periodo, la base de cálculo y las limitaciones de los datos; no incorpora facturación tributaria al sistema.
\end{enumerate}}\\\hline
\huBanda{Descripción del contexto}
\multicolumn{2}{|L{\dimexpr\linewidth-2\tabcolsep-2\arrayrulewidth\relax}|}{Gerencia evalúa la concentración comercial por cliente.}\\\hline
\huCampo{Dado} & hay contratos y valores de un mismo periodo identificados por cliente\\\hline
\huCampo{Cuando} & consulta la distribución y participación por cliente o contrato\\\hline
\huCampo{Entonces} & obtiene porcentajes sobre el total contratado de esa base, con origen de datos y distinción frente a facturación y recaudo\\\hline
\huCampo{Criterios para Verificación de funcionalidad} & \begin{enumerate}[nosep,leftmargin=*]
\item Sumar una muestra de contratos y reconstruir las participaciones por cliente.
\item Cambiar el periodo y verificar contratos incluidos, total base y aclaración sobre el tipo de valor mostrado.
\end{enumerate}\\\hline
\huCampo{Manejo de errores} & \begin{enumerate}[nosep,leftmargin=*]
\item Total base nulo o datos insuficientes: informar que no puede calcularse concentración.
\item Contrato sin cliente o valor válido: señalar el registro a revisar sin asignar porcentajes ficticios.
\end{enumerate}\\\hline
}
{\textbf{Origen:} RF-32; NEC-10.\par
\textbf{Dependencias:} Registro de contratos y cotizaciones adjudicadas.}

\huFichaModelo
{HU-33}
{Alertas de desvío presupuestal en mano de obra}
{SHOULD (3 SP)}
{M6. Seguimiento gerencial y aprendizaje}
{Gerencia General}
{recibir notificaciones ante desviaciones de horas y costos operativos de mano de obra frente a la base aprobada}
{controlar rigurosamente el costo más crítico de la empresa (planilla fija de $\sim 40$ trabajadores)}
{
\huBanda{Criterios de Aceptación}
\multicolumn{2}{|L{\dimexpr\linewidth-2\tabcolsep-2\arrayrulewidth\relax}|}{\begin{enumerate}[nosep,leftmargin=*]
\item Compara horas y costos registrados con la base aprobada para la misma actividad y corte.
\item El aviso identifica proyecto, actividad, presupuesto, valor observado y diferencia, sin liquidar nómina.
\item Permite rastrear los registros que sustentan el aviso y señala información incompleta antes de interpretarlo.
\end{enumerate}}\\\hline
\huBanda{Descripción del contexto}
\multicolumn{2}{|L{\dimexpr\linewidth-2\tabcolsep-2\arrayrulewidth\relax}|}{Gerencia necesita identificar actividades con desviaciones de mano de obra.}\\\hline
\huCampo{Dado} & hay horas y costos operativos registrados y una base aprobada para el mismo corte y actividad\\\hline
\huCampo{Cuando} & se comparan los valores ejecutados con su presupuesto\\\hline
\huCampo{Entonces} & se muestra una alerta con proyecto, actividad, presupuesto, observado y diferencia, trazable a los registros fuente\\\hline
\huCampo{Criterios para Verificación de funcionalidad} & \begin{enumerate}[nosep,leftmargin=*]
\item Construir un caso con exceso de horas y otro de costo y verificar los datos del aviso.
\item Abrir los registros de origen y reconstruir la diferencia sin generar liquidación de nómina.
\end{enumerate}\\\hline
\huCampo{Manejo de errores} & \begin{enumerate}[nosep,leftmargin=*]
\item Tarifa o base faltante: distinguir qué comparación no puede sustentarse.
\item Registros incompletos o inconsistentes: señalarlos antes de interpretar la alerta como una desviación confirmada.
\end{enumerate}\\\hline
}
{\textbf{Origen:} RF-33; NEC-09.\par
\textbf{Dependencias:} HU-07 y HU-30.}

\paragraph{Historias de usuario: Administración de seguridad y perfiles.}\mbox{}\par
\textbf{Épica: Seguridad RBAC y confidencialidad.}\par

\huFichaModelo
{HU-34}
{Gestionar perfiles de acceso (RBAC) y reserva de datos}
{MUST (5 SP)}
{M7. Gestión documental, permisos y auditoría}
{Administrador del sistema}
{administrar usuarios, roles operativos y políticas estrictas de control de acceso}
{proteger conceptos ocupacionales y márgenes comerciales según responsabilidad, excluyendo historias clínicas, diagnósticos y cuentas de clientes externos}
{
\huBanda{Criterios de Aceptación}
\multicolumn{2}{|L{\dimexpr\linewidth-2\tabcolsep-2\arrayrulewidth\relax}|}{\begin{enumerate}[nosep,leftmargin=*]
\item No se crean cuentas de usuario ni portales web para clientes externos (Chevron) ni interventorías; la entrega se realiza por canales formales externos.
\item La información clínica y diagnósticos médicos son estrictamente invisibles para personal operativo, residentes y supervisores (excluida del sistema administrativo, incluso para administradores y coordinación SST).
\item Los márgenes comerciales quedan restringidos a los perfiles autorizados; Comercial y Dirección de Proyectos acceden solo a tarifas y APU necesarios para sus tareas, según la matriz de permisos.
\item Bloqueo automático de cuentas inactivas y registro de sesiones.
\end{enumerate}}\\\hline
\huBanda{Descripción del contexto}
\multicolumn{2}{|L{\dimexpr\linewidth-2\tabcolsep-2\arrayrulewidth\relax}|}{El administrador configura el acceso de personas con responsabilidades distintas.}\\\hline
\huCampo{Dado} & existe una matriz de roles y permisos y el administrador tiene atribución para gestionar cuentas\\\hline
\huCampo{Cuando} & crea o modifica una cuenta interna y asigna sus permisos\\\hline
\huCampo{Entonces} & el usuario solo puede ejecutar operaciones autorizadas, con reserva de conceptos y márgenes y sin acceso a historias clínicas ni cuentas externas\\\hline
\huCampo{Criterios para Verificación de funcionalidad} & \begin{enumerate}[nosep,leftmargin=*]
\item Usar cuentas de SST, obra, Comercial y Gerencia para comprobar lecturas, cambios y descargas permitidos y denegados.
\item Verificar que cambios de permisos queden registrados y que no se habiliten cuentas de clientes externos.
\end{enumerate}\\\hline
\huCampo{Manejo de errores} & \begin{enumerate}[nosep,leftmargin=*]
\item Solicitud de permiso no autorizada: impedir el cambio y conservar la configuración anterior.
\item Intento de acceder a información restringida o crear un portal externo: denegar la operación y registrar el evento aplicable.
\end{enumerate}\\\hline
}
{\textbf{Origen:} RF-40, RF-41; NEC-15.\par
\textbf{Dependencias:} Ninguna.}

\paragraph{Historias de usuario: Coordinación SG-SST y Medicina Laboral.}\mbox{}\par
\textbf{Épica: Control estricto de aptitud laboral y SG-SST (Fase 1).}\par

\huFichaModelo
{HU-35}
{Control de certificaciones críticas y alertas escalonadas (30d y 15d)}
{MUST (5 SP)}
{M8. Talento humano, habilitación y SST}
{Responsable de SST}
{controlar las certificaciones obligatorias del personal operativo (alturas, espacios confinados, soldadura) con alertas preventivas}
{prevenir la presencia de personal no habilitado en frentes de obra mediante avisos automáticos a 30 días y 15 días del vencimiento}
{
\huBanda{Criterios de Aceptación}
\multicolumn{2}{|L{\dimexpr\linewidth-2\tabcolsep-2\arrayrulewidth\relax}|}{\begin{enumerate}[nosep,leftmargin=*]
\item Dado que faltan exactamente 30 días o 15 días calendario para el vencimiento; cuando el sistema ejecuta su rutina de verificación diaria, entonces se genera un aviso prioritario notificando a la Coordinación SST y al residente del frente con la fecha límite de renovación.
\item Dado que la certificación ha expirado; cuando el sistema actualiza el estado del colaborador a \emph{Vencido / No Autorizado}, entonces bloquea inmediatamente su asignación en frentes de obra (HU-09) e impide la generación de permisos de trabajo.
\end{enumerate}}\\\hline
\huBanda{Descripción del contexto}
\multicolumn{2}{|L{\dimexpr\linewidth-2\tabcolsep-2\arrayrulewidth\relax}|}{\textbf{Escenario 1: alerta preventiva escalonada}. Un colaborador cuenta con certificado de trabajo en alturas próximo a caducar.}\\\hline
\huCampo{Dado} & que faltan exactamente 30 días o 15 días calendario para el vencimiento,\\\hline
\huCampo{Cuando} & el sistema ejecuta su rutina de verificación diaria,\\\hline
\huCampo{Entonces} & se genera un aviso prioritario notificando a la Coordinación SST y al residente del frente con la fecha límite de renovación.\\\hline
\huBanda{Descripción del contexto}
\multicolumn{2}{|L{\dimexpr\linewidth-2\tabcolsep-2\arrayrulewidth\relax}|}{\textbf{Escenario 2: certificación vencida}. Se cumple la fecha límite sin cargar el certificado renovado.}\\\hline
\huCampo{Dado} & que la certificación ha expirado,\\\hline
\huCampo{Cuando} & el sistema actualiza el estado del colaborador a \emph{Vencido / No Autorizado},\\\hline
\huCampo{Entonces} & bloquea inmediatamente su asignación en frentes de obra (HU-09) e impide la generación de permisos de trabajo.\\\hline
\huCampo{Criterios para Verificación de funcionalidad} & \begin{enumerate}[nosep,leftmargin=*]
\item Registra tipo de curso, entidad acreditada, soporte PDF y fecha de expiración.
\item Genera alertas preventivas a 30 días y a 15 días sin duplicaciones.
\item Regla de cero tolerancia: ante vencimiento, bloqueo absoluto sin excepciones.
\end{enumerate}\\\hline
\huCampo{Manejo de errores} & \begin{enumerate}[nosep,leftmargin=*]
\item Soporte faltante: el registro se marca en estado \emph{Sin Evidencia} y se inhabilita para obra.
\item Inconsistencia en fechas: validación que fecha de vencimiento sea posterior a la de expedición.
\end{enumerate}\\\hline
}
{\textbf{Origen:} RF-45, RF-46; NEC-18.\par
\textbf{Dependencias:} registro institucional de personal. Fase 1 confirmada.}

\huFichaModelo
{HU-36}
{Gestión de aptitud médica ocupacional y confidencialidad clínica}
{MUST (5 SP)}
{M8. Talento humano, habilitación y SST}
{Coordinador de Gestión / SST (Edward)}
{registrar conceptos y restricciones ocupacionales sin incorporar historias clínicas ni diagnósticos}
{cumplir rigurosamente la Ley 1581 de 2012 y la Resolución 1843 de 2025, asegurando que en obra solo se visualice la habilitación para la tarea y su vigencia}
{
\huBanda{Criterios de Aceptación}
\multicolumn{2}{|L{\dimexpr\linewidth-2\tabcolsep-2\arrayrulewidth\relax}|}{\begin{enumerate}[nosep,leftmargin=*]
\item SST registra el concepto ocupacional, sus restricciones y vigencia emitidos por el prestador. No se cargan historias clínicas ni diagnósticos a la plataforma administrativa.
\item Para los roles operativos (residentes, supervisores, gerencia), el sistema exhibe exclusivamente el estado de habilitación para la tarea (\emph{Habilitado}, \emph{No habilitado} o \emph{Pendiente de validación}) y su fecha de vigencia.
\item Olga custodia únicamente soportes administrativos autorizados; la historia clínica permanece bajo custodia del prestador médico.
\item Ante un concepto de \emph{No habilitado para la tarea}, el sistema activa bloqueo estricto impidiendo la entrada a obra sin posibilidad de dispensas.
\end{enumerate}}\\\hline
\huBanda{Descripción del contexto}
\multicolumn{2}{|L{\dimexpr\linewidth-2\tabcolsep-2\arrayrulewidth\relax}|}{SST registra el concepto emitido por el prestador para determinar habilitación por tarea.}\\\hline
\huCampo{Dado} & la persona está identificada y existe un concepto con fecha, vigencia y restricciones ocupacionales autorizadas\\\hline
\huCampo{Cuando} & SST registra el concepto y obra consulta la habilitación pertinente\\\hline
\huCampo{Entonces} & se conserva el concepto bajo permisos y obra recibe únicamente habilitación, condiciones operativas permitidas y vigencia, sin información clínica\\\hline
\huCampo{Criterios para Verificación de funcionalidad} & \begin{enumerate}[nosep,leftmargin=*]
\item Comparar las vistas de SST y residente y comprobar la exclusión de historias clínicas y diagnósticos.
\item Probar concepto vigente, vencido, faltante y no habilitante y contrastar su efecto en la asignación de HU-09.
\end{enumerate}\\\hline
\huCampo{Manejo de errores} & \begin{enumerate}[nosep,leftmargin=*]
\item Concepto sin soporte o fechas incoherentes: mantenerlo pendiente y solicitar validación.
\item Intento de cargar historia clínica o diagnóstico: no incorporarlo al repositorio administrativo; un estado pendiente no habilita tareas.
\end{enumerate}\\\hline
}
{\textbf{Origen:} RF-47; NEC-18.\par
\textbf{Dependencias:} HU-34; Fase 1 confirmada.}

\huFichaModelo
{HU-37}
{Registrar inducciones y habilitaciones para plantas petroleras}
{SHOULD (3 SP)}
{M8. Talento humano, habilitación y SST}
{Responsable de SST}
{registrar cursos de inducción del cliente (por ejemplo inducción Chevron) y capacitaciones normativas}
{garantizar que todo el personal en frente cumpla los requisitos de acceso exigidos por las operadoras}
{
\huBanda{Criterios de Aceptación}
\multicolumn{2}{|L{\dimexpr\linewidth-2\tabcolsep-2\arrayrulewidth\relax}|}{\begin{enumerate}[nosep,leftmargin=*]
\item Registra fecha, planta/sede, duración, empresa que dictó la inducción y soporte digital.
\item Controla la vigencia anual de la inducción Chevron y notifica su renovación previa.
\item Asocia las inducciones aprobadas al perfil del trabajador para validación en control de acceso.
\end{enumerate}}\\\hline
\huBanda{Descripción del contexto}
\multicolumn{2}{|L{\dimexpr\linewidth-2\tabcolsep-2\arrayrulewidth\relax}|}{SST necesita verificar la formación exigida por una planta o cliente.}\\\hline
\huCampo{Dado} & la persona, la planta y el requisito de inducción o capacitación están identificados\\\hline
\huCampo{Cuando} & registra curso, fecha, entidad, vigencia y soporte\\\hline
\huCampo{Entonces} & el requisito queda disponible para la verificación de habilitación y se identifica su renovación según las reglas del cliente\\\hline
\huCampo{Criterios para Verificación de funcionalidad} & \begin{enumerate}[nosep,leftmargin=*]
\item Recuperar el curso desde la persona y la planta y contrastar fecha, vigencia y soporte.
\item Probar un registro vigente y otro vencido y verificar su estado y aviso de renovación.
\end{enumerate}\\\hline
\huCampo{Manejo de errores} & \begin{enumerate}[nosep,leftmargin=*]
\item Soporte faltante o fecha inválida: identificar el pendiente sin declarar cumplido el requisito.
\item Regla de vigencia del cliente no definida: solicitar su confirmación antes de asignar una habilitación automática.
\end{enumerate}\\\hline
}
{\textbf{Origen:} RF-49; NEC-19.\par
\textbf{Dependencias:} HU-34; Fase 1 confirmada.}

\subsubsection{Revisión con INVEST y validación}
Las historias separan resultados y declaran dependencias (independientes), admiten refinamiento sin cambiar el beneficio (negociables), se vinculan con necesidades (valiosas) y tienen criterios y puntos (estimables). Las de 8 o 13 SP deben dividirse si exceden un sprint (pequeñas). Los criterios de aceptación de las fichas individuales, los escenarios de las 37 historias y los métodos de verificación de los RNF permiten comprobarlas (comprobables).

Las reglas empresariales se retoman del diagnóstico y validaciones declaradas por el equipo. El detalle de pruebas, los umbrales técnicos y el compromiso de entrega se revisarán con los responsables durante el piloto.

\subsubsection{Priorización MoSCoW y planificación por fases}
La prioridad expresa criticidad dentro de cada fase: MUST es indispensable para su dominio; SHOULD admite una alternativa temporal; COULD se aborda si existe capacidad. Quedan fuera la integración contable directa, los pagos, la facturación tributaria, la liquidación de nómina y la regularización de licencias ajenas al sistema.

{\small
\begin{longtable}{@{}L{3.5cm}L{2.0cm}L{2.0cm}L{6.8cm}@{}}
\caption{Distribución del esfuerzo del backlog.}\label{tab:pb-moscow}\\
\toprule\rowcolor{docHead}\textbf{Prioridad} & \textbf{Historias} & \textbf{SP} & \textbf{Interpretación}\\\midrule\endfirsthead
\toprule\rowcolor{docHead}\textbf{Prioridad} & \textbf{Historias} & \textbf{SP} & \textbf{Interpretación}\\\midrule\endhead
\midrule\multicolumn{4}{r@{}}{\footnotesize Continúa en la página siguiente}\\\endfoot
\bottomrule\endlastfoot
MUST & 22 & 110 & 59,5\% de historias y 62,9\% del esfuerzo.\\
SHOULD & 14 & 63 & 36,0\% del esfuerzo; alternativas temporales posibles.\\
COULD & 1 & 2 & 1,1\% del esfuerzo; lecciones aprendidas.\\
Total & 37 & 175 & Estimación relativa; no equivale a horas ni velocidad observada.\\
\end{longtable}
}

La orientación del 60\% de MUST de DSDM se refiere al esfuerzo, no al número de historias~\cite{dsdm-moscow}. La propuesta la supera ligeramente; por ello, antes de comprometer una entrega se debe revisar alcance y capacidad. Los puntos se conservaron según las estimaciones de las fichas.

{\small
\begin{longtable}{@{}L{2.0cm}L{8.4cm}L{1.9cm}L{2.0cm}@{}}
\caption{Alcance y escenario temporal de las fases.}\label{tab:plan-fases}\\
\toprule\rowcolor{docHead}\textbf{Fase} & \textbf{Alcance / historias} & \textbf{Esfuerzo} & \textbf{Ventana}\\\midrule\endfirsthead
\toprule\rowcolor{docHead}\textbf{Fase} & \textbf{Alcance / historias} & \textbf{Esfuerzo} & \textbf{Ventana}\\\midrule\endhead
\midrule\multicolumn{4}{r@{}}{\footnotesize Continúa en la página siguiente}\\\endfoot
\bottomrule\endlastfoot
1. Núcleo & SIG, SST, captura operativa, consumos, equipos y dossier. HU-01--07, 09, 11, 12, 14--23 y 34--37. & 125 SP / 24 HU & Sprints 1--6\\
2. Ampliación & Compras, cotizaciones, lecciones y tablero ampliado. HU-08, 10, 13 y 24--33. & 50 SP / 13 HU & Sprints 7--8\\
\end{longtable}
}

Se plantea un semestre de 16 semanas con ocho sprints de dos semanas. A una referencia de 23 SP/sprint existirían 184 SP nominales y solo 9 SP de holgura frente a 175 SP de alcance. La Fase 1 exige aproximadamente 20,8 SP/sprint y la Fase 2, 25 SP/sprint. Son hipótesis de planificación del equipo de seis integrantes, no cuotas individuales ni capacidad garantizada. Se revisarán con los primeros incrementos y se renegociarán SHOULD/COULD o calendario si resulta necesario.

\endgroup

\paragraph{Estado de validación.}
Las condiciones empresariales se consolidan en la Tabla~\ref{tab:reglas-comunes}; los valores técnicos de prueba son metas del piloto. La custodia clínica queda fuera de la plataforma, la aceptación del mecanismo de firma debe documentarse por contrato y el calendario se revisará con la velocidad observada.

\paragraph{Resultado esperado.}
Se conserva la trazabilidad entre 20 necesidades, 49 RF, 37 HU y 30 RNF. La siguiente entrega debe comprobar criterios y medir resultados con datos comparables; este documento no acredita implementación ni beneficios alcanzados.


\FloatBarrier
\section{Consideraciones éticas y uso de IA}
\begin{itemize}
    \item \textbf{Asistencia de IA:} El equipo declara el uso de Gemini 3.8 Flash como apoyo en SAM, gobernanza y Balanced Scorecard. También se utilizó ChatGPT/Codex para apoyar Porter, cadena de valor, OST, diagramación, revisión de coherencia y edición de LaTeX. Estas herramientas asistieron la elaboración; no entrevistaron a la empresa ni comprobaron por sí mismas las reglas del negocio.
    \item \textbf{Responsabilidad académica:} Los integrantes deben revisar las fuentes, conservar evidencia de las validaciones, comprender las decisiones de diseño y verificar los resultados. Las pruebas, entrevistas y mejoras no ejecutadas no se presentan como hechos cumplidos.
    \item \textbf{Protección de datos:} Los prototipos y pruebas utilizarán datos ficticios o información autorizada y anonimizada. Se excluyen las historias clínicas del repositorio administrativo y no se incorporan imágenes de firma ni información sensible a servicios de IA sin autorización y condiciones de tratamiento adecuadas.
\end{itemize}

\FloatBarrier
\addcontentsline{toc}{section}{Referencias}
\renewcommand{\refname}{Referencias}
\begin{thebibliography}{99}
\bibitem{hbs-cinco-fuerzas}
Institute for Strategy and Competitiveness, Harvard Business School.
\emph{The Five Forces}. Sin fecha.
\url{https://www.isc.hbs.edu/strategy/business-strategy/Pages/the-five-forces.aspx}.
Consultado el 10 de octubre de 2026.

\bibitem{hbs-cadena-valor}
Institute for Strategy and Competitiveness, Harvard Business School.
\emph{The Value Chain}. Sin fecha.
\url{https://www.isc.hbs.edu/strategy/business-strategy/Pages/the-value-chain.aspx}.
Consultado el 10 de octubre de 2026.

\bibitem{chan-reich-2007}
Chan, Y. E., \& Reich, B. H. (2007).
\emph{IT Alignment: What Have We Learned?}
Journal of Information Technology, 22(4), 297--315.
\url{https://doi.org/10.1057/palgrave.jit.2000109}.

\bibitem{torres-ost}
Torres, T. (2021).
\emph{Continuous Discovery Habits: Discover Products that Create Customer Value and Business Value}.
Product Talk LLC.

\bibitem{omg-bpmn}
Object Management Group. (2011).
\emph{Business Process Model and Notation (BPMN), Version 2.0}.
\url{https://www.omg.org/spec/BPMN/2.0/About-BPMN}.
Consultado el 10 de octubre de 2026.

\bibitem{garcia-formato-hu}
Garcia Camargo, J. A. (2026).
\emph{Formato modelo de historias de usuario}.
Material del curso Sistemas de Información, sesión 5. Universidad Nacional de Colombia.

\bibitem{ley1581}
Congreso de Colombia. (2012).
\emph{Ley Estatutaria 1581 de 2012, por la cual se dictan disposiciones generales para la protección de datos personales}.
\url{https://www.saludcapital.gov.co/Normo/gsp/ley_1581_de_2012.pdf}.
Consultado el 11 de octubre de 2026.

\bibitem{decreto1072}
Ministerio del Trabajo de Colombia. (2015).
\emph{Decreto 1072 de 2015, Decreto Único Reglamentario del Sector Trabajo}.
\url{https://www.funcionpublica.gov.co/eva/gestornormativo/norma.php?i=72173}.

\bibitem{res1843}
Ministerio del Trabajo. (2025). \emph{Resolución 1843 de 2025: evaluaciones médicas ocupacionales e historias clínicas ocupacionales}, arts. 30, 31 y 36.
\url{https://normograma.supersalud.gov.co/compilacion/docs/resolucion_mtra_1843_2025.htm}.

\bibitem{dsdm-moscow}
Agile Business Consortium. \emph{What is MoSCoW Prioritization?}
\url{https://www.agilebusiness.org/resource/what-is-moscow-prioritization/}.

\bibitem{api1104}
American Petroleum Institute. \emph{API Standard 1104: Welding of Pipelines and Related Facilities}.
\url{https://www.api.org/products-and-services/standards/important-standards-announcements/1104}.
Se aplicará la edición exigida por cada contrato.

\bibitem{cobit}
ISACA. \emph{COBIT: governance and management of enterprise information and technology}.
\url{https://www.isaca.org/resources/cobit}.

\bibitem{togaf}
The Open Group. \emph{The TOGAF Standard}.
\url{https://www.opengroup.org/togaf}.

\bibitem{swebok}
IEEE Computer Society. (2014). \emph{Guide to the Software Engineering Body of Knowledge (SWEBOK), Version 3.0}.
\url{https://ieeecs-media.computer.org/media/education/swebok/swebok-v3.pdf}.

\bibitem{abet}
ABET. \emph{Criteria for Accrediting Engineering Programs, 2025--2026}, Student Outcome 2.
\url{https://www.abet.org/accreditation/accreditation-criteria/criteria-for-accrediting-engineering-programs-2025-2026/}.

\bibitem{iso25010}
ISO/IEC. (2023). \emph{ISO/IEC 25010:2023: Product quality model}.
\url{https://www.iso.org/standard/78176.html}.


\bibitem{firma-electronica}
Presidencia de la República. (2021). \emph{Decreto 1789 de 2021}, uso de firma electrónica y digital; remisión al Decreto 1074 de 2015.
\url{https://www.funcionpublica.gov.co/eva/gestornormativo/norma.php?i=175127}.

\end{thebibliography}

\end{document}

